% before: asm-tied
Being tied at a scale is the conclusion that all the links aim at. It says
that the invariants of the members at $x$ form numerical data that satisfy
every hypothesis of the comparison inequality of
Chapter~\ref{chapter-inert-prime-balance}. A family tied at all its large
scales therefore satisfies the window bound there.

% before: asm-box-row-tame
The designed families are cut out by conditions on the auxiliary primes $q,d$
and on the level primes $g_0,g_1$ of their members. The box condition asks
for zero-free boxes for the quadratic characters attached to $q$ and $d$, and
the row condition places $q$ and $d$ at the size given by a scale $Y$. The
tame values at the level primes must be distinct, also after removing their
$\ell$-parts for the primes $\ell$ outside $\mathcal B_r(g_k)$.

% before: asm-twist-condition
At a twisted place the local error is controlled by twisting with a ray class
character of $\ell$-power order. The twist condition asks that such a
character be chosen at every twisted place, compatibly with complex
conjugation, so that the twisted theta lattices, Bloch-Kato lengths and theta
values change by at most $\ell^{-2}$. It enters through the Tate-Shafarevich
link of Theorem~\ref{asm-definite-side-exists}.

% before: asm-designed
A family can satisfy all the links only if some of its constants were chosen
with care: the thin-scale bound for host characters, the proviso on the
ranks, and the sizes of the auxiliary primes $T_u$ and of the stabilizer
denominators. Designed families are those for which these choices have been
made. The next theorems show that designed families exist and that the links
hold eventually along them.

% before: asm-duty
Several constants of a family take the same value throughout the proof, and
the ranks must be adapted to the scale. These requirements are collected as
the duty data of a family. They fix in particular the exponent $401$ that
bounds the twist parameter $t$.

% before: asm-sha-bound
The local errors at the twisted places are controlled by Tate-Shafarevich
groups of Tate twists over the host field. The bound needed is the one of
Chapter~\ref{chapter-tate-shafarevich}, with $e^{x^{\varepsilon}/2}$ on the
right. It is imposed on every member of a family at every large scale.

% before: asm-twisted-links
The division link of Chapter~\ref{chapter-integral-structure} is proved under
the division bound at twisted places, which concerns the local defects at the
places where the tuple of Hecke characters is twisted. The Tate-Shafarevich
link bounds these defects by the Tate-Shafarevich sum plus a term of size
$x^{\varepsilon}B_0$. The rank conditions are used by the slope link and by
the frame inequality.

% before: asm-global-inputs
The comparison inequality uses five links at once: gains at inert primes,
near-orthogonality, full row space, division and slopes. The global link is
the statement that these, with the chosen constants, tie the family. It is
proved, with $C_{\mathrm{abs}}=2$, $C_0=3$, $e_D=2$ and $e_L=0$, in
Theorem~\ref{asm-global-link}.

% before: asm-division-from-sha
The bound of Chapter~\ref{chapter-tate-shafarevich} enters the comparison
here. The Tate-Shafarevich sum is at most $\sqrt b$, which is smaller than
$B_0=fmbL$. So the Tate-Shafarevich link turns into a bound on the twisted
defects of the same size as the other error terms.

% before: asm-slope-link
The slope link compares the division lattice $N_u$ with the reference lattice
of the size data. The two differ by an ideal built from the accumulated local
gains, rounded down. The rounding costs at most the logarithms of the primes
involved, and Chebyshev's bound controls their sum.

% before: asm-tied-from-links
The links are proved separately, each eventually along the designed families
of a class. This theorem combines them at one family. The Tate-Shafarevich
bound is assumed on a larger class, the one in which the host requirement
holds, and designed families are found in the intersection.

% before: asm-member-bounds
The comparison inequality of Chapter~\ref{chapter-inert-prime-balance} has
eleven hypotheses and two sets of windows. At a member of a family they become
nine bounds on its invariants: the size of $\log T_u$, the slopes of $M_u$ and
$N_u$, the lattices, the frame, the losses, the weights, the characters and
the windows. Each is proved from its own inputs, and
Lemma~\ref{asm-tie-from-bounds} assembles them into a tie.

% before: asm-tie-from-bounds
Once a family satisfies the nine bounds at all large scales, the tie is a
matter of choosing members and reading off comparison data. The members are
chosen one for each base point, with common exponents, and the constants of
the comparison data are the constants of the bounds and of the division link.

% before: asm-facts-volume-tie
The count of Definition~\ref{asm-member-bounds} uses the volume that a
definite-side datum attaches to its base space, and the proof of the global
link compares this volume with a mass. The comparison uses three properties of
the datum: volumes are measured by fundamental domains for normalized Haar
measures, cutting off a vector of norm $\pm1$ divides the volume by the factor
of Siegel's mass ratio, and the mass formula holds in rank $2$.

% before: asm-facts-base-volume
In rank $2$ the identity between volume and mass comes from unfolding a
fundamental domain over the finitely many classes in a genus. Siegel's mass
ratio (Theorem~\ref{cv-siegel-mass-ratio}) then carries it to every rank,
first through the definite spaces and then one negative direction at a time.

% before: asm-bounds-glued-rank
The lattice and frame bounds need the glued lattice to be nonzero. Its rank is
the number of good rows, and the orthogonality link keeps the space of rows
nonzero: if it were zero, the diagonal Gram matrices would vanish, which the
bound on $\|K_{uu}-1\|_{HS}$ forbids.

% before: asm-bounds-dual-slopes
The lattice conditions bound the minimal slope of the dual of the glued lattice
by its maximal slope. Under duality, the saturated sublattices of the dual
correspond to the saturated sublattices of the lattice, by taking annihilators,
and the Arakelov degree changes sign; this exchanges the two extreme slopes.

% before: asm-bounds-lattice
The lattice conditions of Definition~\ref{asm-member-bounds} are two general
facts about Hermitian lattices. They apply to the division lattices and to the
glued lattice once these are finitely generated, spanning and nonzero.

% before: asm-bounds-frame-realised
Proposition~\ref{hs-frame-at-glued-lattice} needs a frame realisation:
coordinates in which the division pairings become blocks of the nearly
orthogonal maps $P_u$ of the orthogonality link. The Bessel identity provides
these coordinates, and the ramified part of each $M_u$ is the part of the frame
that $N_u$ does not see.

% before: asm-bounds-frame
The parameters of the orthogonality link are of size $\exp(-x^{\varepsilon/8})$,
far below $x^{-C_H}$, and the ramified part of $M_u$ has relative size at most
$2f/b$ because $T_u\ge b^{m/2}$. So a frame realisation with $\sigma=x^{-C_H}$
exists at every large scale, and Proposition~\ref{hs-frame-at-glued-lattice}
turns it into the frame bound.

% before: asm-bounds-nodes-characters
The node bounds and the character conditions need no input from the links. The
first follow from the windows in which the size data place the Harish-Chandra
weights, and the second from attaching to each character of the host group the
logarithmic derivative at $1$ of a primitive Dirichlet character.

% before: asm-bounds-windows
The window conditions are inequalities between the ranks, the scale and the
host data of a member. They follow from the sizes of the family once a few
powers of $x$ lie below the corresponding exponentials and the set $T_2$ has
density at least $x^{-\varepsilon}$.

% before: asm-level-setting
The rank conditions and the slope bounds for $N$ rest on local data at the
auxiliary prime $T_u$ of each word: the dimension and the slopes of the ramified
part of $\Lambda_{u,0}$, and the colength of the local lattices
$N_\lambda(u,w)$ in the ramified local span. The setting of the level primes is
the condition on the size of $T_u$ under which these data are proved.

% before: asm-level-setting-holds
The size condition on $T_u$ puts $T_u$ above $b^{m/2}$, while the weights of
the forms, their spread and the ranks are bounded by small multiples of $b^2$.
So the setting of the level primes is a consequence of the sizes of the family.

% before: asm-level-local-facts
Five local statements at $T_u$ feed the level data. They concern the action of
the groups $\mathbf F_v^\times$ on $M_u$, the slope of the ramified lattice,
and the Hecke operators $U_v$ together with down operators $D_v$ that invert
them on the ramified part up to a power of $T_u$.

% before: asm-level-free-count
When $M_u$ is the space of functions on a free orbit set of the group
$\prod_v\mathbf F_v^\times$, its ramified part is the space of functions whose
sums over the orbits of each factor vanish. Its dimension is then a product of
local counts, one factor $(T-2)/(T-1)$ for each field.

% before: asm-level-division-side
The division side conditions follow from the free basis and the slope line.
The dimension comes from the count of functions on a free orbit set, and the
maximal slope can only decrease on passing to the ramified sublattice once
that sublattice is nonzero.

% before: asm-level-congruence-dvr
The length of $S^{\mathrm{ram}}_w/N_\lambda(u,w)$ is controlled one Hecke
operator at a time. Over a discrete valuation ring, an operator $U$ that has a
partner $D$ with $UD=\varpi^e$ moves a lattice by at most $e$ times its rank,
plus the colength of the lattice through which it factors; a saturated
subspace preserved by both operators inherits the bound.

% before: asm-level-congruence
At a place $w$ over $T_u$ the local lattice $N_\lambda(u,w)$ is the image of
the ramified local span under the $m$-th power of a product of scaled Hecke
operators. The second branch inclusion supplies the lattice through which each
operator factors, and the ramified level zero supplies its partner, so the
lemma over a discrete valuation ring bounds each factor.

% before: asm-level-commensurable
The division lattice $N_u$ is built from $N^\circ_u$ by a fractional ideal,
and $N^\circ_u$ is prescribed place by place relative to $M^{\mathrm{ram}}_u$.
Local inclusions at every place give a global inclusion, and local equality at
a single place gives a common multiplier; so $N^\circ_u$ and $N_u$ are lattices
in $V(N_u)$.

% before: asm-global-second-layer
The local facts at $T_u$ and the values of the columns are drawn from the tied
columns of a member. The second-layer conditions at the real place supply these
columns at all large scales. The stabilizers of the classes in degree $0$ are
trivial once a level prime exceeds $2fk+1$, which the sizes of the family
guarantee.

% before: asm-global-column-values
The lower and upper bounds for the slopes of $M_u$ come from coordinates on the
forms of $u$ in degree $0$. Each form becomes a family of polynomials in the
Weyl modules of its weights, one for each class, and the coordinates describe
its integrality place by place and its norms at the infinite places.

% before: asm-global-count
The auxiliary prime $T_u$ is pinned between $T^*_u$ and $(1+1/(fk))T^*_u$, and
$T^*_u$ is $X_u$ divided by a power of $\Gamma_P$ that is at most $e^2$. So
bounds for $\log X_u$ and for its variation between words become the three
conditions of the count, at the cost of a constant.

% before: asm-global-numerical
The column bounds for the slopes of $M_u$ hold under explicit conditions on the
sizes of $b$, $b'$, the cloud size, the width and $\log T_u$, and on the
entries of the weights of $u$. One computation checks these conditions for the
size data of every member at a large scale, and the bounds are then applied
with them.

% before: asm-global-numerical-hold
The size data of a member satisfy the numerical and weight conditions as soon
as $\log T_u$ and the logarithmic derivative at $1$ are controlled. The
parameter choices place the weights in the required windows, and the sizes of
the family supply the remaining inequalities at large scales.

% before: asm-global-slopes-m
In the column slope bounds every term is a product of a size of the member
with a logarithm. The sizes of the family put each size below $|T|a$, and each
logarithm below a multiple of $L_a$, so the bounds collapse to multiples of
$|T|aL_a$.

% before: asm-global-losses
The losses are bounded at every member whose weight vectors have small shares.
The loss bounds ask for one choice of exponents at each scale, realized by some
member, such that every member with these exponents has bounded losses. Since
every member has small shares at large scales, the exponents of any member
serve.

% before: asm-global-slope-index
The slope bounds for $N$ compare $N_u$ with $M^{\mathrm{ram}}_u$, which differ by
finitely many local indices and by the fractional ideal $I_u$. The slope of a
sublattice of full rank is the slope of the lattice less a weighted sum of these
local lengths, and multiplying by a fractional ideal shifts both the slope and
the maximal slope by its norm.

% before: asm-global-unit-weight
The local data enter the slope of $N_u$ through the rounded values
$\lfloor e_ws(u,w)\rfloor$, the depths $\tau_w$ and the local lengths. The
following computation collects them into the gains, local terms and errors of
the member bounds, and bounds what remains by the index excesses, the lengths at
$T_u$, the rounding and the deficit of the ramified part.

% before: asm-global-level-t
At the places over the auxiliary prime $T_u$ the gain link says nothing, and the
level data take its place. The division side conditions give the rank and the
slopes of $M^{\mathrm{ram}}_u$, and the congruence length bound, with the
weight conditions, bounds the local lengths of $M^{\mathrm{ram}}_u/N^\circ_u$
at these places.

% before: asm-global-slopes-n-member
At one member, the slopes of $N_u$ and of $M^{\mathrm{ram}}_u$ differ by the
rounded local data and the local lengths, and the maximal slopes by the rounded
data less the depths. Lemma~\ref{asm-global-unit-weight} turns these two
comparisons into the slope bounds for $N$.

% before: asm-global-slopes-n
For a family, the hypotheses of Lemma~\ref{asm-global-slopes-n-member} come
from the choices of the family, the level data, the links and the sizes. The
only further input is a bound on the depths, which the size $b^3\le T_u$ of
the auxiliary primes turns into $\sum_w\mathrm{sz}_w(\tau_w)\le b^2$.

% before: asm-global-link-column-values-local-basis
Membership in the lattices of forms at a finite place $w$ is tested on bases of
the lattices $L_u(0,c)$ after inverting the integers prime to the residue
characteristic of $w$. Such a basis is the matrix of a basis of the localized
lattice of coordinate vectors.

% before: asm-global-link-column-values-local-basis-facts
Local bases exist because a Dedekind domain with finitely many primes is a
principal ideal domain. The tests use the dual lattices, whose local bases come
from those of the lattices and the Gram matrix, and they compare two local
bases through a matrix that is integral at the places over $\ell$.

% before: asm-global-link-column-values-weyl-stable
Changing the local basis acts on the polynomials of the forms by a matrix that
is invertible over the local ring at $w$. The lattice spanned by the integral
Weyl modules is stable under such matrices, since over a local ring they are
products of elementary row operations.

% before: asm-global-link-column-values-criterion
The tied forms in degree $0$ describe $\Lambda_{u,0}$ by integrality conditions
at all finite places at once. The localization of $\Lambda_{u,0}$ at one place
$w$ is described by the conditions at $w$ alone, since a denominator of a
vector can be cleared at the other places.

% before: asm-global-link-column-values
The slope bounds for $M$ need the columns of every word to have values. These
come from tied columns at a datum whose classes in degree $0$ have trivial
stabilizers, which the second-layer conditions provide at every member at
every large scale.

% before: asm-global-link-local-facts-free
The first local fact identifies $M_u$ with the functions on a set on which the
group $\prod_v\mathbf F_v^\times$ acts freely. The auxiliary prime $T_u$ is a
norm from $K$, so $E$ splits at the places over $T_u$, and this splitting,
with the trivial stabilizers, gives such a set.

% before: asm-global-link-local-facts-slope
The slope line compares the slopes of $\Lambda^{\mathrm{ram}}_u$ and
$\Lambda_{u,0}$. When the character of the level has finite order, the diamond
operators preserve the metrics at the infinite places, and the comparison
reduces to a count of the ramified part.

% before: asm-global-link-local-facts-integral
Hecke integrality at level zero is a statement at each place over $T_u$. The
lattices $L_u(0,c)$ change only at one place when the representative changes
at that place, which reduces the statement to the Hecke operators at a single
place.

% before: asm-global-link-local-facts-ramified
The ramified level zero says that $U_v$ and $D_v$ preserve the ramified part
and compose to a power of $T_u$ there. Both follow from an explicit formula for
$D_vU_v$ on the forms in degree $0$ and from the way $U_v$ and $D_v$ interact
with the diamond operators at $v$.

% before: asm-global-link-local-facts-branch
The second branch inclusion compares two lattices at a place over $T_u$. In
local coordinates both are described class by class, and the length of the
quotient is then a sum over the classes, which the dimension of $M_u$ counts.

% before: asm-global-link-local-facts
The level data rest on five local facts at the auxiliary prime of a word. They
follow from the forms in degree $0$ that tied columns give at the word, with
more hypotheses for the later facts: the setting of the level primes and a
character of finite order for the slope line, and trivial stabilizers for the
other three.

% before: asm-global-link-count-bounds-weights
The bounds on the auxiliary primes are proved for any weights and
Harish-Chandra parameters that satisfy a list of size conditions. The list
records the shape of the parameters of the construction: equally spaced
entries $v_i$, entries $u_i$ next to them, and weights $w_j$ just above the
entries $b_j$, which are far apart.

% before: asm-global-link-count-bounds-weights-root-ratios
The first size condition asks that the half root discriminant $v$ of the host
field be at least $\pi$, and the host field is a subfield of a cyclotomic field.
The logarithm of the root discriminant of a subfield of $\mathbf Q(\zeta_N)$ is a
sum over the prime powers dividing $N$, each weighted by the proportion of a
congruence subgroup of $(\mathbf Z/N\mathbf Z)^\times$ that lies outside the
subgroup fixing the subfield.

% before: asm-global-link-count-bounds-weights-root-ratio-facts
The outside ratio at $f$ is the proportion of $K_N(f)$ that lies outside $U$. It
is at least one half once $K_N(f)$ is not contained in $U$, because a single
element outside $U$ moves $K_N(f)\cap U$ into the complement.

% before: asm-global-link-count-bounds-weights-root-weights
The ramification sum can also be written as a sum over the elements outside
$U$, each weighted through the additive order of $u-1$. The two forms agree
because the weights of the multiples of a prime power dividing $N$ add up to a
single term.

% before: asm-global-link-count-bounds-weights-root-cyclotomic
The norm of $\zeta-\zeta^a$ in the cyclotomic field depends only on the order $n$
of $\zeta^{1-a}$. It is given by the value at $1$ of the $n$-th cyclotomic
polynomial, which is $p$ when $n$ is a power of a prime $p$ and $\pm1$ when $n>1$
is not a prime power.

% before: asm-global-link-count-bounds-weights-root-different
The discriminant of a subfield $E$ of $\mathbf Q(\zeta_N)$ can be read off from
differents. The ring of integers of $\mathbf Q(\zeta_N)$ is generated by $\zeta_N$
over that of $E$, so the different of $\mathbf Q(\zeta_N)/E$ is generated by the
derivative of the minimal polynomial of $\zeta_N$. Comparing it with the
different over $\mathbf Q$ leaves the conjugates of $\zeta_N$ that move $E$.

% before: asm-global-link-count-bounds-weights-root-discriminant
Together, the last two lemmas express the logarithm of the root discriminant of
$E$ through the subgroup $U$ of $(\mathbf Z/N\mathbf Z)^\times$ that fixes $E$.
Each element outside $U$ contributes through the additive order of $u-1$.

% before: asm-global-link-count-bounds-weights-root-two
The lower bound for the root discriminant of the host field uses two kinds of
divisors of $N$: the powers of $2$ and the prime $d$. At the powers of $2$ the
outside ratios are large, because the subgroup fixing the host field consists
of classes congruent to $\pm1$ modulo $2^{e+2}$.

% before: asm-global-link-count-bounds-weights-root-sum
Adding the contributions of the powers of $2$ and of $d$ bounds the
ramification sum from below. At $d$ the outside ratio is at least one half once
some class congruent to $1$ modulo $N/d$ lies outside $U$.

% before: asm-global-link-count-bounds-weights-root
For the host field, the subgroup fixing it consists of classes congruent to
$\pm1$ modulo $2^{e+2}$, and some class congruent to $1$ modulo $Q_Sq$ moves it.
Otherwise the second sign of $\operatorname{art}(u)$ would be trivial on the
classes congruent to $1$ modulo $Q_Sq$ and on those congruent to $1$ modulo $d$,
hence on all classes, while $\operatorname{art}(-1)=\iota$ has second sign $-1$.
The lower bound for the ramification sum then bounds the root
discriminant from below, and gives the first size condition.

% before: asm-global-link-count-bounds-weights-hold
The members of a family satisfy these size conditions at every scale $x\ge3$,
with the Harish-Chandra weights of the lifts. The conditions come from the
sizes of the family and from its choice of parameters.

% before: asm-global-link-count-bounds-explicit-rows
The bounds on $X_u$ use a few further inequalities between the parameters, the
level primes and the scale. The sizes of the family give them at every member
with weight data, once $x\ge44^4$.

% before: asm-global-link-count-bounds-explicit-log-x
The number $X_u$ is defined by the dimensions $D^*$ and $d^*_u$, which the
masses of the standard lattices, the base volume and the formal degrees of the
discrete series determine. Bounding these on both sides bounds $\log X_u$ below
by the weights and the level primes and above by a multiple of $mL$.

% before: asm-global-link-count-bounds-explicit-variation
The words differ only in the interior entries of their parameters, each moved
by at most $\Delta$. The dimension $d^*_u$ changes with the Vandermonde
determinants of these entries, so $\log X_u$ varies by about $\log k$ between
words.

% before: asm-global-link-count-bounds-explicit
Under the size conditions, the masses of the standard lattices, the volume of
the base space and the formal degrees of the discrete series determine $X_u$ up
to bounded factors. An auxiliary prime in the window around
$X_u\Gamma^{-2/(fk)}$ is then bounded on both sides, and $X_u$ varies little
between words.

% before: asm-global-link-count-bounds
The count of Definition~\ref{asm-member-bounds} follows from the three bounds
of Lemma~\ref{asm-global-count} on the auxiliary primes. These bounds are the
explicit form of the Vandermonde count, which rests on the masses of the
standard lattices, the volume of the base space and the formal degrees of the
discrete series.

% before: asm-global-link-depth-sum
Lemma~\ref{asm-global-slopes-n} asks for a bound on the depths of the size
data. When the cuts and the Weyl exponents of the gain data are small at every
prime below $X'$, the depth at each such prime is small, and Mertens' bound
adds these bounds up to a multiple of $\log X'$ plus a term of size about
$X'\log a$.

% before: asm-global-link-depth-sizes
The cuts and the numbers $D(p)$ are sums of quantities built from the ranks,
the scale and the orbit counts. All of them are at most a power of one number
$S$, and the slack $b'\ge S^{20}$ absorbs these powers. The definition lists
the bounds by $S$ that the estimates below use.

% before: asm-global-link-depth-sizes-hold
For a family the natural choice is $S=\exp(x^{\varepsilon/2})$. The ranks are
powers of $x$, the orbit counts are multiples of $\log x$ by
Chapter~\ref{chapter-period-divisibility}, and $b'$ is about
$\exp(x^{\varepsilon})$, so every bound holds once $x$ is large.

% before: asm-global-link-depth-representative
A served representative $t_0$ at a prime $p<X'$ has a tiled set of positions,
whose size $N(p,t_0)$ is comparable with $n_p\delta_1$ up to a factor
$C_e\log x$. The cut bounds compare everything with $N(p,t_0)a/p$, which
therefore is at least $1$ and dominates the terms $n_pc_G$.

% before: asm-global-link-depth-cuts
The cuts and the numbers $D(p)$ are given by the formulas of
Lemmas~\ref{pd-gain-second-hypotheses} and~\ref{pd-gain-first-hypotheses}.
With the bounds by $S$, each of their terms is at most a constant times
$N(p,t_0)a/p$, the constant depending only on $\varepsilon$ and the constants
of the gain data.

% before: asm-global-link-depth
Every family has small cuts and Weyl exponents at all large scales. The term of
size about $X'\log a$ left by the previous lemma is then at most
$\exp(8x^{\varepsilon/2})$, a power of $b$ whose exponent tends to $0$.

% before: asm-global-link-bad-primes
The bound on the bad tops in the chosen constants carries a factor
$\sum_p\log p/p$ over primes below $X'$. Mertens' bound makes this factor at
most $2x^{\varepsilon/3}$, which turns the bound into one of size
$|T|x^{-2+\varepsilon}$.

% before: asm-global-link-losses-numbers
The inert-prime sums are compared with the gains prime by prime. The errors of
that comparison are explicit expressions in a few numbers of the member, and
the primes fall into four classes according to whether they are dropped and,
if so, why.

% before: asm-global-link-losses-terms
The losses are sums of terms $\max(a_t/p^i-w_i,0)$ and of the lower and floor
terms of the gain link, over sets of primes. Each kind of sum is bounded by a
geometric series or by the bounds of Chebyshev and Mertens for sums over
primes.

% before: asm-global-link-losses-explicit
At a kept prime the gain link bounds the inert-prime sum by the gain, up to the
lower term and the floor term. At a dropped prime the gain is zero, and the
whole inert-prime sum at that prime is a loss. Adding over the primes gives the
following bound, in which three sums over primes remain.

% before: asm-global-link-losses-orders
The primes of the first kind include primes whose class modulo a power of $2$
has small order. Such a class is $\pm1$ modulo a smaller power of $2$, by the
structure of the group of units modulo a power of $2$.

% before: asm-global-link-losses-residues
A sum of $\log l/l$ over integers in two residue classes modulo $M$ is smaller
than the full sum by a factor of order $M$. This is the bound used for the
primes of small order.

% before: asm-global-link-losses-band
The other primes of the first kind have a power close to $X'$. Such a prime is
close to a root of $X'$, and the primes in a short interval around $P$
contribute little to $\sum\log p/p$.

% before: asm-global-link-losses-first
The two previous lemmas bound the sum of $\log p/p$ over the primes of the
first kind. The bound is small because the modulus $2^{e+2}$ is at least
$4x^2$ and the bands around the roots of $X'$ are short.

% before: asm-global-link-losses-sizes
The error terms of the explicit bound involve powers of $x$, the quantities
$\exp(x^{\varepsilon/4})$ and $X'$, and one of $a$, $b'$ or $w''$. On the
windows that the sizes of the family provide, each is at most
$\mathrm{sg}\,|T|mb$ with a factor $\mathrm{sg}\le1$, except one term of size
about $2\log4\cdot C_P|T|mb$.

% before: asm-global-link-losses
The loss bounds are needed only at the members whose weight vectors have small
shares, since the duty data provide such shares at large scales. At these
members they follow from the bound on $|T|$, the bound on the bad tops and the
gain link.

% before: asm-global-link
The global link is the last input of Theorem~\ref{asm-tied-from-links}. Its
proof derives the bounds of Definition~\ref{asm-member-bounds} from the slope
inequalities of Chapter~\ref{chapter-hermitian-slopes}, the masses of
Chapter~\ref{chapter-covolumes} and the other links, and applies
Lemma~\ref{asm-tie-from-bounds}. The resulting tie has $e_D=2$ and $e_L=0$,
which produce the exponent $1/2+9\varepsilon$.

% before: asm-linked-aligned
The proviso of the chosen constants compares $\dim M_u$ with $D^*/\Gamma^2$.
It is proved from the tied data that a twist record contains, so it applies to
families whose members carry tied data at all large scales. The duty data add
the fixed values, the bound on the first modulus and distinct tame data, and a
family with all of these is linked and aligned.

% before: asm-aligned-second-layer
A twist selection chooses a twist record at every member beyond a threshold,
and every twist record contains tied data. The remaining conditions are part
of the duty data. So the classes of Theorem~\ref{asm-design-exists} consist of
linked and aligned families.

% before: asm-proviso-splitting
The class number formula of Proposition~\ref{cv-class-number-formula} applies
only when $g_0$, $g_1$ and $T_u$ split completely in $F$. For a definite-side
datum this follows from its host data and from the congruence that $T_u$
satisfies.

% before: asm-proviso-trivial-stabilizers
The space $M_u$ has the dimension of a product over the classes at the bottom
of the chain only when those classes have trivial stabilizers. A nontrivial
stabilizer would give a unitary matrix of finite order whose order is a power
of a level prime. When the level prime exceeds $2fk+1$, a root of unity of that
order has too large a degree to be an eigenvalue of such a matrix.

% before: asm-proviso-bottom-dimension
The tied columns identify $M_u$ with a space of algebraic modular forms on the
definite unitary group of the bottom space. With trivial stabilizers, a form is
determined by its values at representatives of the classes, and these values
can be chosen freely. Weyl's dimension formula then gives the dimension at the
infinite places.

% before: asm-proviso-window
The rank formula gives $\dim M_u=d^*_u(T_u^k-1)^f$, and the auxiliary prime
$T_u$ lies in a window just above $(D^*/d^*_u)^{1/(fk)}\Gamma^{-2/(fk)}$. Since
$(1+1/(fk))^{fk}\le e$, this window fixes $\dim M_u$ up to a factor $e$ above
and a factor $2$ below.

% before: asm-proviso-reduction
The splitting of the level primes, the trivial stabilizers, the dimension of
$M_u$ and the class number formula combine into the rank formula
$\dim M_u=d^*_u(T_u^k-1)^f$ at every member at a large scale. The window for
$T_u$ then gives the proviso with any constant $C_0\ge e$, and
Theorem~\ref{asm-design-exists} uses it with $C_0=3$.

% before: asm-design-exists
Theorem~\ref{asm-tied-from-links} needs designed families with arbitrarily
large constants. They come from design choices, which a collection of
definite-side hypotheses provides once Minkowski's bound on finite subgroups
of $GL_n(\mathbf Q)$ is known. The remaining chosen constants come from the
class number formula and from the zero-density estimates of
Chapter~\ref{chapter-host-fields}.

% before: asm-host-requirement
The families must live over the host fields of
Chapter~\ref{chapter-host-fields}, which exist only on windows of scales. A
host requirement records what is asked of the host at each member. For a
primitive character $\psi$, it asks for the host conditions of the cyclic
level of $\psi$ and for the Tate-Shafarevich bound.

% before: asm-tied-everywhere
The tie is proved once for all host requirements, and then applied to the
requirement of each primitive character. For a given requirement it needs two
inputs: hosts on windows of scales, and the Tate-Shafarevich bound for the
families built over them.

% before: asm-definite-side-hypotheses
The links of Chapters~\ref{chapter-period-divisibility}
and~\ref{chapter-integral-structure} are proved from inputs that concern a
fixed class of families and the selections made at its members. A collection
of definite-side hypotheses gathers these inputs. Its existence is proved from
published results in Theorem~\ref{asm-definite-side-exists}.

% before: asm-support-local-modules
The support condition of Chapter~\ref{chapter-integral-structure} compares, at
each twisted place $w$ and lift $j$, the ranks of the local modules of a member
with the dimension of its rational forms for $\mu_w$ at level $j$. A local
module is a free module over a complete discrete valuation ring with a faithful
action of a local algebra of Hecke functions, and the rational forms project to
it.

% before: asm-support-hecke-realization
The spaces of forms of a member are realized as algebraic modular forms on a
definite unitary group, with a level and a character. Through this realization
the normalized Hecke operators at split places act on the forms, and their
commutativity and the integrality of their eigenvalues give the support
condition.

% before: asm-support-local-compatibility
Normalized local compatibility attaches to each level and place a Hecke
algebra over a discrete valuation ring, acting on a space that contains the
rational forms. Its maximal ideals correspond to the local indices at that
place, and each local module fills the part of the space at its maximal ideal.
The Hecke ties make the level of the realization unramified away from a small
set of places, where the Hecke operators are standard.

% before: asm-support-hecke-standard
At a split place where the level is unramified, the normalized Hecke operators
preserve the type space, commute, and have eigenvalues that are integral away
from the residue characteristic of the place. These are standard properties of
Hecke operators on algebraic modular forms on definite unitary groups.

% before: asm-support-eigenspaces
A commuting family of split operators with eigenvalues in a local ring
decomposes the space into joint generalized eigenspaces, one for each
homomorphism to the ring of the algebra they generate. Grouping these by the
maximal ideal of each homomorphism gives parts whose dimensions add up to the
dimension of the space.

% before: asm-support-from-ties
With these ingredients the support condition follows at a single member. The
Hecke operators make the generators of each Hecke algebra commute and have
integral eigenvalues, and the decomposition into parts then counts the ranks of
the local modules.

% before: asm-rank-from-ties
The rank equality at small level asks that the space of lifts at the base tuple
have the number of base rows as its dimension. The rows and the projections of
the selected base level reach every basis vector of the lifted space, so the
theta lift at the base tuple is onto and its adjoint is injective. The
dimension of the lifted space is then read off from the parts of the rows and
the counts of the types.

% before: asm-primitive-rows-parts
Lemma~\ref{asm-rank-from-ties} asks that the parts of the definite side at the
base tuple be independent and span. These parts are cut out by the types of the
Hecke eigenvalues, and two types of the base are told apart by some place or
inert index. On the lift side, the lift-distinct types do the same for the
transferred types.

% before: asm-primitive-rows-base
The selection bundle gives the selected base level only through a realization
with primitive rows. The theta lifts of the realization carry the definite side
onto the parts of the top space and of the lifted spaces, and this determines
the matrices of the rows, their parts and the dimensions of the parts. These
are the conditions of hypothesis (2) of Lemma~\ref{asm-rank-from-ties}.

% before: asm-division-kernel-glued
Part (9) of the division kernel bounds asks that the glued lattice of the
member be a lattice. The division rows identify it with the glued lattice of
the columns, which is assembled place by place from finitely many lattices.

% before: asm-division-kernel-errors
Part (6) of the division kernel bounds compares the local errors of the size
data with the exponents $a_t$ at the level primes, and bounds them elsewhere.
For the families of the argument these errors are given by explicit formulas
in the row exponents, the weights and the stabilizer denominators, and the
collection asks for eight bounds on the terms of these formulas, the second
complement inputs.

% before: asm-division-kernel-complement
At a place over a level prime $g_i$ the explicit error is a multiple of the
weight of the place plus a valuation of $\pi_i$, and the places over $g_i$
together carry the weight $\log g_i$. At the other places it is bounded prime
by prime. The second complement inputs then give part (6).

% before: asm-division-kernel-section
Part (8) of the division kernel bounds asks that the Bessel pairings take
values in the coefficient field on the division lattices. The pairings are
tied to the periods of the iterated lift through a section of the map from the
base space to the top forms, and the rationality is read off from this tie.
The section inverts that map on the orthogonal complement of its kernel.

% before: asm-division-kernel-descent
Suppose that a period, on a set of vectors spanning the base space, takes the
values of a linear form over the coefficient field read through the rows map.
Then its composite with the section takes the same values on the division
lattice. The linear form extends to the top forms through the division
lattice, and the section returns to the base space without changing the
period.

% before: asm-division-kernel-multiplicity
The periods of the definite columns are compared with the top space through
the realization. Its multiplicity condition writes each period pairing as an
inner product of a theta lift with a Poincaré vector, up to unitary changes of
variables and a scalar for each word and tuple. The column data of the
realization are built from the periods at the base tuple, with scalars chosen
to give the factorization of a division package.

% before: asm-division-kernel-evaluation
The multiplicity condition evaluates the theta lifts on the row forms with
given coordinates, and the column data use the coordinates of the definite
side. The two evaluations agree: their difference is orthogonal to the
definite subspace, and the adjoint of the theta lift carries the level of the
top space into the span of that subspace.

% before: asm-division-kernel-factorization
With these evaluations, the multiplicity condition at the base tuple computes
the pairings of the column data, and at the other tuples it relates them to
the period pairings. A scalar of the multiplicity condition that lies in the
coefficient field becomes a scalar of the factorization; when it does not, the
period pairing vanishes. The isometry of the row space gives the isometry of
the column data.

% before: asm-division-kernel-rational
The rationality of the Bessel pairings follows from the Bessel tie. The
division rows identify the division pairings with the pairings of the column
data, the compatibility with the periods identifies the periods of the
definite columns with those of the datum at the base tuple, and the descent
lemma carries these identities through the section.

% before: asm-division-kernel-join
Part (7) of the division kernel bounds compares the Bessel pairings with the
standard Bessel periods of the automorphic setting. The join of the indefinite
realization ties its theta lifts, its columns and its Poincaré vectors to the
datum and the setting. With the multiplicity condition it writes the periods
on the base space as inner products with Poincaré series.

% before: asm-division-kernel-bessel
When every top form comes from the base space, the kernel join and the Bessel
tie express each Bessel pairing as an inner product with a Poincaré series.
The dictionaries of the automorphic setting carry these inner products back
to the Bessel-period data of the datum, which define the standard Bessel
periods.

% before: asm-division-kernel-used
Items (3)(e) and (3)(f) of the collection ask for the division rows and for
the compatibility with the periods only at the used tuples of the selected
base level. The indefinite realization of the selection bundle attaches every
tuple to the base or to a place, so both hold at every tuple. It also
identifies the rank of the rows with the dimension of the base rows.

% before: asm-division-kernel-package-passage
A division package must bound the glued lattice of its columns by that of its
rows. The rows and the columns of a realization are the same complex data cut
by the idempotents of two subgroups with characters, and the near-constant
multiplicities of the second in the representations of the passage data
control the loss between them.

% before: asm-division-kernel-package-situation
The rows and the columns of a division package carry metrics at the infinite places and scalars
by which their matrices change these metrics. Both are read off from complex data at one level, so
the next three statements share that setting.

% before: asm-division-kernel-package-rows
The metric of the glued lattice is the Gram metric of the theta lifts at the base place, and a sum
of rescaled pieces of the quotient metric elsewhere. The scalar conditions of a division package
then follow from the norms of the theta lifts on the parts and from the orthogonality of the parts.

% before: asm-division-kernel-package-lifts
The lifted spaces carry the Gram metrics of their coordinate vectors at every infinite place. At the
base place the scalar of a lift is its norm on a part; at the other places it is given directly by
the conjugate scalars of the primitive rows.

% before: asm-division-kernel-package-hecke
A division package also asks for operators over the coefficient field whose joint eigenspaces are
its parts. They come from spectral operators on the complex space, which are rational because the
parts are defined over the coefficient field.

% before: asm-division-kernel-package-normal
Hypothesis (6) of the lemma on the rows asks that distinct parts be orthogonal for the quotient
metric at every infinite place. Over the complex numbers the parts are sums of joint generalized
eigenspaces of commuting operators, and these sums, for disjoint sets of systems, are orthogonal for
any positive definite form for which the operators contain the adjoint of each of them. The quotient
metric inherits this through an operator that commutes with all of them.

% before: asm-division-kernel-package-orth
The orthogonality of the parts now follows at every infinite place. The Hecke and inert Hecke
operators are defined over the subfield on which the involution is complex conjugation, so after
any embedding they still contain the adjoint of each of them, and they descend to the space of
the level and to the quotient by the orthogonal complement of the definite subspace.

% before: asm-division-kernel-package
The division packages of part (1) of the division kernel bounds are built
from the rows and the columns of the realization of the selection bundle.
The definite rows give the numerical bounds that a division package asks of
its rows and columns, and the primitive rows give the Hecke structure on which
the package rests.

% before: asm-division-kernel-record-scale
At large scales the polynomial quantities of a member are small against $b$.
A power $\Pi^e$ with $e\le8$ is at most $X'$, and a fixed power of $x$ times a
constant is at most $b$ once $x$ exceeds a factorial depending on the exponent.

% before: asm-division-kernel-record-numbers
The record of a word collects the numbers of the member that the division
inequality uses. Most of its conditions compare a polynomial in the ranks with
$b$, $a$ or $w''$, and these follow from the sizes of the member at large
scales; the remaining facts concern the weights and the row exponents.

% before: asm-division-kernel-record-hook
Condition (7) of the record numbers compares the sum of $\phi$ over the row
exponents of a lift with its hook sum. The hook sum is the logarithm of the
product of the hook lengths of a Young diagram, and that product is at least
the product of the factorials of the row lengths.

% before: asm-division-kernel-record-phi
The comparison is made row by row, pairing a row exponent with a row of the
hook diagram of nearly the same length. Three elementary bounds for
$y\log y-y$ at integers control the pairs, the unpaired rows and the factorials.

% before: asm-division-kernel-record-shapes
The row exponents of a lift and the rows of its hook diagram interlace up to a
bounded shift. Two shapes describe this pattern: one at the places that are not
top places, and one at the top places, where the first row exponent carries the
exponent $a_t$ of the top class.

% before: asm-division-kernel-record-rows
For a pair of multisets of either shape, the pairing bounds the sum of $\phi$
over the row exponents by the factorial sum of the hook rows, up to a multiple
of the logarithmic size. The interlacing also fixes the total of the row
exponents up to a bounded error.

% before: asm-division-kernel-record-origin
The row exponents and the hook rows of a lift are given by the Harish-Chandra
parameters at an embedding. Their interlacing gives the shapes, once the
weight is shifted down by the lower bound that the sizes give for it.

% before: asm-division-kernel-record
The division inequality also needs record numbers for every word, and the
inputs of Definition~\ref{is-division-inputs} that are inequalities between
numbers. At large scales the sizes of the member give both, with constants
built from those of the family.

% before: asm-division-kernel-inputs-reduction
The covolume of the glued lattice of the rows is controlled place by place by
the Gram determinants of local bases. A Gram determinant of a projected lattice
is bounded once a fixed integer kills the dual of the lattice modulo the
lattice itself.

% before: asm-division-kernel-inputs-gram-spaces
The reduction lemma needs the Hermitian matrices, the projections and the
lattices of the definite side of the family to have the properties it assumes.
These come from the way the family presents its forms: as families of vectors
indexed by classes, fixed by finite stabilizers and paired by a class sum, with
integral and Weyl vectors whose duals are controlled by a level $s$.

% before: asm-division-kernel-inputs-gram-side
The definite side of the family at a level is cut out of the forms by an
element of the group algebra. Its Hermitian matrix is the class-sum pairing on
the image of that element, and its lattice is the set of integral forms there.

% before: asm-division-kernel-inputs-gram-idempotent
The element of the group algebra that cuts out the definite side is the
product of the central idempotent of a set of characters and the idempotent of
a character of a subgroup. Its coefficients lie in the subfield on which the
involution is complex conjugation everywhere, and become integral after
multiplication by the orders of the groups, because the values of the
characters are sums of roots of unity.

% before: asm-division-kernel-inputs-gram-projections
A self-adjoint idempotent of the group algebra acts on the forms as the
orthogonal projection onto its image, because the group preserves the
class-sum pairing. Its matrix is then given by the usual formula for an
orthogonal projection in terms of a basis of the image.

% before: asm-division-kernel-inputs-gram-duality
The lattice bounds compare duals of lattices on a space and on a subspace.
The dual of the part of a lattice in a subspace is the projection of the dual
of the lattice, when the projection is adjoint to the inclusion.

% before: asm-division-kernel-inputs-gram-hermitian
The Hermitian matrix of the definite side has entries in the subfield on which
the involution is complex conjugation everywhere, since the idempotent that cuts
it out does. It is positive definite at every infinite place because it is a
restriction of the class-sum pairing.

% before: asm-division-kernel-inputs-gram-lattice
The duality between integral and Weyl vectors at each class gives the bound
for the dual of the lattice of the definite side. Averaging over a stabilizer
turns the dual at one class into Weyl vectors, and the projection onto the
definite side loses at most the orders of the groups.

% before: asm-division-kernel-inputs-gram
The family conditions supply such an integer for the lattices of the rows: the
order of the group times the level. The reduction lemma then bounds the
valuation of each Gram determinant by the rank times the valuation of that
integer.

% before: asm-division-kernel-inputs-index-lifts
The lift forms enter through the forms of a level, the forms on which a
subgroup of the group acts by a character. Their lattice is the projection of
the integral level forms onto the definite subspace, and it is compared with the
integral and the Weyl forms that lie in that subspace.

% before: asm-division-kernel-inputs-index-general
For a lattice projected orthogonally onto a subspace, the relative index of the
projection and its dual splits into two parts. One is the index of the part of
the lattice in the subspace inside the projection, and the other is controlled by
an integer that moves the second lattice into the dual of the first.

% before: asm-division-kernel-inputs-index-classsum
The duality between integral and Weyl vectors at each class passes to the
forms. The class-sum weights are the orders of the stabilizers, so a common
multiple of them is the integer that is lost.

% before: asm-division-kernel-inputs-index-level
In the coordinates of the level forms, the projection onto the definite
subspace is an orthogonal projection for the restricted class-sum pairing. This
holds because the good part lies in the level forms, and the two projections
compose.

% before: asm-division-kernel-inputs-index
At the twisted places the covolume also depends on the relative index of each
lift lattice and its dual. The family conditions bound it by the local defect,
the lift exponent and the valuations of the stabilizer denominator and of the
level.

% before: asm-division-kernel-inputs
The other division inputs come from the rows condition of the definite-side
condition, which gives the chain of lifts with its height numbers and its
Cauchy-Binet data. The division bound at twisted places controls the
covolume of the glued lattice at the twisted places.

% before: asm-division-kernel-assembly
The package, its record numbers and its inputs live on the rows and the
columns of the realization. The division kernel data identify these with the
lattices, words and tuples of the member, and the three ties used in the
division link compare the numbers of the record with the heights, the losses
and the main term of the member.

% before: asm-division-kernel-sharp
Part (5) bounds the maximal slope of the glued lattice together with the
heights. The passage of a division package bounds the maximal slope of the
glued lattice of its columns by that of its rows, which the functional bound
controls, up to a loss of a multiple of the logarithm of the order of the
group.

% before: asm-division-kernel-mid
The definite rows hold for every large height constant, while the bound for
the glued lattice must be stated with an exponent fixed in advance by the
constants of the family. A property that holds eventually along a class can
be evaluated at a height constant that is a function of the exponent and of
the output constants.

% before: asm-division-kernel-functional
With a height constant fixed in this way the passage bound holds at every
member at all large scales. The logarithm of the order of the group is at
most a multiple of $b$, so the loss is absorbed by one more power of $x$.

% before: asm-division-kernel-explicit-heights
The heights of the division lattices of a member are explicit numbers built
from its size data. They combine the main term at the auxiliary primes, sums
of $\phi$ over the row exponents of the lifts, the logarithmic hook sums of
the weights of the lifts, and correction terms proportional to a constant.

% before: asm-division-kernel-explicit-entries
The bounds of part (5) are sums over the top classes, the lifts, the
embeddings and the rows of terms built from the exponents, the weights and the
other entries of the Harish-Chandra parameters, the row exponents and a few
logarithms. The windows of the size data bound each of these by a power of
$\Pi$ times $b$, and the number of summands is at most a power of $\Pi$.

% before: asm-division-kernel-sizes-local-terms
The local terms at twisted places are also explicit. They depend only on the
residue characteristic: at the level primes and at $d$ they are not positive,
and at an inert prime below $X''$ they are the values of the linear programme
of Chapter~\ref{chapter-period-divisibility} on the weights along an orbit.

% before: asm-division-kernel-sizes-mertens
The gains and the local terms at a prime $p$ are at most a constant over $p$,
and they are weighted by $\log p$. Mertens' first theorem bounds the resulting
sum over the primes below $X''$ by a multiple of $\log X''$.

% before: asm-division-kernel-sizes
The other terms of part (5) are bounded by powers of the scale times $b$. The
ranks, the twist parameter and $x^{\varepsilon}$ are at most powers of the
scale, and the gains and the local terms at twisted places are bounded in the
same way.

% before: asm-division-kernel-explicit-terms
The entry bounds give a bound for each kind of term in an explicit height.
The hook sums need one more input: the sequence of the $v_i$ followed by the
$-b_j$ drops by at least one at each step, so the arguments of the Gamma
factors are at least $2$ and the subtracted logarithms are nonnegative.

% before: asm-division-kernel-explicit
The explicit heights enter part (5) with a negative sign, so they need a lower
bound. Together with the mean of the logarithms of the archimedean scalars
they form the height of the record, which is bounded below by minus a power of
the scale times $b$.

% before: asm-division-kernel-first-lattice-slopes-member
The slopes of the first lattice of a word are first bounded at each member in
terms of the parameters of the word alone. The bounds involve, at each place,
the sum of the entries of the parameter, their excess over the last entry, and
the spread between the first and the last entry.

% before: asm-division-kernel-first-lattice-slopes-members
These bounds come from the slope inequality applied to the coordinates of the
columns of the word. The lower one also uses the exact slopes of the columns and
a bound for hook lengths.

% before: asm-division-kernel-first-lattice-slopes-chain
The numerical and weight conditions bound the entries of the parameters: at the
top places by multiples of $xb$ and $b'$, and elsewhere by multiples of the cloud
size. Inserting these bounds into the member bounds gives the column slope bounds
with constants that depend only on $C_n$.

% before: asm-division-kernel-first-lattice-slopes
The first lattice of the chain of a word is the lattice of its forms in degree
$0$, whose coordinates are the values of the columns. The slope inequality
bounds its slopes through the sizes of these coordinates, the exponents of the
top classes and the auxiliary prime.

% before: asm-division-kernel-first-lattice-arith
The bound on the maximal slope involves the number of top classes, the
quantities of the size data and two logarithms. Each is at most a power of
$\Pi$ times $b$, so the bound is a fourth power of $\Pi$ times $b$.

% before: asm-division-kernel-first-lattice
The slope bounds for the division lattices compare them with the first
lattice of the chain of each word. Its maximal slope is bounded through the
values of the columns, which the tied columns of the twist selection provide.

% before: asm-division-kernel-gauge
The gains, the local terms at twisted places and the local errors enter the
slope bounds for the division lattices as one combination. Part (6) bounds
the local errors from below, so the combination is bounded above by a power
of the scale times $b$.

% before: asm-division-kernel-height
These bounds combine into part (5). The slope bounds for the division lattices
reduce the slope of each division lattice to the first lattice of its chain
and to the local terms, and the quantity of part (5) averages these slopes
over the words.

% before: asm-division-kernel-finite-situation
The finiteness of the spectral data of an automorphic setting is a list of
separate facts about the same data. They concern one definite-side datum and
one setting attached to it, with its rows, column representatives, fundamental
domains, level subgroup and reductions.

% before: asm-division-kernel-finite-folding-domain
The unfolding of a Poincaré series needs a fundamental domain for the rational stabilizer of the Bessel
data. A discrete subgroup of a second countable group has one, cut out of countably many translates of a
small neighbourhood of $1$.

% before: asm-division-kernel-finite-folding-multiplicity
The square integrability of a Poincaré series rests on a finite multiplicity integral over a fundamental
domain of the stabilizer. Its finiteness combines the decay of the Bessel function with a polynomial
bound for the multiplicity and the structure of the nilradical.

% before: asm-division-kernel-finite-folding
The Poincaré series of a word are sums over the rational points of a
stabilizer, and their convergence is reduced to an integral over a fundamental
domain. This reduction needs the level subgroup in a normal form and the decay
of the archimedean Bessel functions.

% before: asm-division-kernel-finite-neat-level
Apart from $1$, no rational point fixing the hyperbolic pairs and the vector $z_u$ of a word may fix a
class of column representatives. This
comes from congruences modulo the product $g$ of the two level primes: modulo $g$ the adelic lattice
has a basis in which the form is standard, and the level subgroup moves its vectors only down a flag.

% before: asm-division-kernel-finite-neat-projection
Modulo $g$ the vectors of a level basis are dual to each other in pairs for the form, with $z^0$ dual
to $-z^0$. Subtracting from a vector of the lattice its components along them, computed with the
form, therefore leaves a multiple of $g$.

% before: asm-division-kernel-finite-neat-unipotent
An element of the level subgroup moves each vector of the flag spanned by the $e^0_a$ and the $b_i$
one step down, modulo $g$. On vectors whose pairings with $e^0_a$, $f^0_a$ and $z^0$ vanish modulo $g$,
the $k$-th power of $\kappa-1$ is then divisible by $g$, by duality with the $b'_i$.

% before: asm-division-kernel-finite-neat-orders
An automorphism of finite order of a lattice that is unipotent modulo a prime above $p$ has order a
power of $p$. Unipotence modulo primes above two different primes therefore forces the order to be
$1$.

% before: asm-division-kernel-finite-neat
The space of Poincaré data of a word is compared with the column space through
the values at the column representatives. When no rational point fixes a
class, these values are free and determine the datum.

% before: asm-division-kernel-finite-series-poincare
A Poincaré datum gives a function on the adelic group through the Bessel function at the infinite
places and the decomposition of the finite part along the unipotent group, the group fixing the flag
and its hyperbolic partners, and the level subgroup. Its Poincaré series sums the left translates by
representatives of the rational points modulo the rational stabilizer of the Bessel data.

% before: asm-division-kernel-finite-series-function-group-level
The topology of the finite adelic group is defined by its principal congruence subgroups. Its local
compactness, separation and countability, and the equality of Haar measures on the two sides, come from
the stabilizer of the adelic lattice and the indices of its conjugates.

% before: asm-division-kernel-finite-series-function-group
The analysis of Poincaré series uses a few properties of the adelic group of the setting. It is locally
compact and second countable with a discrete group of rational points, which is countable, and its Haar
measure is invariant under translations on both sides.

% before: asm-division-kernel-finite-series-function-bessel
The Poincaré function of a word is built from the unipotent groups along its isotropic flag, their
characters, the stabilizers of the flag and its hyperbolic partners, and the Bessel function at the
infinite places. A Poincaré datum transforms under these groups as the Bessel data prescribe.

% before: asm-division-kernel-finite-series-function-unipotent
The unipotent groups and their characters behave as expected. The finite unipotent group is a group on
which its character is multiplicative, and the stabilizer of the flag and its partners normalizes it
without changing the character.

% before: asm-division-kernel-finite-series-function-inputs-level
The comparison of the level character with the unipotent character uses the shape of the setting at the
level primes. Its lattice is built from powers of the different, its level subgroup and character come
from a level basis, and the base points match that basis modulo $g$ up to integers prime to $g$.

% before: asm-division-kernel-finite-series-function-inputs-trace
The finite part of the additive character is trivial on the inverse different. This is the duality of the
trace form, carried to the adeles by the product formula for the character.

% before: asm-division-kernel-finite-series-function-inputs-congruence
Membership in the level subgroup is a condition modulo $g$ on the level basis, while the unipotent group
is defined through the vectors $e_a$ and $f_a$. Clearing the integer by which the two bases differ turns a
lattice vector in the span of the $e_a$ into a combination of the $e^0_a$ modulo $g$.

% before: asm-division-kernel-finite-series-function-inputs-membership
A unipotent element of the stabilizer of the lattice, conjugated by the base point, moves each vector of
the level basis down the flag of the $e^0_a$ modulo $g$. So it lies in the level subgroup.

% before: asm-division-kernel-finite-series-function-inputs-levi
A rational isometry that moves each $e_a$ down the flag and $z_u$ into it is a unipotent element times a
Levi element fixing the hyperbolic pairs and $z_u$. The Levi element acts on the orthogonal complement of
these vectors by the projection of the isometry.

% before: asm-division-kernel-finite-series-function-inputs-transfer
The conditions defining the unipotent and stabilizer groups are linear conditions on the vectors $e_a$,
so they pass between the rational space and its completions. A rational vector in an adelic span of the
$e_a$ lies in their rational span, and the last condition on the unipotent part extends from rational
vectors to all vectors.

% before: asm-division-kernel-finite-series-function-inputs
Four further facts compare the unipotent group with the lattice, the level subgroup and the rational
points. A product of a unipotent element and a stabilizer element preserves the lattice only if both do,
and the archimedean and finite characters are inverse to each other on rational points.

% before: asm-division-kernel-finite-series-function
The Poincaré series inherits its properties from the Poincaré function. The function transforms
under the stabilizer, the level subgroup and the archimedean coefficient as the Bessel data and the
character prescribe.

% before: asm-division-kernel-finite-series-fold-multiplicity
The translates of a compact set by a discrete subgroup overlap only boundedly often. For a fixed base
point the count is bounded uniformly in the second point, and as a function of the base point it is
invariant under the subgroup and measurable.

% before: asm-division-kernel-finite-series-fold-unfold
Integrating a sum over representatives of the cosets of a subgroup against a fundamental domain of
the larger group is the same as integrating once against a fundamental domain of the subgroup. Almost
every point lies in exactly one translate of a fundamental domain.

% before: asm-division-kernel-finite-series-fold-bound
A function bounded by its averages over translates of a fixed set can be summed over the
representatives with a single integral. The translates of the set that meet a given point are then
counted by the multiplicity.

% before: asm-division-kernel-finite-series-fold
The convergence and the square integrability of a Poincaré series are general facts about a discrete
subgroup of a locally compact group. They need a bound of the function by its averages over a compact
set and a finite count of overlapping translates of that set.

% before: asm-division-kernel-finite-series-hecke-representative
On a fundamental domain, right translation is read through a representative of each orbit of the
rational points. Almost every orbit meets the domain once, a representative can be chosen measurably,
and right translation followed by the representative preserves the measure on the domain.

% before: asm-division-kernel-finite-series-hecke-fix
An operator given by integrating the right regular representation against an integrable function
acts on invariant functions by convolution. A function reproduced by the convolution is therefore fixed
by the operator.

% before: asm-division-kernel-finite-series-hecke-union
A sum over representatives of the cosets of a subgroup is invariant under the larger group. Its integral
over a fundamental domain of the larger group is the integral over a fundamental domain of the subgroup,
since the translates of the first domain by the representatives form such a domain.

% before: asm-division-kernel-finite-series-hecke
The space of automorphic forms is cut out of the square-integrable functions by a convolution at the
infinite places and an idempotent at the finite places. A Poincaré series lies in it once both
reproduce it.

% before: asm-division-kernel-finite-series
With the folding, the Poincaré series of every datum is an absolutely
convergent sum whose square is integrable on the fundamental domain. It is then
an element of the space of automorphic forms of the setting.

% before: asm-division-kernel-finite-dimension-situation
The finite dimensionality of the space of automorphic forms needs the matrix coefficients of the
archimedean discrete series to be integrable against a power of the operator norm. The next three
statements take place in an indefinite unitary group with its Haar measure and its polar
decomposition.

% before: asm-division-kernel-finite-dimension-decay-ktype-highest
The model of the lowest $K$-type is shown to be irreducible through its Lie algebra. In a unitary
representation of the Lie algebra of $U(p)\times U(q)$, the vectors obtained from a highest weight
vector by lowering operators span the subspace it generates, and this subspace has no other highest
weight vectors up to scalars.

% before: asm-division-kernel-finite-dimension-decay-ktype-derived
The representation of $K$ is passed to its Lie algebra by differentiation. Since $K$ is compact and
connected, every element is an exponential, so a subspace is stable under $K$ exactly when it is stable
under the derived operators.

% before: asm-division-kernel-finite-dimension-decay-ktype-model
A representation with a prescribed dominant weight is built from functions on $K$. A product of
powers of leading minors, one for each block, is a highest weight vector for right translation, and
its translates span a finite-dimensional unitary representation.

% before: asm-division-kernel-finite-dimension-decay-ktype
The matrix coefficients of the lowest $K$-type are compared with a function built from a model of
that $K$-type. The model is an irreducible representation of $K\cong U(p)\times U(q)$ whose highest
weight is the weight $\Lambda$ of the lowest $K$-type.

% before: asm-division-kernel-finite-dimension-decay-spherical
On the model of the lowest $K$-type there is an operator-valued function on $G$ that transforms
under $K$ on both sides as the model does and is an eigenfunction of the Casimir operator with the
eigenvalue of the discrete series. Under the gap condition it decays exponentially along the
noncompact directions.

% before: asm-division-kernel-finite-dimension-decay-embedding
Such a function determines the representation. The discrete series with lowest $K$-type of weight
$\Lambda$ contains the model of that $K$-type, and its matrix coefficients there are the values of
the function.

% before: asm-division-kernel-finite-dimension-decay
A matrix coefficient of the lowest $K$-type decays exponentially along the noncompact directions,
at a rate proportional to the smallest gap between entries of the Harish-Chandra parameter with
opposite signs. A large gap therefore beats any fixed polynomial weight.

% before: asm-division-kernel-finite-dimension-growth
The balls of the Hilbert--Schmidt norm in an indefinite unitary group have Haar measure bounded by a
power of their radius. So a function that decays fast enough in the radius of the polar
decomposition is integrable.

% before: asm-division-kernel-finite-dimension-weighted
For the signs and the parameters that occur in the setting, the gap condition makes the decay rate
larger than the polynomial weight plus the growth exponent of the balls. The weighted matrix
coefficient is then integrable.

% before: asm-division-kernel-finite-dimension
The space of automorphic forms has to be finite-dimensional for the Gram
matrices to be finite. The finite dimensionality of rows gives this once the
archimedean representations have formal degrees and integrable matrix
coefficients.

% before: asm-division-kernel-finite-fibres-induced
The automorphisms of a row lattice that fix the flag act on the orthogonal complement of the flag
modulo the flag, where the form is positive definite after a reflection in the negative vector. A
lattice with integral values has only finitely many vectors of bounded norm for such a form, and an
induced map is determined by the images of a basis taken in the lattice.

% before: asm-division-kernel-finite-fibres
The unfolding of the Gram matrices sums over terms that refine the terms of the
unfolding. Each term below has only finitely many terms above it, because these
inject into the automorphisms of a row lattice.

% before: asm-division-kernel-finite-pairs-invariants
The corank and the scalarity of a pair of terms are computed from the row lattice, the flag of the
first word and the two moved flags of the second word. They compare the orthogonal complements of
the lattices spanned by the flags, modulo the flags.

% before: asm-division-kernel-finite-pairs-move
Terms over the same term of the unfolding differ by unipotent automorphisms of the rows and by an
isometry fixing the flag of the second word. Each of these moves the spaces that define the corank
and the scalarity by a known isometry, provided the isometry fixing the flag acts trivially modulo
the flag.

% before: asm-division-kernel-finite-pairs-kernel-inputs-lattice
The rows of a word are the rational points of translates of the adelic lattice. Membership in a row
splits along the hyperbolic pairs, and a translate of the adelic lattice is determined by its rational
points.

% before: asm-division-kernel-finite-pairs-kernel-inputs-reduction
The reduction of a term is read from the Levi block of an adelic element at the level primes. On
elements that are parabolic modulo $g$ together with their inverses this reduction is multiplicative,
and an element with trivial reduction and the congruences of the flag lies in the level subgroup.

% before: asm-division-kernel-finite-pairs-kernel-inputs-terms
The rational points fixing the flag and a row of the second word act on the terms by right
multiplication. The reduction of a term can be computed at any representative of its double coset.

% before: asm-division-kernel-finite-pairs-kernel-inputs
The rational points that fix the flag of a word and preserve one of its rows act on the terms with that
row. Each map they induce modulo the flag comes from exactly one of them, and the reduction of its
Levi block at the level primes multiplies the reduction of a term and determines the point, by
neatness.

% before: asm-division-kernel-finite-pairs-kernel
That the isometries fixing the flag act trivially modulo the flag comes from the reductions at the
level primes. Two terms with the same rows have the same reduction, and the reduction of the Levi
block determines the rational point fixing the flag that lies behind the isometry, because the level
subgroup is neat.

% before: asm-division-kernel-finite-pairs
The corank and the scalarity of a pair of terms enter the bounds for the Gram
matrices. They are invariants of the terms below, since terms over the same
term differ by automorphisms that do not change them.

% before: asm-division-kernel-finite-volumes-cocompact-pattern
In a basis adapted to the flag of the isotropic vectors, the unipotent group of the Bessel data
consists of the isometries that differ from the identity by a matrix vanishing below a staircase. The
Cayley transform turns skew matrices of this shape into such isometries, and its coefficients are
integral up to powers of $2$.

% before: asm-division-kernel-finite-volumes-cocompact-cayley
The cocompactness is proved by multiplying an isometry near the identity by two Cayley transforms and
showing that the product is nearer to the identity. Products of staircase matrices go further down
the staircase, and an isometry agrees with the identity plus its skew part to first order, so the
error moves one step down.

% before: asm-division-kernel-finite-volumes-cocompact-integers
The rational points used in that proof come from integral matrices that must approximate a given
archimedean matrix with a bounded error. For a CM field this rests on the ring of integers being a
lattice in the archimedean completion.

% before: asm-division-kernel-finite-volumes-cocompact-approx
A family of skew staircase matrices at the infinite places is approximated by one rational skew
staircase matrix, divisible by a given integer, with an error bounded independently of the family.
Rounding the entries gives the integrality and the staircase shape, and taking the skew part then
restores skewness while keeping both.

% before: asm-division-kernel-finite-volumes-cocompact
The volume $\mathrm{vol}_N(u)$ is the measure of a fundamental domain for the rational points of the
unipotent group of the Bessel data, and its finiteness needs that group to be cocompact. Its rational
points form a lattice in the archimedean group, which is filtered by the flag of the isotropic vectors.

% before: asm-division-kernel-finite-volumes
The volumes of the setting appear as normalizing constants in the Bessel
periods. The normalized level subgroup makes them positive.

% before: asm-division-kernel-finite
Part (7) uses the isomorphism of the Bessel-period data of a member with
those of its automorphic setting. That isomorphism needs the spectral data of
the setting to be finite: the column maps bijective, the Poincaré series
square integrable, the space of automorphic forms finite-dimensional and the
terms of the unfolding finite in number over each term.

% before: asm-tied-from-green
With a collection of definite-side hypotheses at hand, every link of
Theorem~\ref{asm-tied-from-links} can be checked. Green's parametrization of
the cuspidal characters of $GL_d(\mathbf F_g)$ is needed only for the
irreducibility link. The theorem proves the property of
Definition~\ref{asm-tied-everywhere}.

% before: asm-grid-evaluation
The rows of the construction are tested by lifting a row form from $U(Y_0)$
through the theta grid of Chapter~\ref{chapter-unitary-theta} and pairing the
lift with a vector-valued form on the top space by a Bessel period. A grid,
its Weil data and its measures must be tied to the realization being tested,
and the test succeeds when some such pairing does not vanish.

% before: asm-stage-construction
The construction of definite-side data proceeds through stages, and most of
its properties hold only at large stages: small exponent, large scale, and
coefficient fields and twisting primes large enough for finitely many given
requirements. The row test and the availability of evaluation grids are
properties of this kind.

% before: asm-dependence-parameters
A collection of definite-side hypotheses needs bundles at the base points of
the construction, and these are built from parameters: a condition on stages
that decides which cores count as built. Every property the collection needs
is obtained in the form ``linked at the parameters'', which allows a slot of
cores that holds at large base data to be added.

% before: asm-open-parameter-statements
Most hypotheses of Theorem~\ref{is-definite-rows} are properties of the
selected bundles that the construction of bundles provides directly. Gathered
in one definition, they can be passed together from one set of parameters to
a stronger one.

% before: asm-dependence-linked
Linked properties are combined and transported in every step of the
construction. Strengthening the extension removes cores and links but keeps
every linked property, so the open parameter statements, proved for the
parameters of dependence, also hold once the grid property is added.

% before: asm-dependence-howe-unramified
The separation of functionals rests on a local statement at the places where the
Hermitian forms and the additive character are unramified. There the Weil
representation of a pair of unitary groups is generated by one vacuum vector, and
the vectors fixed by the integral points of one group come from that vector by
the action of the other group.

% before: asm-dependence-zero-sums
Besides the zero-density estimates and the zero-free regions of
Chapter~\ref{chapter-host-fields}, the analytic input uses two classical facts
about the zeros of Dirichlet $L$-functions. The explicit formula expresses a
twisted Chebyshev sum through the zeros up to a given height, and the zeros in a
window of bounded height are few.

% before: asm-dependence-genus-mass
The size condition and the volume items need the mass of a positive definite
Hermitian lattice of rank $2$ over the host field. By the definite case of the
volume and mass identity of Chapter~\ref{chapter-covolumes}, the standard lattice
of every positive definite space has the mass of the identity form of its rank.

% before: asm-dependence-projections-onto
The projection matrices of a row-side family send the coordinates of level forms
to the coordinates of a definite subspace. Their surjectivity, one of the items of
a collection of definite-side hypotheses, is linear algebra: the good part of the
forms consists of level forms, so the definite subspace is reached from the level
forms themselves.

% before: asm-dependence-base-vectors
The joins of a collection also ask that the base space of the base tuple map onto
the top forms. The ties of a row-side family express this map through the
columns at the base tuple, so the surjectivity of the columns, which primitive
rows give, carries over to the base space.

% before: asm-dependence-complement-norms
Two of the second complement inputs hold at every member, before any size
condition is used. The auxiliary prime is a norm congruent to $1$ modulo the
level primes, and the valuations of a level prime's factor $\overline{\pi_i}$
split its logarithm evenly between conjugate places.

% before: asm-dependence-findim
The near-orthogonality of the columns is proved inside a space of automorphic
forms on an indefinite unitary group, cut out by a Hecke operator and by a
discrete series at the real places. That this space is finite-dimensional is a
form of Harish-Chandra's finiteness theorem.

% before: asm-dependence-orbit-witness
Part (6) of Theorem~\ref{asm-dependence-inputs} proves the open parameter
statements at the parameters of dependence $c$, one statement at a time. Two of
them need no published result. A bundle linked to a base point is built from a
twisted input, and the input already records primitive rows, lift-distinct
types and the readings of the periods.

% before: asm-dependence-orbit-inputs-transfer
Each part of Lemma~\ref{asm-dependence-orbit-inputs} proves that a property of
the selected bundles is linked at the parameters. In each case the property is
first proved at the bundles linked to base points, or at the bundles built at a
stage with a witness, and then passed to the selections. The passage is the same
every time, and it also gives the eventual forms of the properties.

% before: asm-dependence-orbit-inputs-comparison
Statement (d) asks for a comparison part of the selected bundle. A bundle linked
to a base point is built from a witness, and the witness, with the host frames
hypothesis, normalizes the realization from which the comparison part is made.

% before: asm-dependence-orbit-inputs-join
Statement (c) asks that the indefinite realization be joined to an automorphic
setting. At a bundle with a witness, the comparison part normalizes the
realization of the twisted input, and the join is proved for that realization.

% before: asm-dependence-orbit-inputs-multiplicity
The multiplicity condition at the face of the realization is the first input
that uses the theta correspondence globally. It holds at every bundle linked to
a base point of a good core, once the diagonal rows are known from Siegel's
formula.

% before: asm-dependence-orbit-inputs-division
The division rows come from the multiplicity condition. At a bundle built at a
stage with a witness, the multiplicity condition makes a pairing on the definite
base vanish, and the rationality law of the datum turns this into the division
rows.

% before: asm-dependence-orbit-inputs-multiplicity-one
The rows with $\Lambda_3$, the archimedean scalars and the equality of cuspidal
multiplicities in Lemma~\ref{asm-dependence-bundles-inputs} all use one
comparison of cuspidal multiplicities between the base space and the spaces
$W_i$. Its hypotheses are conditions that the realization of a stage records.

% before: asm-dependence-orbit-inputs-lambda
The rows with $\Lambda_3$ follow from the equality of cuspidal multiplicities at
the realization of the stage. That equality gives the isotypic rank of the
datum, and the isotypic rank gives the condition $\Lambda_3$ at the readings of
the witness.

% before: asm-dependence-orbit-inputs-kappa
The archimedean scalars of the columns are compared with those of the division
data at the infinite place. The comparison needs the rows with $\Lambda_3$ and a
nonzero isotypic part, which the equality of cuspidal multiplicities provides.

% before: asm-dependence-orbit-inputs-good
Statements (e) and (f) rest on Green's parametrization of the characters of the
groups $GL_r(\mathbf F_g)$ at the level places. The goodness of the rows holds at
the stages of the law extension, and the labels of the rows are tied once the
rows are good.

% before: asm-dependence-orbit-inputs
The other open parameter statements rest on the published results on unitary
groups, and several of them on statements already linked. The comparison parts
come before the joins, which use them; the multiplicity condition comes before
the division rows, and the rows with $\Lambda_3$ before the archimedean scalars.

% before: asm-dependence-covolume-ties-points-host
The coefficient field bound at a point compares the field of the datum with the
host. Its first part needs the host to be small in terms of $x^{\varepsilon}$:
the degree, the level primes, the primes $p_{de}$ and $p'_{de}$ and the
discriminant. These bounds follow from the parameter choices that every base
point carries.

% before: asm-dependence-covolume-ties-points-core
The further conditions of a point include the core ties, which the runs of
Chapter~\ref{chapter-integral-structure} use. At a point they follow from the tie
of the core to the frame of its realization. The content is that the level ideal
of the core is the product of the places of the group, and that the places over
$p_{de}$ and $p'_{de}$ in the ramified sets lie over one prime of $K$ each.

% before: asm-dependence-covolume-ties-points-fields-weber
The field ties ask that automorphisms of $\mathbf C$ fixing the coefficient field fix a CM curve with its points of
a given order. This follows once the Weber values of that order lie in the ray class field of the field datum, which
is the theory of complex multiplication for the lattice $\mathcal O_K$.

% before: asm-dependence-covolume-ties-points-fields-size
The size bound for the coefficient field comes from a bound for the field $\Omega_\delta$ that contains all the
fields of the field datum. Its degree and root discriminant are bounded in terms of the ranks and the size bound of
the host.

% before: asm-dependence-covolume-ties-points-fields
Two of the ties at a weighted point concern the coefficient field alone: the
field ties and the condition on small places. Both follow from the existence of
ray class fields, because the coefficient field is the compositum of fields that
an exact field datum lists, among them a ray class field of $E$.

% before: asm-dependence-covolume-ties-points-weighted
The intrinsic conditions of the pinned point reduce to the ties of the side of
the stage at its weighted point. Most of these ties are recorded by the stage,
its rows and its norm choice; the field ties and the small places come from the
previous lemma, and the conjugates of the generators of $g_0$ and $g_1$ are
placed by the column choice.

% before: asm-dependence-covolume-ties-points
The proof of Lemma~\ref{asm-dependence-covolume-ties} combines four facts about
the points of comparison parts. Two come from the stage and its witness: the
normalized embedding with the class-sum pin, and the compatibility conditions.
The other two, the ray condition at large base data and the intrinsic conditions
with the coefficient field bound, use the existence of ray class fields.

% before: asm-dependence-covolume-ties
A comparison part of a selected bundle carries a point in the sense of
Definition~\ref{is-point}. The geometric weak compatibility in part (6) and the
covolume bound in part (7) both need this point to satisfy the class-sum pin and
the further conditions. The covolume bound also needs the coefficient field
bound of the member.

% before: asm-dependence-orbit-geometric
Statement (g) asks for one comparison part whose point has these properties.
Statement (d) supplies a comparison part, and
Lemma~\ref{asm-dependence-covolume-ties} gives the properties at every comparison
part.

% before: asm-dependence-covolume-points-rows
Most of the point conditions concern the realization of a comparison part and
its indefinite realization. They follow from the four results on automorphic
forms, with no further condition on the part.

% before: asm-dependence-covolume-points-export-span
The export condition asks for a lattice in the base change of the space $V_0$ of
the twisting data at a twisted place. That space is cut out of the space of forms
by a level character and a central idempotent, and so are the rows of the
indefinite realization. A semilinear map that is compatible with the group
actions carries the one to a spanning set of the base change of the other.

% before: asm-dependence-covolume-points-export-lattice
Once the span is known, the lattice comes from the lattice of the division data.
Its image under a semilinear map spans a lattice over an integrally closed ring
with the right fraction field, since the integers of the number field map into
that ring.

% before: asm-dependence-covolume-points-export
One point condition, the export condition, bounds the real expansion of the
indefinite realization. It follows from the divisor conditions and the class-sum
pin of the part.

% before: asm-dependence-covolume-points
The covolume bound of Theorem~\ref{cv-covolume-bound} applies at points that
satisfy the point conditions. At the selected parts these conditions come from
the published results on automorphic forms, together with the primitivity,
divisor and goodness statements at the parts.

% before: asm-dependence-covolume-bound
Theorem~\ref{cv-covolume-bound} is stated for families with the family
hypotheses, at points with the point conditions and the coefficient field bound.
For the selected parts of a class whose families carry their duty data, these
hypotheses hold eventually, so the bound reaches the rows of the selected
realizations.

% before: asm-dependence-bundles-member
Part (8) passes from the origin extension to its extension by a property $A$ of
stages. The member property survives this passage when $A$ holds eventually in
the context and at large base data, since $A$ can then be imposed in both passes
of the construction.

% before: asm-dependence-bundles-inputs-move
Parts (2) and (4) of Lemma~\ref{asm-dependence-bundles-inputs}, and part (9) of
Theorem~\ref{asm-dependence-inputs}, change the division lattices and heights of
a later stage. The rows, the norm choice and the further data of the stage pass
to the moved stage, so the conditions of the parameters survive the move.

% before: asm-dependence-bundles-inputs
At the origin extension, the member property and the existence of bundles at
large bases rest on the published results, through the equality of cuspidal
multiplicities at every stage. After a further extension, the existence of
bundles also needs that the positions of the parameters move with the division
data.

% before: asm-dependence-rows-models
The separation of functionals, part (1) of Theorem~\ref{asm-dependence-inputs},
is proved with automorphic forms whose components at the real places are
discrete series. They are recorded by multiplicity maps from a model of the
discrete series to cusp forms, and the iterated theta lifts are followed
through such maps along the first row of the grid.

% before: asm-dependence-rows-sizes
The sizes that size data impose on the rows contain the arithmetic and
archimedean conditions under which the separation is proved. The archimedean
condition is a gap between consecutive $b$-parts of the Harish-Chandra
parameters, large against the ranks.

% before: asm-dependence-rows-local
Two inputs at the real places start the separation. The row and core spaces
have discrete series models with the Cartan involutions of the grid, and one
theta lift along the row carries a multiplicity map of one row space to one of
the next through a map between the models.

% before: asm-dependence-rows-tower
Iterating the row step along the first row of the grid computes the partial
row lifts of a multiplicity map. Each of them is again a multiplicity map,
evaluated at a vector that depends only on the real vector and the starting
vector.

% before: asm-dependence-rows-global
Four further inputs are used: the convergence of the lifts and periods, a
nonzero core period at large scales, the spanning of the core multiplicities by
the ends of column chains, and the uniqueness of Bessel functionals at the real
places. The last one makes the Bessel period of a multiplicity map at the top
row a fixed function of the functional and the vector, up to a scalar.

% before: asm-dependence-rows-good
The theta multiplicities of a realization at the trivial character have the
level of the trivial subgroup of its group. That level contains the principal
congruence subgroup modulo $g_0g_1$, so these multiplicities are row
multiplicities.

% before: asm-dependence-rows-factor
Along the row tower, the uniqueness of Bessel functionals at the top row
separates the dependence of the Bessel period on the real vector and the
functional from its dependence on the multiplicity map, the column and the
coset datum. This lets one real vector and one functional serve for all
multiplicity maps at once.

% before: asm-dependence-rows-weak
For one multiplicity map at a time, some Bessel period of an iterated lift is
nonzero. A core multiplicity has a nonzero core period against the form, the
core multiplicities are spanned by the ends of column chains, and the ascent of
nonvanishing carries the period up a chain to the top Bessel period.

% before: asm-dependence-functional-grid
Part (5) of Theorem~\ref{asm-dependence-inputs} applies the Fourier fit of part
(3) at the points of the selected comparison parts. That needs an evaluation
grid at each such point, which the grid property of the stage supplies once the
point satisfies the compatibility conditions.

% before: asm-dependence-functional-geometric
The comparison parts of the selected bundles carry, besides the properties of
Lemma~\ref{asm-dependence-covolume-ties}, a fitted evaluation grid. Both are
linked at the parameters of the extension by the grid property.

% before: asm-dependence-functional-bound
With the point conditions and the evaluation grids in hand at large scales, the
conclusion of part (3) gives the functional bound for the rows of the selected
parts. The argument is the one that gives the covolume bound in
Lemma~\ref{asm-dependence-covolume-bound}.

% before: asm-dependence-test-assembly-statements
The row test is reached through finitely many test functionals, Bessel periods
of iterated theta lifts of the rows. Their properties are stated for every
realization in an automorphic setting at large scales, with constants that do
not depend on the datum, so that they can be used uniformly along the stages.

% before: asm-dependence-test-assembly-first-layer-rows-pieces
Statements (c) to (f) on test choices rest on the theta correspondence for the twisted Weil
representations at the squares of the grid, and each uses only part of its inputs. Statement (f)
needs only statement (d).

% before: asm-dependence-test-assembly-first-layer-rows
Five of the seven statements on test choices concern the theta lifts of the rows. They follow from
three facts about the twisted Weil representations at the squares of the grid.

% before: asm-dependence-test-assembly-first-layer-convergence-lift-grid
The iterated theta lift of a test functional is computed in a grid of models of discrete series.
The grid has to carry the spaces and measures of the evaluation grid, and its theta kernels have to
be those of the Weil datum of the test choices.

% before: asm-dependence-test-assembly-first-layer-convergence-lift-cuspidal
Along a grid of models each theta lift is an integral over a fundamental domain. For cuspidal
multiplicity maps these integrals converge and give cuspidal multiplicity maps again, which lets the
lifts be iterated.

% before: asm-dependence-test-assembly-first-layer-convergence-lift
A test functional converges when the iterated theta lift converges and its Bessel period converges.
The first is followed down the chain of spaces, along which theta lifts of cuspidal vectors stay
cuspidal. The lift obtained is then an automorphic function, for which the Bessel period converges.

% before: asm-dependence-test-assembly-first-layer-convergence
The real condition compares the test functionals for different real vectors and functionals. This
needs their convergence for every such choice, which the cuspidality of the theta lifts of pairs
gives.

% before: asm-dependence-test-assembly-first-layer-columns
The columns of the test choices are delta functions at the class representatives. For a fixed coset
datum they run over a spanning family of the coefficients at every class.

% before: asm-dependence-test-assembly-first-layer-injective
Injectivity of the test functionals is where the separation of functionals enters. A vector killed
by all test functionals is killed by the Bessel periods of all columns, against which separation
produces a nonzero one.

% before: asm-dependence-test-assembly-first-layer-ratio-inputs-nonvanishing
The ratios of the constants of a square involve values of Hecke $L$-functions at integers. These do
not vanish to the right of the critical strip, and the functional equation carries the nonvanishing
to the left, as long as the gamma factors stay finite.

% before: asm-dependence-test-assembly-first-layer-ratio-inputs-characters
The nonvanishing lemma needs a CM type and the archimedean components of the characters of the
squares. The embeddings of the definite-side datum send $\sqrt{-d}$ to $i\sqrt d$, which gives the CM
type. The level system and the grid fix the archimedean components of the characters, and the grid's
gap makes the exponents larger than half the rank of the square.

% before: asm-dependence-test-assembly-first-layer-ratio-inputs-definite
The theta integrals of a bottom triple are taken over the groups of totally definite spaces. Their
quotients by the rational points are compact, so the integrals converge and depend continuously on
the point of the target group.

% before: asm-dependence-test-assembly-first-layer-ratio-inputs-vanish
The ratio of the factorization constants is read off from values of the theta integrals at points
with archimedean components $1$, where the bottom sums are known to be rational. A continuous
automorphic function is determined by these values, by weak approximation at the archimedean places.

% before: asm-dependence-test-assembly-first-layer-ratio-inputs-standard-waves
A vacuum vector fixed by a lattice and by its dual lattice determines a Schr\"odinger model on the
functions on one half of the space. The functions to which the model applies are spanned by translates of
characters cut off to the lattice, and on these the model is forced by the Heisenberg relations. So the
model exists, and its values on these functions do not depend on the choices made in building it.

% before: asm-dependence-test-assembly-first-layer-ratio-inputs-standard
The functions of a bottom triple are obtained from Fock vacuum vectors by putting a standard function in
the components at finitely many places. A standard function is invariant under translation by a
multiple of the lattice and supported in a dilate of it, so it is a Schwartz-Bruhat function on the part
of the space at these places.

% before: asm-dependence-test-assembly-first-layer-ratio-inputs-weil-root
Characters of the adelic unitary group are built from Hecke characters through determinant roots. Their
continuity needs determinant roots near $1$ for determinants near $1$.

% before: asm-dependence-test-assembly-first-layer-ratio-inputs-weil-twist
A Hecke character trivial on the ideles of $F$ gives a character of each adelic unitary group through
determinant roots. Multiplying a Weil datum by two such characters changes its characters and keeps the
other conditions on a Weil datum.

% before: asm-dependence-test-assembly-first-layer-ratio-inputs-weil-eta
The Weil data are built from one Hecke character of $E$ that extends the quadratic character of $E/F$.
Chapter~\ref{chapter-integral-structure} constructs it for the phase of the conjugate frame: it is unramified at every finite place and is $z/|z|$ at each
embedding of the datum.

% before: asm-dependence-test-assembly-first-layer-ratio-inputs-weil-powers
A Weil datum whose characters are powers of this extension is the product of a lattice model at the finite
places and a Fock model at the archimedean places. The two models commute, and the invariance of the theta
distribution under the rational points gives the remaining condition on a Weil datum.

% before: asm-dependence-test-assembly-first-layer-ratio-inputs-weil
The bottom triples need Weil data with prescribed characters. They are obtained from one Weil datum whose
characters are powers of a fixed unramified extension of the quadratic character, by the twist above.

% before: asm-dependence-test-assembly-first-layer-ratio-inputs-factorization
The theta integrals of a bottom triple at all tuples factor through one space, by the matching of local
maps at the square triples. On it the integrals of the larger space are multiples of those of the
smaller one.

% before: asm-dependence-test-assembly-first-layer-ratio-inputs-towers-stages
The constant of a square is compared with that of a bottom triple along a Witt tower. Its stages are
triples whose first two spaces lose one hyperbolic plane at each step, from the spaces of the square down to
the definite spaces of a bottom triple, with the characters and finite data of the square throughout.

% before: asm-dependence-test-assembly-first-layer-ratio-inputs-towers-inputs
Five facts about the stages carry the comparison. The top stage sees the constant of the square, every stage
has a factorization, each step down produces one value of the $L$-function of the square, and the last
stage sees the constants of a bottom triple. The factors that appear at each step have ratios in the
coefficient field.

% before: asm-dependence-test-assembly-first-layer-ratio-inputs-towers-descent
Going down the tower step by step, by induction on the two indices of the stage, collects one $L$-value
per step. The constants of a stage are a fixed multiple of the numbers of its factorization, so the
constants can be traded for the factorization at each step.

% before: asm-dependence-test-assembly-first-layer-ratio-inputs
The ratio of the constants of a square at two tuples is computed along a tower at the square. The
tower relates it to a ratio of constants of a bottom triple and to values of Hecke $L$-functions.
Each of these lies in the coefficient field once the field conditions hold and the $L$-values do not
vanish.

% before: asm-dependence-test-assembly-first-layer-ratio
The last statement compares the constants of the squares at different tuples. Their ratios come
from the triples and towers at the squares and from the nonvanishing of $L$-values there.

% before: asm-dependence-test-assembly-first-layer-grid-pieces-places
A theta grid is built at a base place above the auxiliary prime and at chosen places of $E$ above the
level primes. The choices must agree with the places of the level data and with the ramification of
the characters of the level system, and the places must not be conjugate to themselves, so that the
local idempotents exist.

% before: asm-dependence-test-assembly-first-layer-grid-pieces-spaces
The spaces of a theta grid have prescribed ranks and signatures, and some of them must contain others with
complements of a given shape. Diagonal forms give the first, and orthogonal sums with antidiagonal forms the
second; the fractional ideals on the added basis vectors are paired so that the sum of the lattices stays
self-dual.

% before: asm-dependence-test-assembly-first-layer-grid-pieces-characters-wild
The characters of a theta grid must have conductor $P^2$ at the places $P$ over the level primes. A
character of the one-units of order $p$, read through the residue of $(x-1)/\varpi$, gives this conductor at a
place over $p$.

% before: asm-dependence-test-assembly-first-layer-grid-pieces-characters-tame
At the base place the characters must be tamely ramified. A character of the residue field that is not trivial
can be chosen with prescribed values on the roots of unity, once the residue field has more than one element
per root of unity; these values make up for the other components on the roots of unity.

% before: asm-dependence-test-assembly-first-layer-grid-pieces-characters-extension
Hecke characters are built from characters of open subgroups of the ideles. A unitary character of an open
subgroup extends to the whole group with prescribed values on the image of a homomorphism, if the two
prescriptions agree where both apply, because the unit circle is divisible.

% before: asm-dependence-test-assembly-first-layer-grid-pieces-characters-anti
A character of the ideles that is trivial on the roots of unity becomes a Hecke character of the CM field after
dividing by its conjugate. The quotient is trivial on the ideles of $F$, which the characters of a theta grid
must be.

% before: asm-dependence-test-assembly-first-layer-grid-pieces-characters-hecke
Local characters at a set of places that contains no pair of conjugate places, with an archimedean type, glue
to a Hecke character of the CM field that is trivial on the ideles of $F$. The only condition is that their
product is trivial on the roots of unity of the field.

% before: asm-dependence-test-assembly-first-layer-grid-pieces-characters
The characters of a theta grid have a prescribed archimedean type, are trivial on the ideles of $F$, and have
prescribed conductors at the distinguished places. They are Hecke characters assembled from local characters
at these places.

% before: asm-dependence-test-assembly-first-layer-grid-pieces-weil-vectors
A grid Weil datum carries, on the pairs of its first column, a Fock vacuum vector and one coset vector for each class
modulo the modulus. The coset vectors come from the place model of the vacuum at the distinguished places, applied to
the indicator function of a coset.

% before: asm-dependence-test-assembly-first-layer-grid-pieces-weil-assemble
Once the Cartan involutions and the Weil data of every grid space and grid pair are chosen, the vacuum and coset
vectors can be built for them. So these choices can be made freely, and in particular taken from a realization on
the first column.

% before: asm-dependence-test-assembly-first-layer-grid-pieces-weil-exists
The Weil data of the grid pairs come from the Weil representation of unitary dual pairs. Its hypothesis is that
the two characters of a pair restrict to the right powers of the quadratic character of $E/F$, which the conditions
of the grid and of the level system give.

% before: asm-dependence-test-assembly-first-layer-grid-pieces-choice-standard
A grid choice is built on rational data: a finite set of primes off which nothing of the grid or of the level system
ramifies, and orthonormal bases of the definite grid spaces with finite adelic points that carry their lattices to the
lattices of the grid. The adelic points are supported at split places off these primes.

% before: asm-dependence-test-assembly-first-layer-grid-pieces-choice-normal
The vacuum vectors of the definite pairs are compared with an odd theta function of a CM lattice of the host. The
normalization fixes, for each definite pair, a constant by which the values at $0$ of the Fock vectors are the
derivatives at $0$ of the odd function, scaled by periods.

% before: asm-dependence-test-assembly-first-layer-grid-pieces-choice-grid
A grid choice puts standard data and a vacuum normalization together with adelic bases of the grid lattices. At the
local places these bases are adapted to the decompositions of the last row and of the last top space, with prescribed
powers of a uniformizer, and a finite set of places records the primes of the ideal of the host.

% before: asm-dependence-test-assembly-first-layer-grid-pieces-choice-primes
The first condition on standard data asks for primes off which the places of $E$ avoid finitely many given places
and integers. The prime factors of a product of norms suffice.

% before: asm-dependence-test-assembly-first-layer-grid-pieces-choice-unramified
The second condition asks that the characters of the tuples be unramified at the split places off a finite set of
primes. A continuous character of the ideles is trivial on a congruence subgroup of the unit ideles, which gives such a
set.

% before: asm-dependence-test-assembly-first-layer-grid-pieces-choice-orthonormal-pieces
The orthonormal bases and adelic points of the definite spaces come from three facts on Hermitian spaces. They are the
existence of an orthonormal basis of a totally definite space, the local conjugacy of self-dual lattices, and a
decomposition of the finite adelic unitary group.

% before: asm-dependence-test-assembly-first-layer-grid-pieces-choice-orthonormal
Conjugating by the rational part of the decomposition moves the orthonormal basis and keeps the support of the adelic
point off the given primes. This gives the data of condition (c) for one definite space.

% before: asm-dependence-test-assembly-first-layer-grid-pieces-choice-top
The last top space $Q_m$ is the sum of $V_{m,m}$ and a line spanned by a unit vector. Its basis and adelic point are
those of $V_{m,m}$ extended by that vector, which gives condition (d).

% before: asm-dependence-test-assembly-first-layer-grid-pieces-choice-standard-exists
Together these give standard data on every finite set of tuples. The definite spaces are totally definite with
self-dual lattices, as the datum, the level system and the grid require.

% before: asm-dependence-test-assembly-first-layer-grid-pieces-choice-odd
The vacuum vectors are normalized against odd vacuum functions of the CM lattice of the host. Such a function exists
for every size with even product, as an averaged tensor power of the theta function of the lattice.

% before: asm-dependence-test-assembly-first-layer-grid-pieces-choice-jets
For a definite pair, the values at $0$ of the Fock vectors of a vacuum vector and the derivatives at $0$ of an odd
vacuum function are proportional, with the periods of the lattice as scaling factors. The factor of proportionality is
the constant of the normalization.

% before: asm-dependence-test-assembly-first-layer-grid-pieces-choice-normal-exists
At each definite pair there are an odd vacuum function and a Fock vacuum vector, and the comparison of their jets
gives the constant. So all standard data have a vacuum normalization.

% before: asm-dependence-test-assembly-first-layer-grid-pieces-choice-adelic
The adelic bases of a grid choice are built from adelic bases of lattices and from ideles with given valuations. A sum
of a lattice and multiples of complement vectors then has an adelic basis that is the sum of the given one and these
multiples.

% before: asm-dependence-test-assembly-first-layer-grid-pieces-choice-rows
The lattices of the last row of the grid are sums with staircase ideals built from the ideal of the host and the
modulus. At a local place prime to the ideal of the host their valuations are the exponents asked of a grid choice.

% before: asm-dependence-test-assembly-first-layer-grid-pieces-choice-frames
With these valuations the adelic bases of the last row and of the last top space can be chosen adapted at the local
places. Standard data and a vacuum normalization then extend to a grid choice.

% before: asm-dependence-test-assembly-first-layer-grid-pieces-choice-different
The ideal $\mathfrak b$ of the host and its conjugate multiply to the different of $E/K$. So its prime divisors lie
over the primes that ramify in $E/K$, which are those dividing $Q_Sq$.

% before: asm-dependence-test-assembly-first-layer-grid-pieces-choice-host-ideal
The adapted bases need local places at which the ideal of the host has exponent zero. A place prime to
$Q_Sq\,\mathrm N\mathfrak c$ is one, by the description of the prime divisors of $\mathfrak b$.

% before: asm-dependence-test-assembly-first-layer-grid-pieces-choice
A grid choice adds to a grid Weil datum adelic bases of the grid lattices adapted at the local places, standard data
and a vacuum normalization. It exists when the local places avoid the primes of the host's ideal $\mathfrak A$.

% before: asm-dependence-test-assembly-first-layer-grid-pieces
An evaluation grid with its Weil datum and grid choice is built in four steps: its places, the grid
on them, the Weil datum from the Weil representations of the realization, and the grid choice. Each
step asks only for arithmetic conditions on the level system and for conditions of the core, which
the automorphic setting supplies.

% before: asm-dependence-test-assembly-first-layer-grid
The number field statement starts from one evaluation grid with its Weil datum and grid choice for a
realization. Such a grid exists once the level primes split and the twisted places avoid $q_0$.

% before: asm-dependence-test-assembly-first-layer-field
Realizations with the same key share their grids, Weil data and grid choices by transport. Since
there are finitely many keys, finitely many number fields contain all the numbers that occur.

% before: asm-dependence-test-assembly-first-layer
The properties of the test functionals come from the theta correspondence for
the twisted Weil representations. Each of them uses only some of the inputs.

% before: asm-dependence-test-assembly-true
The row test at an arbitrary later stage is reduced to the test at a stage with
simpler rows. The reduction keeps the group of the stage, and the group data
fixed by the framing give a common target for both stages.

% before: asm-dependence-test-assembly-realization-equivariant-levels
The finite group of a realization acts on functions on the adelic unitary groups through lifts of
its elements to compact subgroups. These subgroups are cut out by the reductions of a lattice at the
places of the group and at the ramified places of a tuple. A function of the right type transforms
by a character under the kernel of the reduction, and an element acts by right translation by a
lift, normalized by the character.

% before: asm-dependence-test-assembly-realization-equivariant-lifts
Lifts exist once no place of the group is conjugate to a place of the group or to a ramified place
of the tuple. A lift is built place by place from a matrix over the integers of the host field,
through the idempotent of the place. For a realization the places of the group and the ramified
places lie over different primes, which gives the hypotheses.

% before: asm-dependence-test-assembly-realization-equivariant-shift
The action of an element does not depend on the lift that computes it. Shifting a multiplicity map
by a lift gives a multiplicity map of the same level, since conjugation by a lift preserves the level
and its character.

% before: asm-dependence-test-assembly-realization-equivariant-integral
Shifted multiplicity maps are compared through theta integrals over fundamental domains. Right
translation by an element of a compact subgroup preserves Haar measure, so these integrals do not
change under it.

% before: asm-dependence-test-assembly-realization-equivariant-scalars
The theta kernels of the chain are eigenvectors of the Weil data at pairs of lifts of one element of
the group. The eigenvalue is the quotient of the level characters at the two lifts, and with it the
theta lift of a shifted multiplicity map is the shift of its theta lift.

% before: asm-dependence-test-assembly-realization-equivariant-projection
The good part of the multiplicity maps is the image of the idempotent of the family, acting through
the level action. The coefficients of the idempotent are class functions, so this operator commutes
with the level action of every element.

% before: asm-dependence-test-assembly-realization-equivariant-schur-slot-closable-schur
The invariant pairings below are compared through an operator that commutes with a family of unitary
operators. That family contains the adjoint of each member, since the adjoint of a unitary operator is
its inverse. For such a family Schur's lemma follows from the functional calculus of self-adjoint
operators.

% before: asm-dependence-test-assembly-realization-equivariant-schur-slot-closable
The maps from smooth vectors to square-integrable functions given by multiplicity maps are not
bounded a priori. Schur's lemma still applies to invariant pairings of such maps once they are
closable, through the closure of their graphs.

% before: asm-dependence-test-assembly-realization-equivariant-schur-slot-smoothing
Closability is proved by smoothing vectors with test functions on the real unitary group. Smoothed
vectors are smooth, and a functional continuous for a Sobolev norm becomes continuous for the norm of
the Hilbert space after smoothing.

% before: asm-dependence-test-assembly-realization-equivariant-schur-slot
A multiplicity map vanishes once it vanishes at one nonzero vector, if its Petersson pairings are
multiples of the product of inner products. In each slot this is Schur's lemma for the discrete
series representation, in a form for bounded equivariant functions on a space of finite measure.

% before: asm-dependence-test-assembly-realization-equivariant-schur
With rapid decay, cusp forms are bounded, and the Petersson pairing of two multiplicity maps is
defined. Schur's lemma in each slot makes it a multiple of the product of inner products. A
multiplicity map that vanishes at one nonzero vector then has Petersson norm zero, so it vanishes.

% before: asm-dependence-test-assembly-realization-equivariant
The test functionals are evaluated at theta lifts of row forms, while statement (e) describes
them only on the good part of the row forms. The theta lifts commute with the action of the finite
group, and its idempotent acts trivially on the top space. So every theta lift of a row form is
the theta lift of a form in the good part.

% before: asm-dependence-test-assembly-realization
The first step of the assembly gives the row test from the test functionals,
whose values are rational multiples of rational sums. It needs test choices
whose numbers lie in the coefficient field.

% before: asm-dependence-test-assembly-existence
Test choices with numbers in the coefficient field are found once the
coefficient field contains finitely many algebraic numbers. That is a condition
at large base data.

% before: asm-dependence-test-assembly-sizes-explicit
The explicit values fix the size data only up to the choices that the datum itself determines.
When the indices of the rows vanish at the untwisted places, these choices can be made, and size
data with the explicit values exist.

% before: asm-dependence-test-assembly-sizes
At a stage whose rows have vanishing indices, the size data can be chosen at
the scale of the context. With the explicit division data, the rows then have
the sizes of the size data.

% before: asm-dependence-test-assembly-host-gaps
The Harish-Chandra parameters of the construction are arithmetic progressions with two common
steps. The gap conditions of an automorphic setting then reduce to lower bounds for these steps
and for the first $b$-part.

% before: asm-dependence-test-assembly-host-sizes
The steps of the Harish-Chandra parameters are powers of the scale, and the first $b$-part is at
least $w''(n)$. Both beat the polynomial sizes of the ranks once the scale is large.

% before: asm-dependence-test-assembly-host
The gap conditions of an automorphic setting concern only the Harish-Chandra
parameters of the datum, and the frame conditions only its spaces. Both hold at
the stages over a context once the scale is large.

% before: asm-dependence-test-assembly-setting
A realization built at a stage carries most conditions of an automorphic
setting in its records. The remaining ones are the sizes of the rows, the tie of
the series idempotents, which true group data give, and the core conditions,
which follow from existence at every place.

% before: asm-dependence-test-assembly-bare-norm-labels-series-induction
Labels of characters of finite general linear groups are read off from parabolic induction. Induction
from a block upper triangular subgroup and the average over its unipotent radical are adjoint for the
pairing of class functions. A cuspidal character has zero average over every such radical, and the
Mackey formula describes the average of an induced function.

% before: asm-dependence-test-assembly-bare-norm-labels-series-mackey
For a list of cuspidal characters the Mackey formula simplifies. Each cuspidal factor goes whole into
one of the two blocks, so the average over a radical of the induced character is a sum over the ways of
splitting the list between the blocks.

% before: asm-dependence-test-assembly-bare-norm-labels-series
The label of an irreducible character of a finite general linear group comes from the cuspidal
characters from which it is induced. Harish-Chandra's theory of parabolic induction shows that these
cuspidal characters exist and are determined up to order.

% before: asm-dependence-test-assembly-bare-norm-labels-uncertainty
A combination of few characters of a finite cyclic group cannot vanish at most of its points. This
is the tool for recovering a cuspidal character's polynomial from its values.

% before: asm-dependence-test-assembly-bare-norm-labels-count
Green's character formula is evaluated at elements whose characteristic polynomial is irreducible.
Their eigenvalues are the elements of exact degree $d$ over the prime field, and the elements of
smaller degree are few when $g\ge5$.

% before: asm-dependence-test-assembly-bare-norm-labels-cuspidal
For the label to be well defined, a cuspidal character must determine the polynomial whose character
formula it satisfies. Its values at elements with irreducible characteristic polynomial give sums of
characters of a cyclic group, and the two previous lemmas separate them.

% before: asm-dependence-test-assembly-bare-norm-labels
The types of the simpler stage use the series idempotents of the finite general linear groups at
the two level primes. They are indexed by labels: polynomials that record the cuspidal characters
from which an irreducible character is induced, through Green's parametrization.

% before: asm-dependence-test-assembly-bare-norm-series
Summing the central idempotents of the irreducible characters with a given label gives the series
idempotents. Since every irreducible character has one label, they are orthogonal and add up to $1$,
and their coefficients are cyclotomic.

% before: asm-dependence-test-assembly-bare-norm
Over rows without local defects and with vanishing operators, norm data can be built with no local
or twisted places. The group is then the product of the groups $GL_r$ at the places above $g_0$
and $g_1$, and its types come from the label data of these two primes.

% before: asm-dependence-test-assembly-bare
Green's parametrization of cuspidal representations of general linear groups
over finite fields produces group data that are true for a framing. Over rows
with vanishing indices this gives a later stage of the same column choice.

% before: asm-dependence-test-assembly-copy
A realization can be moved between two later stages with the same group, since
its data refer to the group only through the isomorphism of groups. The row
test at the copy gives the row test at the original.

% before: asm-dependence-test-assembly-second-layer
These lemmas assemble into the row test. A realization at a later
stage is copied to a stage with vanishing indices, where the sizes of the rows
hold, and there the test choices give the row test.

% before: asm-dependence-test-assembly
The rationality of the rows, part (2) of Theorem~\ref{asm-dependence-inputs},
rests on a row test for the complex data of the realizations built at the
stages. The test vectors come from theta lifts whose values lie in a number
field.

% before: asm-dependence-rows-parameters
The models of discrete series used along the first row of the grid have
parameters built from the Harish-Chandra parameter of the datum. They are
signed sets assembled from its $v$-, $b$- and $e$-parts and from blocks of
consecutive numbers centred between points $a_l$ spaced by the gap
$\gamma(c)$.

% before: asm-dependence-rows-counts
A model of discrete series with a given parameter exists when the parameter has
as many elements as the rank of the space, as many signs $-$ as its signature
asks for, and integral enumerations. These counts are read off from the
definition of the core and row parameters.

% before: asm-dependence-rows-model-exists
The grid Weil datum fixes a Cartan involution at each real place, and a row or
core model must have this involution. An orthonormal basis that diagonalizes
it, reordered to match the signs of the parameter, turns a discrete series of
the standard unitary group into such a model.

% before: asm-dependence-rows-cusp
The partial row lifts are values of elements of multiplicity spaces. Such
values are cusp forms, so they decay rapidly, and the convergence of the
periods along the row uses this decay.

% before: asm-dependence-rows-bessel-conv
The last lift of a row is paired with a column through a Bessel period. Since
the top space is definite, its automorphic quotient is compact, and the
period of any automorphic function converges.

% before: asm-dependence-grid-exists
The evaluation of the grid in part (4) of Theorem~\ref{asm-dependence-inputs}
asks for grids with full evaluation choices at all large stages. At a fixed
scale, the choices depend on the base data, the framing and a finite set of
characters containing those of the level system only through finitely many
algebraic numbers that the coefficient field must contain.

% before: asm-dependence-functional-twist-bounds
A grid fitted to a point asks two things of its twist choices at each twisted
place. The set of places of the twist must be split and stable under
conjugation, and the theta values twisted by every power of the twisting
character prime to $\ell_i$ must be nonzero with bounded valuation.

% before: asm-dependence-functional-fitted
The Fourier fit of part (3) of Theorem~\ref{asm-dependence-inputs} is stated for
points that carry an evaluation grid adapted to the point. The adaptation has
two halves: conditions on the grid and the twist choices, which make the
scalars of the grid rational over the coefficient field, and the choices that
complete the grid to a full evaluation choice.

% before: asm-dependence-functional-twists
The twist choices of a linked bundle come from a good core, whose conditions on
twists are stated place by place. The fit asks for them in the form of
conditions on pairs of conjugate places and on powers of the twisting
characters.

% before: asm-dependence-functional-pinned
The list extension pins an evaluation grid at the level characters of the theta
datum of every stage. A linked bundle carries that grid to the realization of
its comparison parts, with its values in the coefficient field and its scalars
tied to the twisted places.

% before: asm-dependence-functional-witness-sizes
Several statements below need the rows of the datum to be small against the
scale and far apart in their Harish-Chandra parameters. These requirements are
collected from the parameter choices, the size bounds and a few further
inequalities, so that one hypothesis supplies them all.

% before: asm-dependence-functional-witness-setting
The theta lifts along the evaluation grids are controlled only for
realizations whose rows have these sizes and whose data satisfy further
conditions on gaps, tuples, groups and Weil data. These conditions together
form an automorphic setting.

% before: asm-dependence-functional-witness-numbers-explicit
The construction fixes the ranks, the host, the Harish-Chandra parameters and the size data by
formulas in the scale and a few constants. The size estimates are proved for every datum with
these values, uniformly, as the constants and then the scale grow.

% before: asm-dependence-functional-witness-numbers-parameters-computations
For the explicit values each condition of the parameter choices compares powers, logarithms or
exponentials of the scale, once the constants are chosen in order. The conditions fall into
twenty-one groups, each of which is a computation at large scales.

% before: asm-dependence-functional-witness-numbers-parameters
The parameter choices are a long list of inequalities between the ranks, the host and the
quantities of the size data. For the explicit values each is a comparison of powers of the scale.

% before: asm-dependence-functional-witness-numbers-blocks-computations
The $T$-parameters take the block length $n_W$ and the number of blocks $T_O$ from the explicit
values. What remains is a list of inequalities between them, the ranks and the scale, each a
computation at large scales.

% before: asm-dependence-functional-witness-numbers-blocks
The $T$-parameters split $r+t$ into blocks of a fixed length. For the explicit values the block
length and the number of blocks can be chosen with the required bounds.

% before: asm-dependence-functional-witness-numbers-bounds-computations
The size bounds are windows and inequalities for the quantities of the size data. For the explicit
values they come in eleven groups, each a computation in the scale.

% before: asm-dependence-functional-witness-numbers-bounds
The size bounds are further windows and inequalities for the quantities of the size data. They
too follow from the explicit values once the scale is large.

% before: asm-dependence-functional-witness-numbers-band
The level primes $g_0$ and $g_1$ lie near $(kb)^8$, and the primes $p_{de}$, $p'_{de}$ are bounded
by a power of the discriminant term. Both are small against the scale bounds used later.

% before: asm-dependence-functional-witness-numbers-values
The four estimates combine because each holds for explicit values. Together with the windows of
the ranks they are the inputs of the computation of the sizes of the rows.

% before: asm-dependence-functional-witness-numbers
At a base point the size data take explicit values, chosen in terms of the
scale. At large scales these values satisfy the inequalities of the parameter
choices and of the size bounds.

% before: asm-dependence-functional-witness-computation
The remaining sizes of the rows are inequalities between powers of the scale.
They follow from the size bounds and the parameter choices once the scale is
large.

% before: asm-dependence-functional-witness-twisted
The twisted places of the stage are chosen as the places where an index of the
rows or a local term is nonzero. The witness carries this choice to the
indefinite realization of the bundle.

% before: asm-dependence-functional-witness-core
The weak characterization of a realization asks each theta map to have its
characterizing property only where the map exists. Existence at every place
turns it into the core conditions.

% before: asm-dependence-functional-witness
The grid property of a stage applies to realizations in an automorphic setting
with size data that satisfy the recursion conditions. For the realization of a
linked bundle these hypotheses come from the witness of its construction and
from the sizes of the rows, which hold eventually.

% before: asm-dependence-functional-full
With these hypotheses, the grid property gives, for every evaluation grid of
the realization of a comparison part, another one with the same base place and
characters that carries a full evaluation choice. Since the base place and the
characters determine the scalars, the pinned grid and this one can be combined.

% before: asm-dependence-grid-agree
The grid property is proved stage by stage, but the field conditions must hold
with one finite set of algebraic numbers for all later stages over the same base
data. Two realizations whose level systems agree carry agreeing grids, Weil data
and grid choices, with the same numbers.

% before: asm-dependence-grid-keys
Agreement of level systems is decided by the key, and the grids to be carried
over are classified by their grid keys. Both range over finite sets once the
base data, the twisting characters and the first constant are fixed.

% before: asm-dependence-grid-transport
At a representative stage for a key and a grid key, the numbers of a full
evaluation choice lie in a number field. Carrying the grid, the Weil datum and
the grid choice to another stage with the same keys keeps the numbers, so the
same number field serves there.

% before: asm-dependence-grid-pieces
The representative stages need Weil data, evaluation choices and number fields
for their numbers. These are the analytic and arithmetic inputs of the grid
property, and they hold at large scales.

% before: asm-dependence-rows-bessel-forms
The factorization of the Bessel periods in part (1) of
Theorem~\ref{asm-dependence-inputs} is decided at the real places, where the
Bessel period of a multiplicity map restricts to an invariant bilinear form.
These forms are the archimedean Bessel functionals of the representations
involved.

% before: asm-dependence-rows-bessel-unique
For discrete series the Bessel functionals are unique up to a scalar. This is
what makes the Bessel period of a row lift factor as a function of the real
data times a constant.

% before: asm-dependence-rows-bessel-line
Over several real places, a form that is a Bessel functional in each place
separately is determined up to a scalar by the product of the local
functionals. The uniqueness at each place gives this one coordinate at a time.

% before: asm-dependence-rows-gap
The uniqueness of Bessel functionals is applied to the models of the top row
space. Their parameters are integral, and the sizes of the rows separate the
entries of opposite signs by a gap of the required size.

% before: asm-dependence-rows-coefficients
The coefficients of the grid enter the Bessel period through a linear
functional. Writing their dual as a tensor product over the real places turns
the period into a function of one vector at each place.

% before: asm-dependence-rows-bessel-period
The separate Bessel property of the period rests on properties of the integral:
its additivity and a bound by the values on a compact set, its invariance under
simultaneous right translation, and its equivariance under the unipotent group.
The values of a multiplicity map on a compact set are bounded by a Sobolev norm
at each real place.

% before: asm-dependence-rows-bessel-separate
These properties show that the Bessel period of the value of a multiplicity map
is bilinear and is a Bessel functional at each real place separately. With the
uniqueness, it lies on a single line as the multiplicity map and the column
vary.

% before: asm-dependence-rows-step-pairing
A theta lift between models of discrete series is controlled at each real place
by the parameters of the two models. Alpöge's condition on them adds or removes
a block of consecutive numbers at the centre and moves the other entries, under
gap conditions that make the archimedean theta correspondence explicit.

% before: asm-dependence-rows-step-negated
The row step lifts from one row space to the next with the negated form on the
target. Negating the form of a space changes the signs of a model of discrete
series and of its Cartan involutions, and leaves its vectors and its
multiplicity maps unchanged.

% before: asm-dependence-rows-step-frames
The Fock vectors of the row step are formed with the frames of the grid, while
the local theory at the real places uses the frames of the models. Two frames
with the same Cartan involutions give the same Fock vectors after a linear
change of the polynomials in the creation operators.

% before: asm-dependence-rows-step-models
For models with theta paired parameters and the frames of the models, the theta
lifts of multiplicity maps factor through one map of the local models. This is
where the archimedean theta correspondence enters the row step.

% before: asm-dependence-rows-step-lift
The row tower of Lemma~\ref{asm-dependence-rows-tower} lifts multiplicity maps
one row at a time. Each step is a case of a general lifting of multiplicity
maps along a pair of discrete series models whose parameters correspond under
the archimedean theta correspondence.

% before: asm-dependence-inputs-sha-published
Part (9) of Theorem~\ref{asm-dependence-inputs} derives the Tate-Shafarevich
link in five steps. The first replaces the first and second collections of
published results by the forms in which the later steps use them.

% before: asm-dependence-inputs-sha-preconditions
The Galois preconditions and the fourth tie are first proved at the layer
stages. They concern the datum and its norm choice only, so they can be proved
before any bundle is chosen.

% before: asm-dependence-inputs-sha-transport-parts
A witness that a stage is built from the parameters is moved to new division
lattices and heights part by part. The parts below read the division data only
through objects that move with them, so each passes to the moved stage
unchanged.

% before: asm-dependence-inputs-sha-transport
A property proved at the layer stages must reach the bundles that the link data
attach to base points. This uses that the positions of the parameters of
dependence, and with them the witnesses, survive a change of the division
lattices and heights.

% before: asm-dependence-inputs-sha-cut
The Tate-Shafarevich link is proved through good cuts of the datum. At a good
cut the norm-choice facts hold, and the fourth tie survives the cut.

% before: asm-dependence-inputs-sha-link
The last step combines the Galois preconditions and the fourth tie at the linked
bundles with the good cuts, family by family. The twist condition of the
designed families supplies the remaining hypothesis on the twists.

% before: asm-dependence-inputs
The proof of Theorem~\ref{asm-definite-side-exists} combines results that rest
on the published inputs on unitary groups and on the theta correspondence. They
are gathered here so that the proof can follow its order: grids first, then
the linked properties at the parameters, then the choice of the class.

% before: asm-dependence-all-rows
Once the open parameter statements and the functional and covolume bounds are
linked, the items of a collection that concern the selected bundles follow one
by one. Theorem~\ref{is-definite-rows} supplies the definite rows from nine of
these properties.

% before: asm-class-selection
A property linked at parameters with the member property gives, in the end, a
class of families together with a selection of bundles for it. Siegel's mass
formula enters through the design property of that class.

% before: asm-definite-side-exists
The collection of definite-side hypotheses is built from the published
results on unitary groups and on the theta correspondence and from the
covolume bound of Chapter~\ref{chapter-covolumes}. The only additional input
is a local bound at the places above $2$ and $3$, which yields the
Tate-Shafarevich link.

% before: asm-zero-free-primitive
The tied families give the window bound of
Chapter~\ref{chapter-inert-prime-balance} at every large scale of a family.
For a primitive character, the hosts of Chapter~\ref{chapter-host-fields}
provide such families on windows of exponent $\lambda$. The windowed form of
Landau's theorem in Chapter~\ref{chapter-windowed-landau} then gives the
zero-free half-plane, including the point $s=1$ for the trivial character.

% before: asm-every-character-green
For a primitive character the hosts exist on windows of exponent $3/2$,
provided the Tate-Shafarevich bound holds. Putting the previous theorems
together gives the half-plane for primitive characters. The Euler factors at
the primes dividing the modulus extend it to all characters.

% before: asm-two-three-local-data
The hypothesis on the local bound at $2$ and $3$ in
Theorem~\ref{asm-definite-side-exists} compares, at each local index of a
definite-side datum, the congruence length of a Hecke point with the length of
a Selmer group. Both lengths are read off from the local data that the datum
carries at its twisted places. The killing condition bounds these Selmer
lengths and the values $\zeta_{i,v}$ of the characters $\chi_i$ on tame
inertia.

% before: asm-local-bound-two-three
The bound allows an error of size $fmb\log x$ in the congruence length, besides
the Tate-Shafarevich lengths of $E$ at $\ell$ and the $\ell$-part of its class
number. It is asked only at the local indices over $2$, and over $3$ when $3$
does not divide $Q_S$, and only for hub data that satisfy the setup. A family
satisfies it when one constant serves at all large scales.

% before: asm-two-three-cost
At a single member the bound comes with an explicit cost, polynomial in the
ranks and at most linear in the size $h_B$ and in $s_\ell=\max(1,v_\ell(N))$. Along
a family the ranks grow at most polynomially in $x$ and $b=\exp(x^{\varepsilon})$,
so the cost fits under the error allowed in
Definition~\ref{asm-local-bound-two-three}.

% before: asm-two-three-split-numerics
At a prime that splits in $\mathbf Q(\sqrt{-d})$, the local analysis at the
places over $\ell$ bounds the congruence length by a sum of local terms, indexed
by a partition of these places into row places and column places. The
following computation bounds that sum by the cost of
Lemma~\ref{asm-two-three-cost}, once the local terms are known to be small.

% before: asm-two-three-local-killing
One of the local terms, $h^0_u(\pi)$, is the length of the invariants of a
decomposition group above $\ell$ in a torsion module. The comparison data move
it to a Galois module over a finite extension of $\mathbf Q_\ell$ that is
contained in a cyclotomic field. There, a crystalline representation with
nonzero Hodge-Tate weights in a bounded window has torsion invariants killed by
a bounded power of $\ell$.

% before: asm-two-three-local-invariants
The invariant-killing property, applied at one place above $\ell$, bounds
$h^0_u$ by the rank $r$ times the exponent it gives. The completions of the
host field have at most $4n$ embeddings, so the Hodge-Tate weights of all
embeddings together form one set of at most $4nr$ integers.

% before: asm-two-three-local-zeta
Hypothesis (3) of Lemma~\ref{asm-two-three-split-numerics} asks that the
characters $\chi_i$ be nontrivial on the tame generators at the places of
$Q_i$. Since $\gamma_{i,v}$ generates the inertia group together with the wild
inertia, this follows from the ramification of $\chi_i$ at $v$ once $\chi_i$ is
trivial on the wild inertia.

% before: asm-two-three-killing-wach-frame
The torsion invariants of a crystalline lattice are bounded after extending
scalars to $\mathbf C_\ell$. There a matrix $P$ conjugates the action of an open
subgroup of the Galois group on the lattice to an action through a matrix
cocycle $U$ with entries in a cyclotomic extension, close to the identity at a
generator of the cyclotomic quotient. Each vector of a Hodge-Tate basis is fixed
by this action twisted by a power of the cyclotomic character.

% before: asm-two-three-killing-wach-tate-sen
In a Hodge-Tate basis, a vector fixed by the twisted action $*_i$ has
coordinates fixed by the Galois action twisted by the characters
$\chi^{i-u_k}$. When $i$ is not a weight of the basis these twists are
nontrivial, and the theorem of Tate and Sen shows that such coordinates vanish.

% before: asm-two-three-killing-wach-torsion-trace
The torsion bound of Lemma~\ref{asm-two-three-killing-wach-torsion} first
passes to the subgroup $G_{K_\infty}$, where the twisted action is the
coordinate action. There the input is Tate's theorem that the cyclotomic tower
is deeply ramified: traces from finite extensions of $K_\infty$ reach every
element of the maximal ideal of $K_\infty$.

% before: asm-two-three-killing-wach-torsion-average
A trace element $b$ with trace $a$ lets one average a cocycle over coset
representatives. The average is a coboundary up to the values of the cocycle on
a small open subgroup, which continuity makes as small as one likes.

% before: asm-two-three-killing-wach-torsion-tate
Repeating the averaging step with errors that halve each time gives a
convergent sequence. Its limit shows that every element of the maximal ideal of
$K_\infty$ kills the first cohomology of $G_{K_\infty}$ with integral
coefficients, with coboundaries of integral vectors.

% before: asm-two-three-killing-wach-torsion-fixed
The same trace element, with trace $\ell$, approximates an invariant element of
$\mathbf C_\ell$ by elements of $K_\infty$. So the integral invariants of
$G_{K_\infty}$ form the closure of the integers of $K_\infty$, as in the theorem
of Ax, Sen and Tate.

% before: asm-two-three-killing-wach-torsion-cyclotomic
On $G_{K_0}$ the cyclotomic character measures the levels of the cyclotomic
tower: an element fixes $K_j$ when its value is congruent to $1$ modulo
$\ell^j$. The remaining parts are the $\ell$-adic estimates for powers that the
next lemmas use.

% before: asm-two-three-killing-wach-torsion-dense
An element $g$ of $G_{K_m}$ with $|\chi(g)-1|=\ell^{-m}$ generates the quotient $G_{K_m}/G_{K_\infty}$ topologically. Its powers
approximate the value of $\chi$ at any element of $G_{K_m}$, and compactness
turns this into density.

% before: asm-two-three-killing-wach-torsion-procyclic
By density, a continuous cocycle is determined on $G_{K_m}$ by its values on
$G_{K_\infty}$ and at the generator $g$. This is fact (T4) of the torsion bound.

% before: asm-two-three-killing-wach-torsion-decomposition-ramification
Tate's decomposition of the completed integers of $K_\infty$ is built one layer
$K_n/K_{n-1}$ at a time, and it rests on the absolute values in each layer.
Over the unramified field $K_0$, the $\ell^j$-th roots of unity generate a
totally ramified extension with uniformizer $\zeta-1$, so the absolute values of
$K_j$ are the powers of $|\zeta-1|$.

% before: asm-two-three-killing-wach-torsion-decomposition-layer
Over $K_{n-1}$ the field $K_n$ is generated by a primitive $\ell^n$-th root of
unity $\zeta$, whose $\ell$-th power lies in $K_{n-1}$. In an expansion in
powers of $\zeta-1$ the terms have distinct absolute values, so the integers of
$K_n$ are the integral combinations of $1,\zeta,\ldots,\zeta^{\ell-1}$ over
$K_{n-1}$, and these powers are independent over $K_{n-1}$.

% before: asm-two-three-killing-wach-torsion-decomposition-complement
The terms with a nonzero power of $\zeta$ give a complement of the integers of
$K_{n-1}$ in those of $K_n$. Summing these complements over the layers above
$K_m$ and taking the closure gives the group $X_m$, which meets $K_m$ only in $0$
and is stable under multiplication by the integers of $K_m$ and under $G_{K_0}$.

% before: asm-two-three-killing-wach-torsion-decomposition-operator
A power $g^{\ell^b}$ of the generator multiplies a primitive
$\ell^{m+b+1}$-th root of unity by a primitive $\ell$-th root of unity. Summing
over the powers of $g$ therefore inverts $g-1$ on each layer up to the factor
$\zeta-1$, and an estimate that loses one factor of $\ell$ carries this to the
closure $X_m$.

% before: asm-two-three-killing-wach-torsion-decomposition
The action of $g$ on the completed integers of $K_\infty$ is studied through
Tate's decomposition: the integers of $K_m$ plus a closed complement $X_m$
spanned by roots of unity of higher level. On $X_m$ the operator $g-1$ can be
inverted after multiplication by $\zeta-1$, for a primitive $\ell$-th root of
unity $\zeta$.

% before: asm-two-three-killing-wach-torsion-perturbation
The twisted action of $g$ differs from the plain action by a matrix that is
congruent to $1$ modulo $\zeta-1$ times a small matrix. A geometric series
transfers the invertibility of $g-1$ on $X_m$ to the twisted action.

% before: asm-two-three-killing-wach-torsion-eigen
On the integers of $K_m$ the twisted action of $g$ is linear, given by the
matrix $U(g)$. The vectors of a Hodge-Tate basis are eigenvectors of $U(g)$,
with eigenvalues the powers $\chi(g)^{-u}$ of the weights, so a product over the
weights annihilates $U(g)$.

% before: asm-two-three-killing-wach-torsion-coinvariants
The coinvariants of $g$ are now split along Tate's decomposition. On the integers
of $K_m$, the annihilating polynomial divides out every eigenvalue other than
that of the twist, at the cost of the factor $\ell^{\nu_m(J,i)}$; on $X_m$ the
perturbation lemma inverts $g-1$ directly. This is fact (T3).

% before: asm-two-three-killing-wach-torsion
A vector that is invariant modulo a lattice gives a cocycle with integral
values, which becomes a coboundary after multiplication by a power of $\ell$.
The power needed depends only on the level and on the weights of the
Hodge-Tate basis. The proof reduces to the quotient $G_{K_m}/G_{K_\infty}$, which
is procyclic with generator $g$.

% before: asm-two-three-killing-wach-twisted
Together, the two results bound the denominators of a vector that is invariant
under the twisted action modulo a lattice between $\mathcal O_{\mathbf C}^r$
and $\ell^{-1}\mathcal O_{\mathbf C}^r$. One factor of $\ell$ makes the cocycle
integral, the torsion bound contributes $\ell^{1+\nu_m(J,i)}$, and the
remaining invariant vector vanishes by the theorem of Tate and Sen.

% before: asm-two-three-killing-wach-frame-exists-data
Alpöge builds the frame from a Wach module of the lattice, which the local
Hodge-Tate data of the datum provide abstractly. To each lattice representation
they attach a module over a ring of power series with a Frobenius and an action
of a quotient of $G_K$, a lattice in it, the Wach module, and a comparison with
the lattice over a period ring.

% before: asm-two-three-killing-wach-frame-exists-clauses
The construction of the frame uses five properties of a Wach module: a basis
compatible with the comparison, bounds by powers of $Q$ for the Frobenius in
both directions, stability under $G_{K_s}$ with $\gamma-1$ divisible by $\pi$,
and functoriality. Together with crystallinity over $G_{K_s}$ and a window for
the Hodge-Tate weights, they form the hypothesis from which the frame is built.

% before: asm-two-three-killing-wach-frame-exists-inputs-basis
The Wach module of a crystalline lattice is free. It is projective over
$\mathbf Z_\ell[[\pi]]$, so $\ell$ and $\pi$ act on it as a regular sequence,
and over the two-dimensional regular local ring $\mathcal O[[\pi]]$ this forces
freeness.

% before: asm-two-three-killing-wach-frame-exists-inputs-heights-expansions
The bounds for the Frobenius of a Wach module are read off from expansions at
roots of unity of $\ell$-power order. The datum expands each letter of $A$ as a
power series, and each element of the Wach module as a sequence of crystalline
periods, compatibly with Frobenius.

% before: asm-two-three-killing-wach-frame-exists-inputs-heights-points
At a primitive $\ell^s$-th root of unity the element $Q$ vanishes to order
exactly one. This makes the order of vanishing of an expansion a count of the
powers of $Q$.

% before: asm-two-three-killing-wach-frame-exists-inputs-heights-division
Divisibility by powers of $Q$ in $A$ is tested at the points where $Q$
vanishes. An element whose expansions all vanish to order $k$ at these points
is a multiple of $Q^k$.

% before: asm-two-three-killing-wach-frame-exists-inputs-heights-letters
The crystalline periods split into letters, one for each embedding of $K_0$,
and Frobenius moves each letter onto the next. Each letter has the dimension of
$V$, which lets a spanning family of the right size serve as a basis.

% before: asm-two-three-killing-wach-frame-exists-inputs-heights-standard-values
The Frobenius formula for the expansions involves the expansions of powers of
$Q$. At a root of unity of order dividing $\ell^{s-1}$ these are units, and at
every root of unity of $\ell$-power order they are nonzero; both facts are read
off from the constant term and the coefficient of $X$.

% before: asm-two-three-killing-wach-frame-exists-inputs-heights-standard-generator
The group $\Gamma$ acts on the expansions in degree $u$ through the character
$\chi_\Gamma^u$. Separating the degrees of an expansion needs an element whose
character is not a root of unity, and the generator $\gamma_0$ is one, because
$\chi_\Gamma(\gamma_0)$ is close to $1$ but different from it.

% before: asm-two-three-killing-wach-frame-exists-inputs-heights-standard-descent
Fullness of the expansions passes down along Frobenius, from $(i,\zeta)$ to
$(i-1,\zeta^\ell)$. The Frobenius formula relates the two lattices of
expansions, and a power of $Q$ moves the whole Wach module into the part where
the formula applies.

% before: asm-two-three-killing-wach-frame-exists-inputs-heights-standard-graded
When the expansions are full, the lattice of expansions is determined by its
single terms: a sequence lies in it exactly when each of its terms does. The
generator $\gamma_0$ scales the term in degree $u$ by a factor that depends on
$u$, and this separates the terms of a finite sum.

% before: asm-two-three-killing-wach-frame-exists-inputs-heights-full-embedding
At a primitive $\ell^s$-th root of unity $\zeta$ the expansions are compared
with de Rham periods through a map $\kappa$ from $L$ to $B_{\mathrm{dR}}$. It
must invert an embedding $\tau$ of $K$ that lies over $\tau_i$ and send $\zeta$
to $\zeta_{\ell^s}$; the irreducibility of the cyclotomic polynomial over the
unramified field $K_0$ provides one.

% before: asm-two-three-killing-wach-frame-exists-inputs-heights-full-evaluation
Power series over $L$ can be evaluated in $B_{\mathrm{dR}}$ at the uniformizer
$t_s$, compatibly with the filtration. This turns expansions, which are
sequences, into de Rham periods.

% before: asm-two-three-killing-wach-frame-exists-decompletion-frobenius-tilt
Integrality over $A_{\mathrm{inf}}$ will be read off from Witt components in
$\mathcal O_{\mathbf C}^\flat$ and its fraction field. Two properties of the
tilt enter: it is a valuation ring, and the Frobenius is bijective on it and on
its fraction field.

% before: asm-two-three-killing-wach-frame-exists-decompletion-frobenius-residue
On first Witt components, a Frobenius relation becomes an equation between
$\ell$-th powers of the entries of a matrix and affine expressions in them. Over
a valuation ring, the entry of largest value shows that such a matrix is
integral.

% before: asm-two-three-killing-wach-frame-exists-decompletion-frobenius-witt
Integrality of first components lets one split off a Teichmüller matrix and
divide the rest by $\ell$, and the quotient satisfies a relation of the same
kind. Repeating this matches the Witt components of the matrix, one at a time,
with components of matrices over the valuation ring.

% before: asm-two-three-killing-wach-frame-exists-decompletion-frobenius
Integrality over $A_{\mathrm{inf}}$ is detected by Frobenius: a matrix over
$W(\mathbf C_\ell^\flat)$ whose Frobenius is an affine expression in it with
coefficients in $A_{\mathrm{inf}}$ already has entries in $A_{\mathrm{inf}}$.
After a choice of basis of $\mathcal O_L$ over $\mathbf Z_\ell$, the same holds
with coefficients in $\mathcal O_L$.

% before: asm-two-three-killing-wach-frame-exists-inputs-heights-full-image-sandwich
Up to a power of $\varpi=\iota_A(\varphi^{s-1}(\pi))$, the Wach module and the
lattice generate the same module over $R_A$. A basis of the Wach module obeys
bounds by a power of $Q$ for Frobenius in both directions, and as
$\varphi(\varpi)=\varpi\,\iota_A(Q)$, Frobenius integrality turns these bounds
into the two inclusions.

% before: asm-two-three-killing-wach-frame-exists-inputs-heights-full-image
The comparison of the Wach module with de Rham periods sends each element of
$N(T)$ to a period whose expansion in powers of $t_s$ is its sequence of
expansions. These periods generate the de Rham lattice of $V$ along $\kappa$.

% before: asm-two-three-killing-wach-frame-exists-inputs-heights-full
At a primitive $\ell^s$-th root of unity, the expansions contain every sequence
in the letter that vanishes below a fixed bound. This lets single terms be
moved within the lattice of expansions.

% before: asm-two-three-killing-wach-frame-exists-inputs-heights-basis
The expansions of a basis of the Wach module form a basis of the lattice of
expansions. A shift of an expansion by $k$ inside this lattice then measures
the order of vanishing of the coordinates.

% before: asm-two-three-killing-wach-frame-exists-inputs-heights-standard-one
At $\zeta=1$ the expansions are described directly. Fullness descends there
from a primitive $\ell^s$-th root of unity, and the action of $\gamma_0$ shows
that the expansions in each nonzero degree are controlled by those in lower
degrees, which forces them to vanish in negative degrees.

% before: asm-two-three-killing-wach-frame-exists-inputs-heights-standard
At a root of unity of order dividing $\ell^{s-1}$, the expansions of the Wach
module carry no information beyond the letter: they are all sequences in the
letter that vanish in negative degrees. The Frobenius formula moves the
expansions at level $s$ to this level.

% before: asm-two-three-killing-wach-frame-exists-inputs-heights-filtration-component-split
The parts of de Rham periods along the embeddings of a field are cut out by
idempotents. When $L'$ receives $[E:F]$ embeddings of $E$, the algebra
$L'\otimes_FE$ is a product of copies of $L'$, and its idempotents split any
space with commuting actions of $E$ and $L'$ into the parts on which $E$ acts
through a single embedding.

% before: asm-two-three-killing-wach-frame-exists-inputs-heights-filtration-component-descent
An element of the $\tau$-part of $D_{\mathrm{dR}}(V)$ that lies in step $a$ of
the filtration of $V\otimes_{L,\kappa}B_{\mathrm{dR}}$ already lies in
$V\otimes\operatorname{Fil}^aB_{\mathrm{dR}}$. The parts of
$V\otimes B_{\mathrm{dR}}$ along the embeddings of $L$ into $\bar{\mathbf
Q}_\ell$ that lie over $\tau$ are permuted by $G_K$, which fixes the element and
preserves the filtration, and the other parts do not meet the $\tau$-part.

% before: asm-two-three-killing-wach-frame-exists-inputs-heights-filtration-component
On the letter $D_i(V)$, the Hodge filtration along $\tau$ can be read in
$V\otimes_{L,\kappa}B_{\mathrm{dR}}$, where $\kappa$ inverts $\tau$. This
turns the filtration conditions on expansions into conditions on de Rham
periods.

% before: asm-two-three-killing-wach-frame-exists-inputs-heights-filtration-single
A single term $\delta_kd$ lies in the lattice of expansions exactly when the
period of $d$ lies in the step $-k$ of the de Rham filtration along $\kappa$.
The periods of a basis of the Wach module form a basis of the de Rham lattice,
so the filtration condition can be solved for the coefficients.

% before: asm-two-three-killing-wach-frame-exists-inputs-heights-filtration
At a primitive $\ell^s$-th root of unity the lattice of expansions is the de
Rham lattice cut out by the Hodge filtration along one embedding of $K$. This
is where the Hodge-Tate weights enter the bounds.

% before: asm-two-three-killing-wach-frame-exists-inputs-heights
The window of Hodge-Tate weights controls the Frobenius of the Wach module, one
letter $e_i$ at a time. Alpöge proves this by expanding the Wach module at
roots of unity and comparing the expansions with the de Rham lattice.

% before: asm-two-three-killing-wach-frame-exists-inputs-ratio
The Wach data give crystallinity over $K$, and the construction needs it over
the larger field $K_s$. Fontaine's linear algebra for $B_{\mathrm{cris}}$
transfers it, starting from the fact that a ratio of crystalline periods fixed
by $G_F$ is a constant in $F_0$.

% before: asm-two-three-killing-wach-frame-exists-inputs-independent
From the ratio property, a relation with the fewest nonzero coefficients shows
that invariant crystalline periods independent over $F_0$ stay independent over
$B_{\mathrm{cris}}$. Counting them in the free module
$V\otimes B_{\mathrm{cris}}$ bounds the dimension of $D_{\mathrm{cris},F}(V)$.

% before: asm-two-three-killing-wach-frame-exists-inputs-restrict
A basis of $D_{\mathrm{cris},F}(V)$ lies in $D_{\mathrm{cris},F'}(V)$ for a larger
field $F'$ and stays independent over its maximal unramified subfield. Together
with the dimension bound over $F'$, this carries crystallinity from $F$ to $F'$.

% before: asm-two-three-killing-wach-frame-exists-inputs
The frame is built from four properties of the Wach data of the datum:
crystallinity over the open subgroup $G_{K_s}$, a basis of the Wach module of a
crystalline lattice compatible with the comparison, bounds for its Frobenius
one letter $e_i$ at a time, and the decomposition of the Wach module into its
letters. The window of Hodge-Tate weights enters only through the bounds.

% before: asm-two-three-killing-wach-frame-exists-wach
The bounds for the letters $e_i$ add up to bounds for the whole Wach module,
and the remaining clauses are properties of the Wach data. So the Wach module
of a crystalline lattice with Hodge-Tate weights in a window satisfies the Wach
hypothesis.

% before: asm-two-three-killing-wach-frame-exists-decompletion-sandwich
The bounds by powers of $Q$ in the Wach clauses turn into a comparison over
$A_{\mathrm{inf}}$ between the image of the Wach module and the lattice. The
element $x=\iota_A(\varphi^{s-1}(\pi))$ satisfies $\varphi(x)=x\,\iota_A(Q)$,
so its powers absorb the denominators that Frobenius integrality excludes.

% before: asm-two-three-killing-wach-frame-exists-decompletion-lift
The comparison places $x^{M_0}$ times the image of the Wach module inside the
lattice over $A_{\mathrm{inf}}$. The resulting map keeps the module structure
and, up to the factor $x^{M_0}$, the Galois action, and its image generates
$x^{M_0+h_0}$ times the lattice.

% before: asm-two-three-killing-wach-frame-exists-decompletion-specialization
Specializing $A_{\mathrm{inf}}$ to $\mathcal O_{\mathbf C}$ by a Frobenius twist
of Fontaine's map sends $x$ to $\zeta_{\ell^{n+1}}-1$. On the lattice over
$A_{\mathrm{inf}}$ it gives an integral, Galois-equivariant map to
$\mathbf C_\ell^d$ that sends the lattice to the standard coordinates.

% before: asm-two-three-killing-wach-frame-exists-decompletion-tower-surjective
A generator of $G_{K_s}$ modulo $G_{K_\infty}$ is found by prescribing its
cyclotomic character. Over the unramified field $K_0$ the cyclotomic polynomials
stay irreducible, so the character takes every value in $\mathbf Z_\ell^\times$.

% before: asm-two-three-killing-wach-frame-exists-decompletion-tower-equivariant
The specialized idempotents of $A$ are fixed by $G_K$ because $\iota_A$ is
equivariant. The trivial lattice representation of rank one, whose Wach module
is spanned by an element mapping to $1\otimes e_0$, shows this.

% before: asm-two-three-killing-wach-frame-exists-decompletion-tower-idempotent
An idempotent of $\mathcal O_{\mathbf C}\otimes\mathcal O_L$ has algebraic
coordinates, since every embedding of $L$ sends it to $0$ or $1$. If its
coordinates are fixed by $G_K$, they therefore lie in $K$.

% before: asm-two-three-killing-wach-frame-exists-decompletion-tower-gamma
On the ring $A$, an element of $G_{K_s}$ acts through a power series in $\pi$
whose linear coefficient is congruent to $1$ modulo $\ell$. So it moves every
power of the ideal $(\ell,\pi)$ into the next one.

% before: asm-two-three-killing-wach-frame-exists-decompletion-tower-iterate
An $\ell$-power iterate of a map that raises a filtration by one step is close
to the identity. The binomial coefficients of $(1+X)^{\ell^b}$ below degree
$\ell^{b-1}$ are divisible by $\ell^2$, and the higher terms raise the
filtration by $\ell^{b-1}$ steps.

% before: asm-two-three-killing-wach-frame-exists-decompletion-tower
The construction also uses arithmetic of the cyclotomic tower: a generator of
$G_{K_s}$ modulo $G_{K_\infty}$, values of the specialized ring $A$ in the
valuation ring of $K_m$, and the size of $\ell$ against powers of
$\zeta_{\ell^{n+1}}-1$. The level $m=2s+n+1$ is chosen so that these fit
together.

% before: asm-two-three-killing-wach-frame-exists-decompletion-matrix
A matrix $\Delta$ that conjugates the action of $G_{K_m}$ into matrices over
the valuation ring of $K_m$, close to $1$ at a generator and trivial on
$G_{K_\infty}$, turns this action into a frame. The integrality of
$\ell\Delta^{-1}$ comes from a relation $\Delta X=\ell$ with $X$ integral.

% before: asm-two-three-killing-wach-frame-exists-decompletion
Under the Wach hypothesis, specializing the Wach module at a root of unity of
level $m$ conjugates the action of $G_{K_m}$ on the lattice to an action
through a frame of rank $d$. The conjugating matrix is integral, its inverse is
integral up to a factor $\ell$, and it makes the action of $\mathcal O_L$
integral over $K_m$.

% before: asm-two-three-killing-wach-frame-exists-component-idempotent
The part of $V$ along an embedding $\tau$ of $K$ is cut out by an idempotent of
$K\otimes_{\mathbf Q_\ell}L$. Written through a separability element of $K$
over $\mathbf Q_\ell$, it is a matrix with entries in $K$.

% before: asm-two-three-killing-wach-frame-exists-component-factor
Passing to the $\tau$-part requires factoring an idempotent matrix through
integral matrices of smaller size. Over a valued field, a nonvanishing minor of
maximal size and largest absolute value gives such a factorization.

% before: asm-two-three-killing-wach-frame-exists-component-symbols
De Rham periods, shifted by powers of $t$ and specialized by Fontaine's map
$\theta$, give a basis of $\mathbf C_\ell^d$ on which $G_K$ acts through powers
of the cyclotomic character. The weights that survive the idempotent of the
$\tau$-part are Hodge-Tate weights along $\tau$.

% before: asm-two-three-killing-wach-frame-exists-component-cut
Cutting the frame of rank $d$ by the factored idempotent gives a frame of
smaller rank. The idempotent commutes with the action of $G_{K_m}$, which makes
the cut action multiplicative, and the Hodge-Tate basis projects onto one of
the cut frame.

% before: asm-two-three-killing-wach-frame-exists-component-coordinates
Condition (3) asks for integrality over $\mathbf Z_\ell$ to be detected after
the cut. Coordinates in an $\mathcal O_L$-basis of $T$, moved by an
automorphism $\iota$ of $\bar{\mathbf Q}_\ell$, detect integrality, and $L$
acts on them through $\iota$.

% before: asm-two-three-killing-wach-frame-exists-component
The frame of rank $d$ mixes the embeddings of $K$. An idempotent cuts out the
part along one embedding $\tau$, and de Rham periods give a Hodge-Tate basis of
that part, with weights the Hodge-Tate weights along $\tau$.

% before: asm-two-three-killing-wach-frame-exists-derham-descent
The de Rham property can be checked after a finite Galois extension. A
trace-dual basis of the extension splits the de Rham periods over the larger
field into copies of those over the smaller one.

% before: asm-two-three-killing-wach-frame-exists-derham
The Wach hypothesis gives crystallinity only over the open subgroup $G_{K_s}$,
while the Hodge-Tate basis of the $\tau$-part needs $V$ to be de Rham over $K$.
Over $K_s$ a crystalline representation is de Rham, and Galois descent from
$K_s$ to $K$ finishes the argument.

% before: asm-two-three-killing-wach-frame-exists
For a crystalline lattice over a subfield of $K_s$, Alpöge builds the frame
from a Wach module of the lattice. The level $m$ depends on $s$ and on the
width $M_0+h_0$ of the window of Hodge-Tate weights, and the weights of the
Hodge-Tate basis are Hodge-Tate weights of the representation. Condition (3)
brings an integrality statement over $\mathbf C_\ell$ back to the lattice over
$\mathbf Z_\ell$.

% before: asm-two-three-killing-wach
The invariant-killing property is proved over a field generated by roots of
unity of $\ell$-power order over an unramified field. Over such a field,
Alpöge bounds the exponent that kills the torsion invariants of a crystalline
lattice in terms of its Hodge-Tate weights, the level $s$ and the degree of a
cyclotomic extension.

% before: asm-two-three-killing-count
The exponent of Lemma~\ref{asm-two-three-killing-wach} grows with the number
of Hodge-Tate weights and with their $\ell$-adic valuations. When the weights
lie in a set of at most $n$ integers in $[-h,h]$, it is at most the exponent
$c_\ell(N,h,n)$ of the invariant-killing property.

% before: asm-two-three-killing-holds-lift-block
The split-case bound of Lemma~\ref{asm-two-three-rows-split-bound} needs, besides
the invariant-killing property, two further local properties at $\ell$. One
lifts crystalline classes along a surjection of lattices at the cost of a power
of $\ell$; the other controls the crystalline classes of the twists by
$\chi^{-e}$, $0\le e<t$, that a block decomposition produces.

% before: asm-two-three-killing-holds-exponents
The lifting property is proved over the field generated by the roots of unity of
order $\operatorname{lcm}(\ell,N)$, where Alpöge's lifting theorem holds with an
explicit power of $\ell$. Its exponent has a part that grows with the window of
the Hodge-Tate weights and a sum over the weights shifted by $1,\ldots,M$.

% before: asm-two-three-killing-holds-exponent-bound
The exponent depends on the Hodge-Tate weights only through their number and
their $\ell$-adic valuations. Both are controlled by the length of the window
that contains the weights, so the exponent is at most linear in the window and
in $s$.

% before: asm-two-three-killing-holds-wach-lift
Over a field generated over its unramified part by the $\ell^s$-th roots of
unity, a crystalline class with values in the target lattice of a surjection
lifts, after multiplication by a power of $\ell$, to a crystalline class with
values in the source lattice. Alpöge proves this with the Wach modules of the two
lattices, and the power of $\ell$ is $\ell^{e^P_\ell}$.

% before: asm-two-three-killing-holds-transfer
Passing from the cyclotomic field back to $K$ uses the transfer from an open
subgroup of finite index. Its formulas on cocycles show that restriction
followed by transfer is multiplication by the index, up to a coboundary.

% before: asm-two-three-killing-holds-descent
Restriction and transfer both preserve crystalline classes, so a lifting
statement over a finite extension descends to the base field. The cost is the
$\ell$-adic valuation of the degree, the part of the index that is not a unit.

% before: asm-two-three-killing-holds-block-top
The block property concerns the class of $g\mapsto\varphi(ge_0-e_0)$ for a
vector $e_0$ on which an invariant functional $p$ takes the value $1$. Over the
cyclotomic field, Alpöge bounds the power of $\ell$ that makes this class
crystalline, using the Wach module of the kernel of $p$.

% before: asm-two-three-killing-holds-block-transfer
The bound for the block classes holds after restriction to the cyclotomic
field. A transfer brings the class back to $K$, at the cost of the
$\ell$-adic valuation of the degree.

% before: asm-two-three-killing-holds-block-twists
The second half of the block property concerns the twists of the cyclotomic
character themselves. For a twist by $\chi^{-e}$ with $e>0$ every crystalline
class is a coboundary, for the trivial twist the crystalline classes are those
that are coboundaries on inertia, and the divisible classes that become
crystalline after multiplication by $\ell^m$ form a module of length linear in
$m$.

% before: asm-two-three-killing-holds
In an $N$-cyclotomic pair $(K,L)$, the field $K$ lies in the field of
$\operatorname{lcm}(\ell,N)$-th roots of unity, which is generated over its
unramified part by roots of unity of $\ell$-power order. Restriction to this
larger field keeps a representation crystalline with the same Hodge-Tate
weights, so Lemma~\ref{asm-two-three-killing-wach} applies there, and
Lemma~\ref{asm-two-three-killing-count} compares the two exponents.

% before: asm-two-three-crystalline-jump-sets
The Hodge-Tate weights of the hub representations at a place over $\ell$ are
read off from the Harish-Chandra parameter of the datum at one embedding. After
the twist by $\chi_i^{-1}$, they fill the cloud set, the low set and a block of
$t$ consecutive integers, possibly reflected about $(t-1)/2$. Removing the
block, which comes from the characters $\chi_i\varepsilon_\ell^{1-c}$, leaves the
weights of $W'$.

% before: asm-two-three-crystalline-jump-numerics
When the size data are admissible, the cloud set lies far above the block and
the low set lies on both sides of it at distance at least $(t+3)/2$ from its
centre. This separates the three pieces, so each weight occurs once. The windows
of the size data then bound all weights by $2b$.

% before: asm-two-three-crystalline
Lemma~\ref{asm-two-three-local-invariants} needs crystallinity and a window
for the Hodge-Tate weights of $W'$ at every place over $\ell$ and for every
extension of the coefficients. Crystallinity comes from that of the hub
representations $r_{i,x}$, since $\rho_{i,\pi}$ and $\chi_i$ are constituents
of $r_{i,\pi}$. The window comes from Alpöge's jump conditions, which locate
the Hodge-Tate weights.

% before: asm-two-three-wild-powers
At a place of residue characteristic $p$, an element $\sigma$ of the wild
inertia behaves like an element of $p$-power order: if $\sigma^n$ lies in an
open normal subgroup for some $n$ prime to $p$, then so does $\sigma$. So a
continuous character whose value at $\sigma$ has order prime to $p$ is trivial
at $\sigma$.

% before: asm-two-three-wild
Lemma~\ref{asm-two-three-local-zeta} needs $\chi_i$ to be trivial on the wild
inertia at the places of $Q_i$. The twisting data write $\chi_i$ as a
character of $\ell_i$-power order times a character whose ramification away
from $\ell_i$ is tame, of order dividing $15$ or $g_k-1$.

% before: asm-two-three-rows-partition
The local terms of Lemma~\ref{asm-two-three-split-numerics} are indexed by a
partition of the places over $\ell$ into column places and row places. When
$\ell$ splits in $\mathbf Q(\sqrt{-d})$, the image of $\sqrt{-d}$ in a
completion is one of two distinct square roots of $-d$ in $\mathbf Z_\ell$, and
complex conjugate places receive different ones. This sign decides on which
side of the CM type the embeddings at a place fall.

% before: asm-two-three-rows-split-bound-lifting
The split-case bound asks that crystalline classes lift along surjections at
the places over $\ell$, for representations built from the hub representations
by sums, subrepresentations and quotients. Such representations stay
crystalline, with Hodge-Tate weights among those of the pieces, so the lifting
property at $\ell$ applies to them.

% before: asm-two-three-rows-split-bound-block
The split-case bound also asks for the conclusions of the block property at
the row places, for the Hodge-Tate weights of $r_{i,\pi}\otimes\chi_i^{-1}$.
Collecting these weights over all embeddings of $E_w$ gives one set $U$ of
bounded size, to which the block property at $\ell$ applies.

% before: asm-two-three-rows-split-bound-invariants
At the complex conjugates of the column places, the split-case bound uses the
invariants of the dual of $W'$ twisted by $\varepsilon_\ell$. When this
representation is crystalline with no Hodge-Tate weight $0$, the
invariant-killing property bounds them as it bounds those of $W'$ in
Lemma~\ref{asm-two-three-local-invariants}.

% before: asm-two-three-rows-split-bound
For such a partition, Alpöge bounds the congruence length by the Selmer length
and local terms at the column and row places. The bound holds at any local index
whose prime is not a level prime, splits suitably in the host field and has the
three local properties, for any $N$ for which the local fields are
$N$-cyclotomic. It uses global Artin reciprocity, local Tate duality and the
local Euler characteristic formula, together with a long list of local
conditions on the hub representations, which the setup supplies. The same
hypotheses bound the invariants at the column places and make $\chi_i$
nontrivial at the tame generators of $Q_i$, the two remaining inputs of
Lemma~\ref{asm-two-three-split-numerics}.

% before: asm-two-three-rows-cyclotomic
The split-case bound is applied with $N=Q_Sqd$. The host field lies in the
field of $N$-th roots of unity, and the comparison data put the
$\operatorname{lcm}(\ell,N)$-th roots of unity into the coefficient fields, so
every local pair at $\ell$ is $N$-cyclotomic.

% before: asm-two-three-rows
The primes $2$ and $3$ are treated together: both split in
$\mathbf Q(\sqrt{-d})$, neither is a level prime, and at both
$s_0(i)=s_\ell=\max(1,v_\ell(N))$. The published local results at $\ell$ give
the bound of the congruence length by local terms, and the setup makes those
terms small.

% before: asm-local-bound-two-three-holds
The bound at a member and the polynomial growth of the ranks along a family
give one constant for each family. A property that holds for every family holds
eventually along every class, and in that form
Theorem~\ref{asm-definite-side-exists} uses the bound.

% before: asm-every-character-theta
The hypotheses of Theorem~\ref{asm-every-character-green} reduce to published
results. Green's parametrization is contained in the published results on
unitary groups, and the local bound at $2$ and $3$ follows from them by
Proposition~\ref{asm-local-bound-two-three-holds}. What remains is the
list of results from which the collection of definite-side hypotheses is
built.

% before: asm-theta-agree-doubling
Theorem~\ref{asm-every-character-theta} assumes the agreement of global and
local theta integrals and the local matching. Alpöge derives both from two
published results on doubled unitary groups: the meromorphic continuation of
Siegel Eisenstein series and the Siegel-Weil formula. Their statements use
the doubled group of a lattice space, its frames with their standard compact
subgroups, sections with their Eisenstein sums, and a model of the Weil
representation for a doubled pair.

% before: asm-theta-agree-eisenstein
The Eisenstein sum of a flat family converges only in a right half-plane.
Langlands' theorem continues it meromorphically to the whole plane. The form
used here takes sections at the level of a standard compact subgroup and
bounds the order of the pole at each point uniformly against growth
functionals.

% before: asm-theta-agree-siegel-weil
When one space of the pair is definite, the theta integral over its unitary
group is an integral over a compact quotient. The Siegel-Weil formula
identifies the average of the theta sum over this quotient with the value
of the continued Eisenstein series at a point determined by the two ranks,
up to a factor $2$ when the rank of the definite space is at most that of
the other.

% before: asm-theta-agree-models
Definition~\ref{th-theta-hypotheses} describes the agreement of theta
integrals in words. Its derivation needs the objects behind those words: an
explicit Schrödinger model of the Weil representation for a pair of lattice
spaces, the polynomial Fock vectors at the infinite places, local vacua at the
finite places, and the theta integral over a fundamental domain. The full-rank
data, whose component at one split place has a linear model supported on the
matrices of maximal rank, are the functions of the second space on which the
agreement is stated.

% before: asm-theta-agree-triple
A triple consists of two Hermitian spaces and a skew-Hermitian space whose rank
is the sum of theirs, with characters that fix the splittings of the two Weil
representations. The rank and signature conditions put the first pair in the
range where its theta integrals converge, and they leave room to remove
hyperbolic planes. Re-choosing the archimedean data and the bases of a triple
changes its polynomial Fock spaces but not its theta integrals.

% before: asm-theta-agree-matched
The agreement compares the two theta integrals of a triple on functions whose
images under local maps agree. These maps take values in irreducible
representations of the unitary group of the skew-Hermitian space, and at the
ramified places they are pinned down by the zeta functions of
Lemma~\ref{th-bottom-split-zeta}. Correspondence is the polynomial form of this
comparison, and nondegeneracy is the property that the derivation proves by
induction.

% before: asm-theta-agree-planes
The induction behind the derivation removes a hyperbolic plane from one of the
spaces of a triple. Its tools are the constant term along the parabolic
subgroup of an isotropic line, which splits into two parts with prescribed
laws on the Levi elements, and the cones, which average a function over the
lattice stabilizer and over the ideles to pass to the smaller triple.

% before: asm-theta-agree-parts
With these objects the agreement and the local matching of
Definition~\ref{th-theta-hypotheses} can be stated as six properties of triples
and matched quadruples. The step
factors and the partial $L$-values measure how the constant of a triple
changes when a hyperbolic plane is removed.

% before: asm-theta-agree-real-fock
The real maps of a triple are built at each infinite place from the Fock model
of the Weil representation, in which Schwartz functions are replaced by
polynomials. The two unitary Lie algebras act on the polynomials by quadratic
expressions in raising and lowering operators, and the target of the real maps
is described by the action of the Lie algebra of the skew-Hermitian space.
Frames adapt the matrices of the lattice spaces at an infinite place to the
standard form of the right signature.

% before: asm-theta-agree-real-fock-relations
The two actions on the Fock space are the infinitesimal forms of the two
members of a dual pair. They commute, and the diagonal matrices act on
monomials with rational weights, so their quotients are modules of the kind
just defined.

% before: asm-theta-agree-real-quotient-lie-module
The two modules of Lemma~\ref{asm-theta-agree-real-quotient} are quotients of
Fock spaces on which $M_{R+S}(\mathbf C)$ acts, and each is generated by one
vector of highest weight for the block-diagonal matrices. The argument for a
common quotient uses only a finite-dimensional space around that vector, with
an invariant inner product, and a gap in the spectrum of the Casimir operator
of the block-diagonal matrices on the vectors one off-diagonal step away.

% before: asm-theta-agree-real-quotient-filtration
A module generated by a finite-dimensional space $\iota(A)$ is filtered by the
number of off-diagonal matrices applied to $\iota(A)$. The Casimir gap makes
the compact Casimir operator, shifted by its value on $A$, injective on each
step of the filtration. So a vector that the shifted operator sends into
$\iota(A)$ already lies in $\iota(A)$.

% before: asm-theta-agree-real-quotient-local
In a module generated by a highest weight vector, the vectors killed by some
polynomial in the compact Casimir operator that does not vanish at the
Casimir value of the highest weight form a subspace. The highest weight vector
is not in it. By Lemma~\ref{asm-theta-agree-real-quotient-filtration}, every
proper stable subspace is.

% before: asm-theta-agree-real-quotient-common
Two modules generated with Casimir gap in the same weight are both quotients
of the stable subspace of their product generated by the pair of generating
vectors. By Lemma~\ref{asm-theta-agree-real-quotient-local} that subspace has
a largest proper stable subspace, and the quotient by it is irreducible.

% before: asm-theta-agree-real-quotient-harmonics
The modules of Lemma~\ref{asm-theta-agree-real-quotient} are quotients of
polynomial rings, and the proof that they are generated with Casimir gap works
with the grading by degree. Each off-block operator splits into parts that
raise and lower the degree by two, and the polynomials killed by the lowering
parts are the harmonics.

% before: asm-theta-agree-real-quotient-shapes
The highest weights that occur among joint harmonics are described by four
sequences of integers. The Casimir gap comes down to an inequality between the
Casimir value of the generating weight and those of the weights one root away
from it.

% before: asm-theta-agree-real-quotient-fock-facts
The Fock space is spanned by the harmonics and by what the raising operators
produce, and the joint harmonics carry each pair of highest weights at most
once, in a degree fixed by the weights. These are results of Howe and of
Kashiwara and Vergne.

% before: asm-theta-agree-real-quotient-fock-rules
The proofs below also use how the two actions interact with the Fock inner
product, with the grading and with the harmonics. These are direct
computations with the operators of the variables.

% before: asm-theta-agree-real-quotient-split
In the quotient by the twisted first action, every homogeneous polynomial can
be replaced by one on which the block-diagonal matrices act by the character.
Below a certain degree there are none, and two degrees above it the harmonics
suffice.

% before: asm-theta-agree-real-quotient-generator
When the joint shape is unique, the joint harmonics on which the first action
is the character have one highest weight for the second action. A
highest-weight vector of that weight generates all of these joint harmonics
under the block-diagonal part of the second action.

% before: asm-theta-agree-real-quotient-generation
The image of these joint harmonics generates the whole quotient. The proof goes
up by degree, writing each homogeneous polynomial through the harmonic
decomposition.

% before: asm-theta-agree-real-quotient-root
The Casimir gap concerns the vectors one off-block step above the generating
degree. A highest-weight vector among them has a weight that differs from the
generating weight by one root between the two blocks.

% before: asm-theta-agree-real-quotient-gap
For a gap weight, the Casimir value of such a weight exceeds the generating
value by more than $R+S$. So the shifted Casimir operator is injective on the
relevant vectors, also after passing to the quotient.

% before: asm-theta-agree-real-quotient-side
The generating vector, the polynomials of the generating degree and the gap
lemma give a highest-weight datum with a Casimir gap. This is the form in which
Proposition~\ref{asm-theta-agree-real-quotient-common} uses the quotients.

% before: asm-theta-agree-real-quotient-weights
For the two quotients of Lemma~\ref{asm-theta-agree-real-quotient} the
generating weights are explicit. Dominance and the sign bounds reduce the
inequality for gap weights to a root from the last element of one block to the
first element of the other.

% before: asm-theta-agree-real-quotient-intertwiner
The quotients must also be shown to be nonzero. This uses a map between two
modules with invariant inner products, built from two vectors with the same
eigenvalues that the lowering matrix units kill.

% before: asm-theta-agree-real-quotient-nonzero
A linear form on the Fock space that vanishes on the span defining the quotient
and not everywhere shows that the quotient is nonzero. Such a form comes from
splitting the variables into two groups and pairing harmonics of the two groups
through the map of the previous lemma.

% before: asm-theta-agree-local-mult-archimedean-sides
The gap weights of Lemma~\ref{asm-theta-agree-real-quotient-weights} and the
nonvanishing of Lemma~\ref{asm-theta-agree-real-quotient-nonzero} apply to two
quotients of Fock spaces, one for each Hermitian space of a triple. Under
conditions on the ranks and on an even integer $k$, both quotients are
generated with Casimir gap, and the twist by $\tfrac k2\operatorname{tr}$ in the
second makes their weights equal. The common quotient of the next lemma and
the radicals of the theta forms are both built on these two modules.

% before: asm-theta-agree-real-quotient
At one infinite place, the two Hermitian spaces of a triple give two Fock
spaces for the same skew-Hermitian space. Dividing each by the action of the
Lie algebra of its Hermitian space, twisted by a character, leaves a module for
the Lie algebra of the skew-Hermitian space. The conditions on the ranks and on
$k$ are those under which the two modules have a common irreducible quotient.

% before: asm-theta-agree-real-one-place
The common quotient, with the composite maps from the two Fock spaces, is the
local datum at one infinite place. The maps are lifts with the twists that the
two characters of a triple impose at that place.

% before: asm-theta-agree-real-parity
The local datum at an infinite place uses half the difference of the exponents
of the two characters of a triple. This difference is even, because the
characters restrict to powers of the quadratic character whose exponents have
the same parity.

% before: asm-theta-agree-real-frames
The Fock model at an infinite place is defined relative to a basis in which the
Hermitian form is standard and the Cartan involution is diagonal. Such a basis
exists by the spectral theorem and Sylvester's law of inertia.

% before: asm-theta-agree-real-lines
The conditions on real maps involve the one-parameter subgroups of the unitary
groups at the infinite places and their maximal compact subgroups. In the
coordinates given by frames, a one-parameter subgroup is the exponential of an
element of the Lie algebra, and every element of the maximal compact subgroup
lies on such a subgroup.

% before: asm-theta-agree-real-weak
The conditions on real maps are stated with actions on all functions on the
adeles, while the Fock model gives derivatives against admissible linear forms.
The two formulations agree, and each determines the image of a given Schwartz
function.

% before: asm-theta-agree-real-derivative
The Heisenberg operators at the infinite places are the derivatives of the
Schrödinger operators. They are the operators of
Lemma~\ref{asm-theta-agree-density-derivative}, here in the weak form.

% before: asm-theta-agree-real-exponential
An element of a maximal compact subgroup acts on Fock vectors through the
exponential of the action of its Lie algebra. On polynomials this exponential
is a finite computation, since the operators involved preserve the polynomials
of bounded degree.

% before: asm-theta-agree-real-tensor
The target of the real maps is the tensor product of the local data over the
infinite places. Its irreducibility follows from that of the factors by the
Jacobson density theorem, since each factor has scalar endomorphisms.

% before: asm-theta-agree-real-fock-map
The Fock space of a model embeds into the Schwartz space at infinity: a
polynomial goes to the corresponding polynomial in the creation operators,
applied to a Gaussian vacuum. Its image is the polynomial Fock space of
Definition~\ref{asm-theta-agree-models}(4), on which the conditions on real
maps are imposed.

% before: asm-theta-agree-real-metaplectic
The derivatives of the Weil representation at the infinite places are compared
with those of the metaplectic representation of a real symplectic group on the
Schwartz space. Near the identity its operators exist and are differentiable,
with derivative a quadratic expression in the Heisenberg operators.

% before: asm-theta-agree-real-projective
A homomorphism into operators that agrees, up to scalars, with the metaplectic
operators has the derivative of the metaplectic representation up to a scalar
on each one-parameter subgroup. These scalars form a linear form on the Lie
algebra that vanishes on commutators; on a unitary Lie algebra such a form is a
multiple of the trace.

% before: asm-theta-agree-real-oscillator
The operators of a Weil representation implement the symplectic action on the
Heisenberg operators, so along a one-parameter subgroup their derivative
differs from the quadratic operator of the metaplectic representation by a
scalar. This scalar is linear on the Lie algebra and vanishes on commutators,
hence is a multiple of the trace at each infinite place.

% before: asm-theta-agree-real-lie
With these tools the action of the two Lie algebras on Fock vectors can be
computed. The scalar of the previous lemma is fixed by the doubling: the
doubled representation on a block-diagonal and on a diagonal one-parameter
subgroup, together with its rule on the Siegel parabolic subgroup, shows that
it is half the exponent of the character.

% before: asm-theta-agree-real
The derivation rests on twenty-six intermediate properties of triples, and on
two further properties that are themselves parts of the agreement. The first
two give the local matching: a triple has local maps matched to it, at the
infinite places and at each finite place.

% before: asm-theta-agree-finite-unique
The finite maps are built one place at a time, from operators that act at a
single place. An operator acts at a place by at most one function, which a
single test function at the other places detects.

% before: asm-theta-agree-finite-split-models
At a place where $E$ has an $F$-embedding into $F_v$, the unitary groups of the
triple become general linear groups, and the Weil representations become actions
on functions of matrices. The finite maps at such a place take values in a space
of functions on $GL_d(F_v)$ induced from the block upper triangular matrices.

% before: asm-theta-agree-finite-bernstein
The image of the finite maps at a split place must be irreducible under the
second unitary group. The induced space has this property by a theorem of
Bernstein, as it is induced from unitary characters of the determinants.

% before: asm-theta-agree-finite-iwasawa
At a split place outside $T$ the vacua of the two pairs must have the same
image. An element of the induced space is determined by its values on
$GL_d(\mathcal O)$, so it suffices to compare these values.

% before: asm-theta-agree-finite-zeta
The finite maps at a split place are zeta integrals over the first general
linear group. At the places of $T$ they are the zeta integrals of the zeta
condition. At the other split places they are normalized by the unramified
$L$-factor, so that both vacua go to the function with value $1$ on
$GL_d(\mathcal O)$.

% before: asm-theta-agree-finite-characters
The zeta integrals need the local components of the characters of the triple to
be unitary and smooth. Whether $\nu$ is ramified at the place decides which of
the two kinds of zeta integral is used.

% before: asm-theta-agree-finite-local-objects
At an inert place the Weil representations of the two pairs are studied through
the stabilizers of the lattices, a Lagrangian adapted to the lattice of $W$ with
its parabolic subgroup, and linear forms invariant along the Lagrangian. The
action of a global operator at one place is described by local operators.

% before: asm-theta-agree-finite-inputs-slice
Local operators at a place $v$ are obtained by restricting adelic functions to
the slice through a base point, on which only the component at $v$ varies. The
operators below multiply by a function of the $v$-component and move the
argument within such a slice; they span enough of the Schwartz-Bruhat space to
test membership and admissibility.

% before: asm-theta-agree-finite-inputs-schur
An adelic operator that commutes with the Schrödinger operators away from $v$
should act at $v$ by a local operator. The proof builds, for each point of the
slice, an operator that commutes with every Schrödinger operator, and the
irreducibility of the Schrödinger representation makes it a scalar.

% before: asm-theta-agree-finite-inputs-powers
The smoothness of the local operators reduces to a statement about matrices. A
matrix all of whose powers stay within $\tfrac12$ of the identity is the
identity: comparing the sum of its first $N$ powers with $N$ bounds $N$ times its
distance from the identity.

% before: asm-theta-agree-finite-inputs-smooth
The Weil representation is only known to be weakly continuous. On a group with
small subgroups, weak continuity forces each vector to be fixed by one of them,
since the orbit of the vector under a small subgroup spans a finite-dimensional
space.

% before: asm-theta-agree-finite-inputs-subgroups
The unitary groups of the lattice spaces have small subgroups made of elements
at $v$. They are cut out by bounds on the coordinates of the entries at $v$.

% before: asm-theta-agree-finite-inputs-bruhat-tits
At an inert place the unitary group of $W$ is an unramified unitary group, and
the stabilizer of the lattice is a hyperspecial maximal compact subgroup. The
Cartan and Iwasawa decompositions for it are the inputs to the statements on
the stabilizers at inert places.

% before: asm-theta-agree-finite-inputs-modulus
The index identity at an inert place compares the stabilizer in the parabolic
subgroup with its conjugates. The ratio of the two indices is the modular
character of the parabolic subgroup, which is computed on one diagonal element.

% before: asm-theta-agree-finite-inputs-frame
To apply the local facts, the unitary group of $W$ at an inert place is written
in a basis of the completed lattice in which the form is a multiple of the
antidiagonal form. Such a basis exists because unimodular Hermitian lattices
over an unramified extension are determined by their rank.

% before: asm-theta-agree-dictionary-change
The split-place dictionary is first obtained for one linear model and then moved
to the coordinates given by bases of the lattices. Replacing the matrix variable
by $PYQ$ turns one linear model into another and carries the readings of the
Weil operators along.

% before: asm-theta-agree-dictionary-conjugate
At a split place the two embeddings of $E$ give two local characters for each
Hecke character. When the restriction of the character to the ideles of $F$ is a
power of $\eta$, it is trivial on norms, and the two local characters are
inverse to each other.

% before: asm-theta-agree-dictionary-vacuum
The vacuum is fixed, up to a character, by the Schrödinger operators of the
lattice at $v$. In the coordinates of bases of the completed lattices, these
operators become translations by integral matrices and multiplication by
characters trivial on integral matrices, so a linear model carries the vacuum
to a multiple of the indicator function of the integral matrices.

% before: asm-theta-agree-dictionary-models
Combining the dictionary with a change of coordinates and the computation of
the vacuum gives linear models of both kinds in the bases of the lattices.
They carry the vacuum to the indicator function of the integral matrices and
turn the Weil operators into the two explicit actions of general linear groups.

% before: asm-theta-agree-finite-inputs-invariant
At an inert place an evaluation form along a Lagrangian is a linear form on the
local Schwartz-Bruhat functions that is invariant under the Schrödinger
operators of an isotropic subspace. Such a form exists for every isotropic
subspace: it integrates a function over the first coordinates of the subspace,
against a phase that makes the integral invariant.

% before: asm-theta-agree-finite-inputs-criterion
To show that an invariant form does not vanish at the vacuum, it suffices to
show that an invariant form vanishing there vanishes everywhere. By the
irreducibility of the Schrödinger representation this reduces to the translates
of the vacuum, which a decomposition of the symplectic space controls.

% before: asm-theta-agree-finite-inputs-uniqueness
For a Lagrangian subspace the invariant forms are unique up to a scalar, so a
group that preserves the subspace acts on them by a character. The proof
restricts functions to the first coordinates of the subspace and reduces to
forms invariant under translations.

% before: asm-theta-agree-finite-inputs-averaging
A second isotropic subspace, transverse to the Lagrangian, gives a new
invariant form by integrating the translates of a first one along it. The new
form changes under operators that move the transverse subspace linearly by the
modulus of that linear map, and so measures the character of the previous lemma.

% before: asm-theta-agree-finite-inputs-local-data
At an inert place the global Weil representation restricts to a Weil
representation over the unramified quadratic extension. Its local structure
consists of Darboux coordinates, a splitting of the unitary group, and a doubled
splitting that satisfies the Siegel rule for the diagonal Lagrangian.

% before: asm-theta-agree-finite-inputs-transport
The Siegel rule is part of the data only for the diagonal Lagrangian of the
doubled space. An element of the doubled unitary group carries the diagonal to
the sum of two copies of any Lagrangian, which moves the rule there; the box
property then brings it back to the space itself.

% before: asm-theta-agree-finite-inputs-fixed
When the character of the second space is trivial, the unitary group of the
first space fixes the evaluation forms along a Lagrangian of the second. Through
the doubled group it acts on them by a character, and the Siegel rule and the
averaged form show that this character is trivial on the stabilizer of the
diagonal Lagrangian, hence trivial.

% before: asm-theta-agree-finite-inputs-local-pair
The local lemmas apply to the Weil pairs of the triple through their components
at an inert place. Lifting elements of the local unitary groups to elements at
$v$ turns the global operators into a local Weil pair, and a global form
invariant along the diagonal Lagrangian, built from a local one, transfers the
Siegel rule.

% before: asm-theta-agree-finite-inputs
The constructions at the finite places use facts about the local Weil
representations. At a split place they are the matrix models of the two pairs;
at an inert place, the decompositions of the unitary group of $W$ and the
behaviour of the evaluation forms under its parabolic subgroup.

% before: asm-theta-agree-finite-split-core
Through the matrix models, the conditions on matched finite maps at a split place
become conditions on maps from Schwartz-Bruhat functions of matrices to an
irreducible space of functions on $GL_d(F_v)$. Maps invariant under the first
general linear groups and equivariant under the second give finite maps at $v$.

% before: asm-theta-agree-finite-ramified
At a place of $T$ the zeta condition fixes the finite maps: they are the zeta
integrals of that condition, read through the matrix models. Bernstein's theorem
shows that they map onto the induced space.

% before: asm-theta-agree-finite-split
At a split place outside $T$ the characters are unramified. The zeta integrals
normalized by the unramified $L$-factor take the vacua of both pairs to the same
function, which gives the condition on the vacua.

% before: asm-theta-agree-finite-inert-situation
At an inert place outside $T$ there are no matrix models, and the finite maps
come from spherical functions on the unitary group of $W$. These are built from
local operators for the two pairs at the place.

% before: asm-theta-agree-finite-howe
The spherical vectors of the unramified Weil representation are described by
Howe. The vacuum is spherical for both groups, and the spherical vectors for one
group are spanned by translates of the vacuum under the other.

% before: asm-theta-agree-finite-vacuum
Averaging over the compact stabilizers turns Howe's description into statements
modulo the augmentation space. Every function is then congruent to a combination
of translates of the vacuum, and a sum of its translates by the stabilizer of $W$
is congruent to a multiple of the vacuum.

% before: asm-theta-agree-finite-average
A linear form can be averaged over a compact group by finite sums, since every
Schwartz-Bruhat function is fixed by a subgroup of finite index. The average is
invariant, and it keeps any invariance of the form under operators commuting
with the group.

% before: asm-theta-agree-finite-count
The spherical functions of the two pairs are averages, over the stabilizer of
$W$, of functions that transform by characters of the parabolic subgroup. Their
equality is a counting identity in the group, which uses the symmetry of the
double cosets and the modulus of the parabolic subgroup.

% before: asm-theta-agree-finite-spherical
The finite maps at an inert place need spherical forms for the two pairs with
the same spherical function. They are averages of evaluation forms, and the
counting lemma compares their values.

% before: asm-theta-agree-finite-inert
With equal spherical functions, the finite maps at an inert place send a function
to its coefficient against the spherical form. Their common target is the span
of the translates of the spherical function.

% before: asm-theta-agree-finite
At a finite place the conditions on matched maps involve only that place.
The construction therefore splits into the places of $T$, the split places
outside $T$ and the inert places outside $T$.

% before: asm-theta-agree-decay-full
A full-rank datum is a finite sum of products of a Schwartz function at the
infinite places and a full-rank finite part. The theta function of each product
decays faster than any power of the quotient height in the first variable,
with a power of the height in the second variable that does not depend on the
product. Summing over the products gives the same decay for the datum.

% before: asm-theta-agree-decay
The theta integrals of a triple must converge before they can be compared.
For the full-rank data of the second space they also have moderate growth,
and the theta functions of these data decay rapidly on the unitary group of
the second space.

% before: asm-theta-agree-coinvariants-local
At a split place the local Weil representation can be written on functions of
matrices, where the unitary groups act as general linear groups twisted by
characters. A dictionary is a linear map that carries one action to the
other.

% before: asm-theta-agree-coinvariants-step
Modulo the augmentation space, a Schwartz-Bruhat function on matrices can be
moved, one row at a time, onto the matrices whose rows are linearly independent.
Each step uses row operations from the general linear group, and averaging over
finitely many of them; the nontrivial character makes the averaging invertible.

% before: asm-theta-agree-coinvariants-kill
A subspace that is stable under multiplication by enough locally constant
functions contains every Schwartz-Bruhat function that vanishes at a given point.
The functions used separate that point from every other one.

% before: asm-theta-agree-coinvariants-eigen
Through a linear model of the second kind, half of the Schrödinger operators
act at a split place by multiplication by characters. A linear form on which
these operators act by a fixed character is then determined by the value of the
model at a single matrix.

% before: asm-theta-agree-coinvariants-slices
Membership in the full-rank data can be checked one slice at a time. A function
all of whose slices at $v_1$ have full rank is a sum of pure tensors with a
full-rank factor at $v_1$.

% before: asm-theta-agree-coinvariants-fibre
The $d\times n$ matrices of rank $n$ fibre over the $n$-dimensional subspaces of
$k^d$, through their column spaces, and $GL_d$ acts transitively on these
subspaces with a parabolic subgroup as stabilizer. So the functions on these
matrices form a representation induced from that parabolic subgroup, and the
functions on $GL_n$, the fibre, carry the inducing representation.

% before: asm-theta-agree-coinvariants-haar
The twisted coinvariants of the functions on $GL_n$ under right translation are
one-dimensional. Integration against the twisting character is the linear form
that realizes them.

% before: asm-theta-agree-coinvariants-exact
Parabolic induction should commute with taking these coinvariants. It does,
because an induced function is a combination of translates of functions
supported near the parabolic subgroup, and averaging over compact subgroups
preserves the relations.

% before: asm-theta-agree-coinvariants-irreducible
After passing to the coinvariants of $GL_n$, the induced representation becomes
one induced from two characters of the Levi factor. Bernstein's irreducibility
theorem then leaves only two possibilities for a stable subspace.

% before: asm-theta-agree-dictionary-split
The dictionary at a place $v$ where $E$ has an $F$-embedding $t$ into $F_v$
rests on the splitting $E\otimes_FF_v=F_v\times F_v$. An element at $v$ is then
a pair of elements of $F_v$, and an element of a unitary group at $v$ is
determined by its image under $\pi_t$, which can be any invertible matrix.

% before: asm-theta-agree-dictionary-stone-von-neumann
Over a nonarchimedean local field the Schrödinger representation on
Schwartz-Bruhat functions is irreducible, and any two Darboux coordinates on the
same symplectic space give isomorphic Schrödinger representations. These two
facts produce the linear models and pin down the operators of the Weil
representation up to a scalar.

% before: asm-theta-agree-dictionary-evaluation
The scalar is found by evaluating a linear form at one point. A linear form on
Schwartz-Bruhat functions that is unchanged by multiplication with every
additive character is a multiple of evaluation at $0$.

% before: asm-theta-agree-dictionary-situation
The dictionary is proved for a pair of lattice spaces and again for the doubled
pair and for the pair with the two spaces exchanged. The setup that these share
consists of the place, the embedding, two matrices and an isomorphism of
Darboux coordinates.

% before: asm-theta-agree-dictionary-coordinates
At $v$ an element of $M_{n\times d}(\mathbf A_E)$ is a pair of matrices over
$F_v$, read through $t$ and through $t\circ c$. In these coordinates the
symplectic form becomes the standard one, and the right action of the unitary
group of $W$ at $v$ becomes right multiplication by its image under $\pi_t$.

% before: asm-theta-agree-dictionary-chart
A Heisenberg chart transports the local Schrödinger operators of the model to
the standard Schrödinger operators on functions of matrices. It is the form of
a linear model that the computation of the Weil operators uses.

% before: asm-theta-agree-dictionary-charts
The Stone-von Neumann theorem gives a Heisenberg chart, and a Heisenberg chart
is the same as a linear model of the second kind. Exchanging the two spaces of
the pair transposes the matrices, which lets the computation for the group of
$W$ also serve for the group of $X$.

% before: asm-theta-agree-dictionary-schur
An adelic operator that moves the Schrödinger operators by a symplectic map which
is the identity away from $v$ acts at $v$. Composed with the inverse of a local
operator that moves them in the same way, it commutes with all Schrödinger
operators, so it is a scalar by the adelic form of Schur's lemma.

% before: asm-theta-agree-dictionary-diagonal
The scalar is computed with the doubled pair, through a linear form invariant
under the diagonal Lagrangian. Integration over the diagonal is such a form, and
it does not vanish on a suitable product function.

% before: asm-theta-agree-dictionary-scalar
For an element of the unitary group of $W$ at $v$ the Weil operator moves the
Schrödinger operators by right multiplication. In a Heisenberg chart the same is
done by right translation of the matrix variable, so the Weil operator acts at
$v$ by a multiple of right translation.

% before: asm-theta-agree-dictionary-action
The multiple is found through the doubled group. The element $(h,1)$ of the
doubled group is conjugate to an element of the Siegel parabolic subgroup by an
element at $v$, and the Siegel rule for that element gives the value of $\chi$
and the power of the absolute value at the determinant.

% before: asm-theta-agree-coinvariants-dictionary
At a place where $E$ has an $F$-embedding into $F_v$, the unitary groups at $v$
become general linear groups. A linear model of the second kind then turns the
Weil representation into the twisted action on functions of matrices.

% before: asm-theta-agree-coinvariants-inputs
The description of the coinvariants rests on four local statements. Two concern
the twisted action of general linear groups on functions of matrices, one
supplies dictionaries at split places, and one is the stability of the
full-rank data.

% before: asm-theta-agree-coinvariants-split
At a clean place the unitary groups should become general linear groups. This
needs an embedding of $E$ into the completion of $F$ at that place.

% before: asm-theta-agree-coinvariants-nontrivial
The local statements apply when the first character of the dictionary is
nontrivial. At a clean place this comes from the infinity type of $\nu$, whose
exponents are positive at every infinite place.

% before: asm-theta-agree-coinvariants-kernel
Given a dictionary with a nontrivial character, a linear map that vanishes on
the augmentation space and is compatible with the action of $U_W$ has the
augmentation space as its kernel. Every local component is also congruent to
one of full rank.

% before: asm-theta-agree-coinvariants
At the clean places the matched maps should be the coinvariants of the local
Weil representations. This identification lets a component at $v_1$ be
replaced by one of full rank without changing its image.

% before: asm-theta-agree-convergence-tori
Convergence of a theta integral is decided on a Siegel set, which up to
rational elements is a torus times a compact set. On the torus the theta
function is bounded by a lattice sum: a sum over the rational points of a
majorant that decreases rapidly at the archimedean places and has compact
support at the finite ones.

% before: asm-theta-agree-convergence-siegel
Reduction theory supplies the Siegel sets. For the unitary group of a
nondegenerate form it gives a torus from a Witt decomposition whose translates
by a compact set meet every orbit of the rational points, and it bounds
integrals over a fundamental domain by integrals over the cone of the torus,
weighted by the inverse of the modulus character.

% before: asm-theta-agree-convergence-signature
The signature of a lattice space bounds the size of its tori. The first vectors
of a Witt torus span a totally isotropic subspace at every infinite place, and
the dimension of such a subspace is at most the smaller entry of the signature.

% before: asm-theta-agree-convergence-theta
The theta function of a Weil pair is the theta sum of a translate. Weil's theta
distribution is continuous in the archimedean component, so the theta function
is continuous in the group variables and in the archimedean Schwartz function.

% before: asm-theta-agree-convergence-poisson
The product of two theta sums is bounded by a sum of coefficients over the
rational points of the symplectic space. This is Poisson summation, applied to
the product of one Schwartz-Bruhat function and a rational translate of the
other.

% before: asm-theta-agree-convergence-coefficients
The coefficients transform under the Weil representation by the linear action
on the symplectic space, since the operators intertwine the Schrödinger
operators and preserve the $L^2$-norm. For pure tensors they decay rapidly at
the archimedean places and have compact support at the finite ones.

% before: asm-theta-agree-convergence-square
Together, Poisson summation and the two properties of the coefficients bound
the square of a theta function by a lattice sum. The bound depends on the
archimedean component only through a continuous seminorm, which gives the
continuity in the archimedean component later.

% before: asm-theta-agree-convergence-lattice
Lattice sums are invariant under rational elements and change by a bounded
factor under compact translations. On the tori of the two groups they are
computed coordinate by coordinate in a basis adapted to both tori, by counting
the points of a lattice in each coordinate.

% before: asm-theta-agree-convergence-cone
On the torus of the first group the lattice sum grows exponentially in
$d\sum_is_i$, and the Siegel weight decays exponentially in $\delta_n(s)$. The
square root of the first is integrable against the second over the cone in the
range of signatures of Weil's criterion.

% before: asm-theta-agree-convergence-dominated
Reduction theory for the second group makes the bound by a lattice sum uniform
in the second variable, at the cost of a power of its height. In the first
variable the square root of the lattice sum is integrable over a fundamental
domain, by the Siegel integration bound and the cone integral.

% before: asm-theta-agree-convergence
For a single Weil pair, convergence holds in a range of signatures. In that
range the theta integral is linear on the Schwartz-Bruhat functions and
continuous in the archimedean component, which allows limits.

% before: asm-theta-agree-continuity-sums
The theta function of full-rank data sees only the rational matrices of full
rank, and this gives convergence in every signature. The estimates use lattice
sums over these matrices, the excess of one torus coordinate over the other,
and the finite parts whose products with Schwartz functions stay full rank.

% before: asm-theta-agree-continuity-pieces
Full-rank data are built from Schwartz functions and full-rank finite parts,
so estimates for a single finite part suffice. A full-rank finite part is in
turn a combination of factorized functions whose component at $v_1$ has full
rank, and this local factor is where the full rank enters.

% before: asm-theta-agree-continuity-cut-local
At the place $v_1$ the full-rank condition is a condition on the support in a
linear model. A function supported on the matrices of full rank can be cut into
pieces on each of which a suitable character is constant, while that character
separates the piece from every matrix of lower rank.

% before: asm-theta-agree-continuity-eigen
A function that is an eigenvector of a Heisenberg operator has coefficients
that change by a character under the corresponding translation. At a point
where that character is not $1$, the coefficient vanishes.

% before: asm-theta-agree-continuity-cut
The cut and the eigenvector computation together show that each piece of a
full-rank local function has coefficients vanishing at every point whose
component at $v_1$ has rank less than $n$. In the square bound for the theta
function, this removes the matrices of lower rank from the lattice sum.

% before: asm-theta-agree-continuity-fourier
The square bound pairs a function with its reflection, which must then have the
same theta function. The first step is that the coefficients of two
Schwartz-Bruhat functions determine the product of their values, by Fourier
uniqueness on the adeles.

% before: asm-theta-agree-continuity-parity
The element $-1$ of $U_X$ acts by the reflection up to a scalar, because the
two operators have the same coefficients. As $-1$ is central, a function and
its reflection then have the same theta function.

% before: asm-theta-agree-continuity-square
For a full-rank finite part, the square of the theta function is bounded by a
lattice sum over the rational matrices of full rank. The theta function times
that of the reflection is bounded by a sum of coefficients, and the full-rank
pieces kill the terms of lower rank.

% before: asm-theta-agree-continuity-full-sums
The full-rank lattice sums keep the two properties of the full sums that are
used with reduction theory. They are invariant under rational elements, and
moving the arguments within compact sets changes them by a bounded factor.

% before: asm-theta-agree-continuity-minor
On the Siegel torus, a matrix of full rank cannot be small in all its
coordinates when the first torus coordinate outruns the second. A nonzero
minor of bounded denominator bounds its size from below in terms of the excess.

% before: asm-theta-agree-continuity-lattice
The lower bound for the size becomes a saving in the full-rank lattice sum.
Raising the exponent of the majorant gains any power of the exponential of the
excess over the bound for the sum over all matrices.

% before: asm-theta-agree-continuity-decay
These estimates make the theta function of a full-rank finite part decay
rapidly in the first variable, with a power of the height in the second that
does not depend on the finite part. The excess pays for the growth along the
first torus, so no condition on the signature is needed.

% before: asm-theta-agree-continuity-volume
Decay in the first variable gives integrability over a fundamental domain once
that domain has finite volume. For the unitary group of a nondegenerate form
this is the theorem of Borel and Harish-Chandra.

% before: asm-theta-agree-continuity-domination
Decay and finite volume together dominate the theta function of a full-rank
finite part by an integrable function of the first variable. The power of the
height of the second variable does not depend on the finite part.

% before: asm-theta-agree-continuity-integrals
Such a domination makes theta integrals converge, add and depend continuously
on the archimedean component, by dominated convergence. The statement allows
finitely many finite parts at once, as the continuity lemma needs.

% before: asm-theta-agree-continuity
The full-rank data are approximated by sequences of polynomial data. Their
theta integrals converge to that of the limit, so identities proved for
polynomial data pass to the full-rank data.

% before: asm-theta-agree-density-moments
Density of the polynomial Fock space reduces to a statement about Schwartz
functions on a real vector space. A continuous linear form that kills a
Gaussian times every polynomial kills everything, by the Fourier inversion
formula.

% before: asm-theta-agree-density-derivative
The polynomial Fock space is built from the first-order actions of the
Schrödinger operators along lines at the infinite places. In coordinates these
actions are first-order differential operators with linear coefficients.

% before: asm-theta-agree-density-gaussian
For the polynomial Fock space to be large it must contain a Gaussian. The
Cartan involutions of the two spaces give a complex structure on the
archimedean part of the symplectic space, and the Gaussian is the vacuum of
that structure.

% before: asm-theta-agree-density-linear
Multiplication by a linear coordinate is one of the first-order actions of the
Schrödinger operators. So the polynomial Fock space contains a Gaussian times
every polynomial.

% before: asm-theta-agree-density
A theta integral that vanishes on the polynomial data for one choice of
archimedean data should vanish for every choice. Density of the polynomial
Fock space in the Schwartz space, with the continuity above, gives this.

% before: asm-theta-agree-partners
Polynomial data of either space have partners in the other space with the
same images under the matched maps. The proof only chooses preimages under
surjective maps, except at $v_1$, where the description of the coinvariants
supplies a full-rank component.

% before: asm-theta-agree-levi
The parts of a constant term are compared at Levi elements. Every element of
the unitary group of the smaller skew-Hermitian space extends to one, by a
direct construction in the basis given by the plane.

% before: asm-theta-agree-stabilizer
The cones integrate over the stabilizer of a lattice, which must be compact to
carry a finite Haar measure. Its compactness combines the definiteness given
by a Cartan involution at the infinite places with the compactness of the
lattice at the finite places.

% before: asm-theta-agree-evaluation-heisenberg
The evaluation maps are built by comparing Schrödinger models of one pair in
different Darboux coordinates. The comparison rests on Weil's results on the
adelic Heisenberg group: the theta distribution is the only invariant one, and
an operator commuting with all Schrödinger operators is a scalar.

% before: asm-theta-agree-evaluation-metaplectic
Two Darboux coordinates on the same symplectic space differ by an adelic
symplectic transformation. Weil's metaplectic operators intertwine the
Schrödinger operators along it, and they factor over the places, which is what
the definition of an evaluation map asks for.

% before: asm-theta-agree-evaluation-transfer
The metaplectic operators are only defined up to a scalar, so the comparison
is phrased for any operator that intertwines two Darboux coordinates. When the
coordinates are rationally related, such an operator multiplies the theta sum
by a nonzero constant, and it has an inverse of the same kind.

% before: asm-theta-agree-evaluation-rational
Regular models send rational matrices to rational coordinates. A dimension
count then shows that the coordinates of any two regular models of one pair
are rationally related.

% before: asm-theta-agree-evaluation-lattice
The vacua of a model are eigenvectors for the Schrödinger operators of a
lattice at each finite place. Two arithmetic properties of this lattice are
needed: the character it sees is trivial on integral traces, and it is its own
annihilator under the symplectic form.

% before: asm-theta-agree-evaluation-vacuum
A metaplectic operator carries the vacua of one model to vacua for the
coordinates of another, but not necessarily to the given ones. The vacuum at a
finite place is unique up to a scalar, so the operator can be normalized place
by place.

% before: asm-theta-agree-evaluation-change
These results give, for any two regular models of a pair, an evaluation map
between them that is bijective on Schwartz-Bruhat functions. It lets the
evaluation maps be built on a model adapted to the planes and then moved to the
given model.

% before: asm-theta-agree-evaluation-derivative
The polynomial Fock space is defined through derivatives of Schrödinger
operators at the real places. To compare Fock spaces across a change of model,
these derivatives must pass through the continuous component of the evaluation
map.

% before: asm-theta-agree-evaluation-pure
A change of model must carry the polynomial Fock space of one model onto that
of the other, so that polynomial data correspond. This follows once its
component at the real places is shown to be bijective and to intertwine the
derivative operators.

% before: asm-theta-agree-evaluation-compose
The evaluation maps are composites of a change of model and an evaluation map
on an adapted model. Composites of evaluation maps are evaluation maps, and a
map that factors over the places carries polynomial data onto polynomial data
when its components do.

% before: asm-theta-agree-evaluation-glue
On a model whose coordinates split along the plane, the evaluation map is
restriction to the coordinates that remain. The Heisenberg and invariance
conditions become identities between dot products.

% before: asm-theta-agree-evaluation-adapted-frame
The adapted models are built in coordinates attached to the planes. The vectors
of a plane and the columns of its embedding form a basis, and in this basis the
Hermitian form splits into the smaller form and a hyperbolic plane, compatibly
with the lattices.

% before: asm-theta-agree-evaluation-adapted-vacuum
At a finite place the vacuum of the adapted model is the vacuum of the smaller
model times the indicator function of a compact open subgroup in the remaining
coordinates. Three conditions on the coordinates make this product an
eigenvector of the lattice operators with the right eigencharacter.

% before: asm-theta-agree-evaluation-adapted-model
The coordinates of an adapted model are those of the smaller model together with
coordinates of a part of the space on which the Hermitian form is a pairing.
When such a splitting of the form is compatible with the lattices, it gives a
regular model whose vacua restrict to the vacua of the smaller model.

% before: asm-theta-agree-evaluation-adapted-signs
A Gaussian is characterized by a relation between pairs of elements at the
infinite places, twisted by the Cartan involutions. The relation respects the
symplectic form, and on rank-one elements the form has a sign fixed by the
positivity of the Cartan involutions.

% before: asm-theta-agree-evaluation-adapted-commutator
The positivity has an analytic counterpart. Two first-order operators whose
commutator is a negative multiple of the identity cannot satisfy the Gaussian
relation on a nonzero Schwartz function.

% before: asm-theta-agree-evaluation-adapted-fock
On an adapted model, restriction to the coordinates of the smaller pair commutes
with the derivative operators along the smaller space. In the remaining
coordinates a Gaussian is a Gaussian function, so restriction carries the
polynomial Fock space onto that of the smaller pair.

% before: asm-theta-agree-evaluation-adapted
The adapted models are regular models of the pair whose coordinates split along
the given planes. For the plane in the skew-Hermitian space alone, the
restriction must also carry the polynomial Fock space onto that of the smaller
pair.

% before: asm-theta-agree-evaluation-lagrangian-normal
A rational Lagrangian is described by an idempotent matrix and a symmetric
matrix over the field. Invariance under the Lagrangian then becomes invariance
under two explicit families of Schrödinger operators.

% before: asm-theta-agree-evaluation-lagrangian-finite
At the finite places the uniqueness is a statement about linear forms on
locally constant functions with compact support. Invariance under translations
and under characters that separate the points off the translations forces a
form to vanish once it vanishes on one indicator function.

% before: asm-theta-agree-evaluation-lagrangian-real
At the infinite places the same statement concerns continuous forms on a
Schwartz space. It holds for multiplications by characters, where the form is a
multiple of evaluation at zero, and the Fourier transform and multiplication by
quadratic phases carry that case to every Lagrangian.

% before: asm-theta-agree-evaluation-lagrangian
The uniqueness of the evaluation map at zero reduces to a statement about
linear forms on Schwartz-Bruhat functions. A form invariant under the
translates of a rational Lagrangian is determined up to a scalar.

% before: asm-theta-agree-evaluation-unique
Two evaluation maps at zero differ by a nonzero scalar on the polynomial data.
This is what makes the image of the polynomial data independent of the
evaluation map chosen.

% before: asm-theta-agree-evaluation
Removing a hyperbolic plane from the skew-Hermitian space is done on the
Schrödinger models by an evaluation map. It must carry the polynomial data of
the larger pair onto those of the smaller pair, so that nonvanishing can be
transported.

% before: asm-theta-agree-evaluation-cone
When the plane is removed from both spaces of a pair, two evaluation maps are
needed, one for each of the two parts of the constant term. Each is invariant
under the Schrödinger translations by its own set of matrices.

% before: asm-theta-agree-tower-isotropic
The induction removes a hyperbolic plane from a Hermitian and from a
skew-Hermitian lattice space. Such a plane starts from an isotropic vector,
and a form that is indefinite at every real place and has rank at least three
has one by the Hasse principle.

% before: asm-theta-agree-tower-pair
An isotropic vector has to be completed to a hyperbolic pair that is compatible
with the self-dual lattice. The vector is rescaled so that its pairings with
the lattice are integral at the places of $P$, and a partner is corrected by a
multiple of it until the partner is isotropic.

% before: asm-theta-agree-tower-cartan
An adapted plane also asks the Cartan involutions at the infinite places to
preserve the plane and its complement. Over the complex numbers, a hyperbolic
pair and a diagonalizing basis of the complement diagonalize the whole form,
and the sign involution of this basis is such a Cartan involution.

% before: asm-theta-agree-tower-basis
At the places of $P$ a lattice space carries a basis of the completed lattice,
and an adapted plane prescribes two of its columns. The splitting of the
lattice along the plane shows that these columns, together with a basis of
the smaller lattice, still form a basis.

% before: asm-theta-agree-tower-plane
The last three results give the plane, once the Cartan involutions and the
bases of the space are chosen again. The complement of the plane is a smaller
lattice space whose lattice is still self-dual and whose signature has lost
one positive and one negative direction.

% before: asm-theta-agree-tower-eigen
The smaller pairs of spaces need regular models, and a model needs a vacuum
at every finite place. A vacuum is a common eigenvector of the Schrödinger
operators of a compact lattice group, and it is found by averaging a
translate of a test function over a finite quotient of that group.

% before: asm-theta-agree-tower-cocycle
For a Hermitian and a skew-Hermitian lattice, the scalars in the definition of
a vacuum satisfy the rule of the previous lemma. The self-duality of the two
lattices makes the standard character trivial on the relevant values of the
trace form.

% before: asm-theta-agree-tower-vacuum
The two previous lemmas give a vacuum at every finite place. At almost all
places the coordinates are integral, and the indicator function of the
integral points is a vacuum, as a regular model requires.

% before: asm-theta-agree-tower-model
A regular model consists of rational Darboux coordinates, vacua, a Haar
measure and a fundamental domain. All four exist for every Hermitian and
skew-Hermitian lattice space.

% before: asm-theta-agree-tower-weil
The smaller triples also need Weil representations of the new models, with
the characters of the character datum. Their existence is that of the Weil
representation of a unitary dual pair, through Kudla's splitting.

% before: asm-theta-agree-tower-alpha
The induction lowers $p+q$. When the second Hermitian space is definite and the
first is not, a hyperbolic plane can be removed from the first space and the
skew-Hermitian space, after re-choosing the archimedean data and bases.

% before: asm-theta-agree-tower-beta
When the second Hermitian space is not definite, the plane is removed from it
and the skew-Hermitian space. This lowers $q$ by one and keeps $p$.

% before: asm-theta-agree-split-notions
The constant term of a theta integral along a plane is computed by sorting the
rational points of the theta kernel by a vector of the first space. The
isotropic vectors are moved to one of them by rational elements, and the
resulting sums are unfolded over an Iwasawa decomposition into integrals over
the ideles, the cone integrals.

% before: asm-theta-agree-split-low
An adapted plane is a hyperbolic plane with compatible lattice and archimedean
data. Removing it lowers both parts of the signature by one, so a space with an
adapted plane is indefinite.

% before: asm-theta-agree-split-witt
A definite Hermitian space has no isotropic vectors, and in an indefinite one
the rational unitary group moves every isotropic vector to every other. One or
two transvections suffice.

% before: asm-theta-agree-split-inputs-characters
An absolutely convergent series of characters of an abelian group determines
the coefficient of the trivial character. Averaging the series over translates
removes the nontrivial characters one at a time. This separates the terms of
the theta expansions below.

% before: asm-theta-agree-split-inputs-polar
A theta sum on a larger space can be expanded along a family of rational
directions, each term an evaluation map applied to a translate. The directions
must pair with a complementary family that the evaluation map does not see,
and both families must be orthogonal to the image of the smaller space.

% before: asm-theta-agree-split-inputs-expansion
Along a polar system, the theta sum of a function is a constant times the sum,
over the rational points of the polar directions, of the theta sums of its
evaluations. The same argument shows that an evaluation map is determined up
to a scalar by its data, and that the composite of two evaluation maps is
again one.

% before: asm-theta-agree-split-inputs-systems
For adapted planes in the two spaces, the sets of matrices that the evaluation
maps along the lines of $W$ and of $X$, and along $\mathcal J_\alpha$ and
$\mathcal J_\beta$, do not see carry polar systems. Their directions are built
from the second vectors $y$ and $y'$ of the planes.

% before: asm-theta-agree-split-inputs-kfinite
The unfolding over an Iwasawa decomposition needs data whose translates under
the lattice stabilizer span a finite-dimensional space. For polynomial data
this holds because the stabilizer commutes with the Cartan involutions, so it
preserves each step of the filtration of the polynomial Fock space by degree.

% before: asm-theta-agree-split-inputs-kill
Full-rank data are detected at the place $v_1$ through a linear model, in
which a Schrödinger operator at $v_1$ acts on the image by multiplication by a
character. A linear form with prescribed eigenvalues under enough such operators
therefore vanishes on the full-rank data, as soon as the eigenvalues force the
relevant matrices to have rank below $n$.

% before: asm-theta-agree-split-inputs-cone-haar
The unfolding of a sum over isotropic vectors leads to an iterated integral over
the parabolic subgroup fixing a line and over a compact subgroup. To compare it
with Haar measure on the whole group, it suffices that the measure it defines be
invariant on the left under the parabolic subgroup and on the right under the
compact one, when together they give the whole group.

% before: asm-theta-agree-split-inputs-cone-parabolic
Haar measure on the parabolic subgroup is built from Haar measures on its
unipotent group and on the two factors of its Levi elements. Conjugation by the
Levi elements rescales Haar measure on the unipotent group, and the rescaling
is the power of the norm that weights the unfolded integral.

% before: asm-theta-agree-split-inputs-cone-unfold
The sum over isotropic vectors runs over the cosets of the stabilizer of a line
in the rational unitary group, so its integral over a fundamental domain
unfolds. By the two previous lemmas the result is an integral over the compact
group, the ideles and a fundamental domain of the smaller unitary group.

% before: asm-theta-agree-split-inputs-cone-tate
The line integrals in the cone run over the ideles against a power of the norm.
Their convergence reduces to the absolute convergence of global zeta integrals
to the right of the critical strip.

% before: asm-theta-agree-split-inputs-cone-line
Moving a Schwartz-Bruhat function along a line of Schrödinger operators and then
evaluating gives a family indexed by the ideles. Its integral against a
character and a power of the norm is again a Schwartz-Bruhat function, and its
theta integrals are dominated, which allows the integrals in the cone to be
exchanged.

% before: asm-theta-agree-split-inputs-cone-exchange
The cone integral integrates theta integrals over the compact group and the
ideles, while the theta integral of the cone integrates last. For data with
finitely many translates under the compact group the two agree, so the
unfolded sum becomes a theta integral of the smaller pair.

% before: asm-theta-agree-split-inputs-xunique
Maps that intertwine the Schrödinger operators along the columns of $\iota_H$
and are invariant under the matrices $x_Hq^{\mathsf T}$ are determined up to a
scalar, as on the side of $W$. This is how the unipotent group of $H$ is shown
to act trivially through an evaluation map.

% before: asm-theta-agree-split-inputs-xcommutator
The unipotent group of $H$ is generated by elements fixing $x_H$ and moving the
columns of $\iota_H$ along $x_H$. A commutator with a scaling element keeps only
the part of such an element that moves $y_H$ along $x_H$, and conjugation by the
scaling elements multiplies that part by squares.

% before: asm-theta-agree-split-inputs-heis
The laws of an evaluation map under the two unitary groups come from a
uniqueness argument for maps that intertwine Schrödinger operators. Such a map
is described by a symplectic change of coordinates followed by a restriction to
some of the coordinates, and the conditions below are what that argument uses.

% before: asm-theta-agree-split-inputs-heis-axioms
Four properties of Schwartz-Bruhat functions on a power of the adeles carry the
argument: closure under the basic operations, the description of evaluation at
the origin by its invariance, Weil's operators for symplectic automorphisms, and
Haar integration. Each follows from a result stated earlier.

% before: asm-theta-agree-split-inputs-heis-lines
From these properties, a map that obeys the law of a symplectic automorphism
and a restriction is unique up to a scalar, and such maps exist. The same
argument shows that the forms they define see every nonzero function after a
Schrödinger operator.

% before: asm-theta-agree-split-inputs-heis-double
A group acting on one side is compared with the evaluation map through the
doubled space, where the diagonal is a Lagrangian. An evaluation map for the
doubled data splits off one for the original data, and an average over the
coordinates outside the restriction measures how an operator scales the
diagonal.

% before: asm-theta-agree-split-inputs-heis-scalars
If an operator and an operator on the smaller space transport the law of an
evaluation map into itself, uniqueness forces the evaluation map to intertwine
them up to a scalar. For a group acting on both sides, these scalars form a
continuous character.

% before: asm-theta-agree-split-inputs-side
The arguments use one of the two unitary groups at a time, together with the
doubling that the Weil representation carries. Both groups of a Weil pair give
such data, with the roles of the two spaces exchanged.

% before: asm-theta-agree-split-inputs-siegel-situation
Over a finite product of local fields, a character of a unitary group that is
trivial on the stabilizer of a Lagrangian is shown to be trivial in three
steps: on the opposite unipotent elements, on a big cell, and then everywhere.
The first two steps use only an anti-involution $\sigma$ of a ring, the
elements $x$ with $\sigma(x)x=1$ and a pair of complementary idempotents.

% before: asm-theta-agree-split-inputs-siegel-unipotent
Conjugation by a diagonal element of the group carries a unipotent element
$1+s$ to $1+s/4$, whose fourth power is $1+s$ again. So a character takes a cube
root of unity at $1+s$, and the same holds at its cube root $1+s/3$, which forces
the value $1$.

% before: asm-theta-agree-split-inputs-siegel-cell
An element $g$ for which $ge+f$ is invertible factors as a unipotent element
$1+s$ times an element of the parabolic subgroup, with $s$ built from the
inverse of $ge+f$. A character trivial on both kinds of element is therefore
trivial on this big cell.

% before: asm-theta-agree-split-inputs-siegel-finite
Over an Artinian algebra, a matrix that depends quadratically on a parameter and
is invertible at one value is invertible at all but finitely many values in the
base field. The argument passes to the finitely many residue fields.

% before: asm-theta-agree-split-inputs-siegel-generic
Every element of the group is the quotient of two elements of the big cell.
They are found as matrix polynomials in the element and its inverse with a
parameter from the base field, which is infinite, so that a value avoiding the
finitely many bad ones exists.

% before: asm-theta-agree-split-inputs-siegel-artinian
For a unitary group over an Artinian algebra over a field of characteristic
$0$, the three previous lemmas apply to the anti-involution
$x\mapsto c(J^{-1}x^{\mathsf T}J)$ and to the idempotents that project onto a
Lagrangian and onto a complementary one. The next lemma applies this to the
doubled group over a finite product of completions.

% before: asm-theta-agree-split-inputs-siegel
The characters produced by doubling are continuous and trivial on the Siegel
parabolic subgroup of a doubled group, and the arguments need them trivial on
the whole group. Elements that are trivial away from finitely many places
approximate every element, so the question reduces to a doubled group over a
finite product of local fields.

% before: asm-theta-agree-split-inputs-seesaw
Doubling turns the action of a unitary group on one side into a character of
the doubled group, which the Siegel rule and the previous lemma make trivial.
When the plane lies in the other space, an average over the complementary
coordinates compares the exponents $a/2$ and $a_1/2$ of the two Siegel rules;
when it lies in the same space, the Levi element contributes its own factor.

% before: asm-theta-agree-split-inputs-charts
The evaluation maps along planes fit the abstract setting after a change of
coordinates adapted to the plane. In these coordinates the columns of $\iota$
give the coordinates to restrict to, and the vectors along $x$ and $y$ give the
two complementary halves.

% before: asm-theta-agree-split-inputs-xvacuum
An evaluation map along a plane of $X$ must carry vacua to vacua. In the
adapted coordinates, a vacuum restricts to a nonzero multiple of the vacuum of
the smaller model.

% before: asm-theta-agree-split-inputs-xexists
An evaluation map along a plane of $X$ is built in two steps. Weil's operator
moves the functions to the adapted coordinates, and restriction to the
coordinates of the smaller space, normalised at the vacua, finishes the map.

% before: asm-theta-agree-split-inputs
The computation rests on structural facts about the groups and on the constant
term of the theta kernel itself. The kernel's constant term along the plane of
$W$ is a theta function of the smaller space plus a sum over the isotropic
vectors of the first space, and that sum unfolds into a cone integral.

% before: asm-theta-agree-split-kernel
Expanding the theta kernel along the line of $x_{H'}$ writes it as a sum over
the vectors $x$ of the first space. Under the integral over the unipotent
group the anisotropic vectors drop out, the zero vector gives the smaller
theta function, and the isotropic vectors are all moved to $x_H$.

% before: asm-theta-agree-split-unique
The two laws of a splitting differ in their exponents, $a'/2<a/2$. An idele of
absolute value larger than one separates them, so a splitting is unique.

% before: asm-theta-agree-split-fubini
The constant term of a theta integral is computed under the integral over the
first group. Exchanging the two integrations needs a bound on the theta kernel
over a compact set of unipotent elements, which the moderate growth gives.

% before: asm-theta-agree-split-domains
The theta integrals of two models of the same first space use different Haar
measures and fundamental domains. For automorphic functions the two integrals
differ by a positive constant.

% before: asm-theta-agree-split-zero
The term of the kernel's constant term at the zero vector integrates to a theta
integral of the smaller pair. Its law comes from the evaluation formula at the
Levi elements of $W$.

% before: asm-theta-agree-split-shift
The terms at nonzero isotropic vectors are theta functions of shifted
Schwartz-Bruhat functions. They are invariant under the rational stabilizer of
$x_H$, and along an orbit of $E^\times$ they become the shifts by the multiples
of a fixed rank-one matrix.

% before: asm-theta-agree-split-aniso
When the first space is definite, the kernel's constant term has only the term
at the zero vector. The constant term of the theta integral is then a theta
integral of the smaller pair.

% before: asm-theta-agree-split-big
When the first space is large compared with $W$, the theta integral converges
absolutely and the sum over isotropic vectors unfolds directly into a cone
integral. This gives the two parts of the constant term, with their laws.

% before: asm-theta-agree-split-small
When the first space is small, the sum over isotropic vectors is compared with
a second sum, along the other rank-one direction, by two expansions of the same
theta sum. The difference is a sum over the orbits of $E^\times$ that has the
law of the zero term, and the second sum unfolds into a cone integral.

% before: asm-theta-agree-split-any
For polynomial data the two cases combine into one statement about the part
with the law of the first character. A definite first space is treated
directly; otherwise an adapted plane is chosen in the first space after
re-choosing its archimedean data and bases.

% before: asm-theta-agree-split
The constant term of a theta integral along a plane splits into a part with
the law of $\alpha$ and a part with the law of $\beta$. The comparison uses
that this splitting is unique.

% before: asm-theta-agree-ct-beta-zero
For the removal of a plane from the second space, the $\beta$-part of the
constant term of the first theta integral is a theta integral of the smaller
triple. The evaluation map computes it at the Levi elements.

% before: asm-theta-agree-ct-beta-cone
On the second space the $\beta$-part of the constant term is a theta integral
of a cone. The cone averages over the lattice stabilizer and over the ideles,
weighted by $\nu$ and a power of the absolute value.

% before: asm-theta-agree-ct-alpha-cone
For the removal of a plane from the first space, the roles are exchanged. The
$\alpha$-part of the constant term of the first theta integral is a theta
integral of a cone of the smaller triple.

% before: asm-theta-agree-ct-alpha-zero
In the same situation the constant term of the second theta integral has no
$\beta$-part. Its $\alpha$-part is a theta integral of the smaller triple
through the evaluation map.

% before: asm-theta-agree-factors
The constants of two triples related by a plane differ by a step factor and a
partial $L$-value. Both must be nonzero for nondegeneracy to pass from the
smaller triple to the larger one.

% before: asm-theta-agree-invariance-heisenberg
Theta functions are sums over rational points of Schwartz-Bruhat functions. Their
invariance under the rational points of the unitary groups comes from comparing
linear forms on the Schwartz-Bruhat space that are unchanged by the Schrödinger
operators of rational points, or of the adelic points of a rational subspace.

% before: asm-theta-agree-invariance-tate
The analysis of the theta sum uses the self-duality of the adeles. The quotient
by $F$ is compact, its characters are given by elements of $F$, and
multiplication by an element of $F^\times$ preserves Haar measure.

% before: asm-theta-agree-invariance-theta
The theta sum converges on Schwartz-Bruhat functions, depends continuously on
the archimedean component, and is not identically zero. A rational translation
permutes its terms, and a character trivial on rational points does not change
them.

% before: asm-theta-agree-invariance-poisson
Fix the finite component of a function to be the indicator of a lattice. A linear
form invariant under rational translations and characters then becomes a
tempered distribution on $F_\infty^I$ invariant under a lattice and under the
characters dual to it, and such a distribution is a multiple of the sum over the
lattice.

% before: asm-theta-agree-invariance-modulation
The characters dual to the lattice must be supplied by rational Schrödinger
operators. On the adeles whose finite part lies in $N\widehat{\mathcal O}$, each of
them agrees with a character $x\mapsto\psi(xy)$ with $y\in F$.

% before: asm-theta-agree-invariance-unique
Comparing the levels of the finite component gives the uniqueness of the theta
sum. An admissible linear form that is unchanged by every rational Schrödinger
operator is a multiple of it.

% before: asm-theta-agree-invariance-average
A second linear form is attached to a rational Lagrangian subspace: the average
of the theta sum over the adelic points of the subspace modulo its rational
points. It is invariant under those adelic points, and it inherits every
eigenvalue of the theta sum under an operator that normalizes the Schrödinger
operators through a map preserving the subspace.

% before: asm-theta-agree-invariance-character
The eigenvalues of the theta sum under rational elements of a unitary group form
a character of its rational points. A character that is trivial on the
stabilizer of a rational Lagrangian is trivial, because every rational element
is a product of elements of the stabilizer and of unipotent elements opposite to
it, on which every character is trivial.

% before: asm-theta-agree-invariance-datum
The unitary groups of a Weil pair act by operators that normalize the
Schrödinger operators and satisfy a Siegel rule on a doubled space. These
conditions make sense for any pair of matrices and any scalar in the trace form,
so the form can be rescaled until both matrices are Hermitian.

% before: asm-theta-agree-invariance-split
When the Siegel rule holds for a rational Lagrangian, the rational points act
trivially on the theta sum. By uniqueness each rational element multiplies the
theta sum by a scalar, and the Lagrangian average shows that the scalar is $1$
on the stabilizer of the Lagrangian.

% before: asm-theta-agree-invariance-doubling
For a Weil datum the Siegel rule holds on the doubled space, where the diagonal
is a rational Lagrangian. Embedding $U_{J_B}$ in the doubled group and taking
the product with a function of nonzero theta sum carries the invariance back to
$U_{J_B}$.

% before: asm-theta-agree-invariance-rational
For a Weil pair the trace form is rescaled so that both matrices become
Hermitian, and each unitary group is then treated by the doubling argument. The
result is Weil's invariance of theta series under rational points, for the
operators of a Weil pair \cite{Wei64}.

% before: asm-theta-agree-anisotropic-automorphy
Theta integrals are taken over a fundamental domain for the rational points of
$U_X$. Their behaviour under translations rests on the invariance of the theta
function under the rational points of both unitary groups.

% before: asm-theta-agree-anisotropic-invariance
Because $U_X$ is unimodular, integration over a fundamental domain turns left
invariance under $U_X(F)$ into invariance under right translations. For theta
integrals this means that the action of $U_X$ on the first factor leaves them
unchanged.

% before: asm-theta-agree-anisotropic-augmentation
The augmentation space at a finite place is spanned by the changes that the
action of $U_X$ at that place makes to a local component. By the invariance
just proved, changing one local component by an element of this space does not
change the theta integral.

% before: asm-theta-agree-anisotropic
A polynomial constant is defined by the comparison on full-rank data of the
second space. When the second Hermitian space is definite, the comparison
holds for all corresponding polynomial data.

% before: asm-theta-agree-rechoice-terms
Re-choosing the archimedean data changes the Cartan involutions, and with them
the polynomial Fock spaces. Elements of the unitary groups at the infinite
places carry one choice of involutions to the other, and the Weil
representation at such an element acts on the archimedean component alone. The
comparison also uses the subspaces of Howe's theorem, at the infinite places
and at each finite place.

% before: asm-theta-agree-rechoice-schur
Two systems of maps with the same kernels and the same equivariance differ by
a linear map on one side and by a further scalar on the other. Dixmier's lemma
supplies the scalar, since the commutant of an irreducible family on a space of
countable dimension consists of scalars.

% before: asm-theta-agree-rechoice-operators
An element of the unitary groups at the infinite places commutes with the
Schrödinger operators of finite vectors. So the Weil representation at it
factors through an operator on archimedean Schwartz functions, and this
operator is continuous because it implements a symplectic map.

% before: asm-theta-agree-rechoice-intertwine
The derivatives of the Schrödinger translates are first-order differential
operators. An archimedean operator carries the derivative along one vector to
the derivative along its image under the group element.

% before: asm-theta-agree-rechoice-fock
The polynomial Fock space is generated by Gaussians under these derivatives.
An archimedean operator at an element that moves one choice of Cartan
involutions to another carries Gaussians to Gaussians, and so carries one Fock
space into the other.

% before: asm-theta-agree-rechoice-lie
Along a line in the unitary groups the archimedean operators form a
one-parameter group implementing symplectic maps. Up to a scalar factor it is
the metaplectic representation, whose derivative is a quadratic expression in
the Schrödinger derivatives, so it preserves the polynomial Fock space.

% before: asm-theta-agree-rechoice-count
Dixmier's lemma needs countable dimension. The polynomial Fock space is the
image of a polynomial ring in finitely many variables.

% before: asm-theta-agree-rechoice-linear
The actions of lines and of the maximal compact subgroups at the infinite
places are then defined on the polynomial Fock space and are unique. The
actions of the second group are linear maps of it, as the irreducibility
arguments below require.

% before: asm-theta-agree-rechoice-howe
Howe duality at the real places and at the finite places gives, for each local
component, a unique maximal proper theta subspace. At the finite places it is
derived from the two alternatives of Theorem~\ref{th-howe-finite}.

% before: asm-theta-agree-rechoice-kernels
The kernels of matched maps are proper theta subspaces, and the irreducibility
of the targets makes them maximal. By the uniqueness just proved, any two
systems of matched maps have the same kernels.

% before: asm-theta-agree-rechoice-compare
With equal kernels and the same equivariance, the abstract lemma applies at
the infinite places and at each finite place. Off the finite set of the
character datum the vacua fix the scalar to be $1$.

% before: asm-theta-agree-rechoice-cartan
Two Cartan involutions of a Hermitian matrix of given signature have frames of
the same type. The frames give a unitary matrix at each infinite place carrying
one involution to the other.

% before: asm-theta-agree-rechoice-move
For the comparison, the element moving the Cartan involutions of the Hermitian
space must be chosen together with the one for the skew-Hermitian space. Its
archimedean operators then carry one polynomial Fock space onto the other.

% before: asm-theta-agree-rechoice-transport
Composing the real maps of the original data with these archimedean operators
gives real maps for the re-chosen data. The conditions transfer because the
operators intertwine the actions of conjugate lines and conjugate compact
elements.

% before: asm-theta-agree-rechoice
The towers of Lemmas~\ref{asm-theta-agree-tower-alpha}
and~\ref{asm-theta-agree-tower-beta} re-choose the archimedean data and the
bases. A nonzero polynomial constant for the re-chosen data must give one for
the original data.

% before: asm-theta-agree-nonvanishing-schwartz
The towers re-choose the matrices $\theta_w$, and with them the Gaussians and
the polynomial data. Over a definite space the theta integral is continuous in
the archimedean component, and polynomial Fock vectors are dense. So its
vanishing on polynomial data for one choice gives its vanishing on all
Schwartz-Bruhat functions, which do not depend on the choice.

% before: asm-theta-agree-nonvanishing-full-rank
When the second space is definite, its theta integral is an integral over a
compact quotient, so it converges and is linear on all Schwartz-Bruhat
functions. The coinvariants at the place $v_1$ then let a polynomial datum be
replaced by one of full rank at $v_1$ with the same theta integral.

% before: asm-theta-agree-nonvanishing-doubling
When both spaces are definite, nonvanishing comes from the doubling zeta
integral. It integrates a matrix coefficient of the Weil representation
against a power of a height function on the doubled group, and its factors at
the places outside a finite set are integrals of spherical functions over the
local unitary groups.

% before: asm-theta-agree-doubling-iwasawa
The doubling zeta integral and the Eisenstein series of the doubled group are
built from a height function, which needs every element of the doubled group to
factor through the Siegel parabolic subgroup and a compact subgroup. At each
place this is the Iwasawa decomposition, for the stabilizer of the doubled
lattice at the finite places and for the centralizer of the doubled Cartan
involution at the infinite places.

% before: asm-theta-agree-doubling-stabilizer
The local decompositions assemble to one for the lattice stabilizer of the
doubled space, which is compact. Compactness makes the height of a
factorization independent of the choice, and continuous.

% before: asm-theta-agree-doubling-count
The zeta integrals converge for large real part because the height on the
embedded unitary group is small away from sets of finite measure. The estimate
below bounds the measure of the set where the height is at least $2^{-j}$ by a
geometric sequence in $j$.

% before: asm-theta-agree-doubling-integrable
Cramer's rule expresses an element of the unitary group, scaled by the norm of
the determinant of its Iwasawa component, through the entries of an element of
the compact lattice stabilizer. With the measure estimate this makes high
powers of the height integrable on the unitary group.

% before: asm-theta-agree-doubling-frame
The continuation of Siegel Eisenstein series and the Siegel-Weil formula are
stated for the standard compact subgroups of frames. The lattice stabilizer is
conjugate to such a subgroup by an element of the Siegel parabolic subgroup,
which changes the height function by a constant factor.

% before: asm-theta-agree-doubling-continuation
Transported through this conjugation, the continuation of Siegel Eisenstein
series holds for sections of the lattice stabilizer. The weak holomorphy in the
published statement becomes, by uniform boundedness, a bound of moderate growth
that is uniform near each point.

% before: asm-theta-agree-doubling-datum
The Siegel-Weil formula is applied to the box product of two polynomial data on
the doubled space. The doubling datum built from the Weil representation of a
pair has a theta sum that splits as a product of two theta functions; this
seesaw identity turns the value of the Eisenstein series into an average of
products of theta functions.

% before: asm-theta-agree-doubling-level-gaussian
The Gaussian relations say that a complex multiple of one first-order operator
of the Schrödinger representation agrees with another on the function. When
the relations hold for every direction, the logarithmic derivative of a nonzero
solution is linear, and the solution is the exponential of a quadratic form
with negative definite real part. Two such Gaussians cannot satisfy relations
of opposite signs in a nonzero direction.

% before: asm-theta-agree-doubling-level
The section given by the box product must have right translates under the
lattice stabilizer that span a finite-dimensional space. Two facts about single
functions give this: a Schwartz-Bruhat function is fixed by the Schrödinger
operators of a stable set of finite vectors, and a solution of the Gaussian
relations is a nowhere vanishing Gaussian.

% before: asm-theta-agree-doubling-section
For polynomial data the box product is finite under the lattice stabilizer,
since products of Gaussians span finite-dimensional stable spaces and
derivatives preserve finiteness. So the box product defines a section of the
degree at which the Siegel-Weil formula is evaluated.

% before: asm-theta-agree-doubling-siegel-weil
The Siegel-Weil formula for a standard compact subgroup transfers to the lattice
stabilizer through the conjugating element of the Siegel parabolic subgroup.
The transfer uses that this subgroup commutes with the action of the other
unitary group in the doubling datum.

% before: asm-theta-agree-doubling-value
The continuation of the Eisenstein series and the Siegel-Weil value combine into
one continued function that is holomorphic at the Siegel-Weil point. The
maximum principle on a small circle gives the bounds and the continuity there.

% before: asm-theta-agree-doubling-kernel
When the unitary group is definite, every coset of the Siegel parabolic subgroup
in the rational doubled group meets the image of the rational unitary group in
exactly one element. So the integral of the Eisenstein series against a function
on both copies of the unitary group unfolds to an integral of the section
against a matrix coefficient.

% before: asm-theta-agree-nonvanishing-zeta-local-dichotomy
At a split place the local component of the polynomial datum will be the
indicator function of a small neighbourhood of a matrix of full rank. An
element of the general linear group either moves that neighbourhood into
itself, when it is close to $1$, or is detected by a character that is trivial
on the neighbourhood and on its image but not on the displacement of its
centre.

% before: asm-theta-agree-nonvanishing-zeta-local-vacuum-skew
At an unramified place an element of the lattice stabilizer of the doubled
space is to be split into an involution and elements of the Siegel parabolic
subgroup, which needs a unit built from its columns. Over the integers of the
place, a correction by a skew matrix makes a determinant a unit, as one can
raise the rank of its residue one step at a time.

% before: asm-theta-agree-nonvanishing-zeta-local-vacuum-columns
The columns of an element of the doubled group span an isotropic subspace, and
this makes the hypothesis of the previous lemma hold over the integers of a
finite place. The resulting correction gives the unit with integral inverse
from which the decomposition is built.

% before: asm-theta-agree-nonvanishing-zeta-local-vacuum-decomposition
With the unit in hand, an element of the doubled group factors as an involution
conjugate of a unipotent element times an element of the Siegel parabolic
subgroup. The Siegel rule then computes the action of each factor on the
doubled vacuum.

% before: asm-theta-agree-nonvanishing-zeta-local
The polynomial datum of the last part of the next lemma is chosen place by
place. At a place of $P$ its component either is fixed by an integral element
or makes the matrix coefficient vanish there, and at an unramified place the
doubled vacuum is fixed by the lattice stabilizer, which makes the zeta
integrand spherical.

% before: asm-theta-agree-nonvanishing-zeta
Four properties of the doubling zeta integral are used. A height function
exists, and so does a polynomial Fock vector fixed by the maximal compact
subgroup at infinity. The zeta integral continues to the point where the
Siegel-Weil formula holds and vanishes there with the theta integral, and a
polynomial datum can be chosen whose integral over the elements integral
outside $T$ factors at those places and is nonzero at that point.

% before: asm-theta-agree-nonvanishing-euler
The zeta integral of a function that factors at the places outside $T$ is a
product of local integrals and an integral over the elements integral outside
$T$. This holds for every integrable function with such local factors, once
the local Haar measures give the integral elements volume one.

% before: asm-theta-agree-nonvanishing-normalizer
The local factors $d_v(s)$ normalize the unramified integrals. Their product
over the places outside $T$ converges on the half-plane
$\operatorname{Re}s>0$ to a nonzero holomorphic function, by comparison with
the Dedekind zeta function of $F$.

% before: asm-theta-agree-nonvanishing-unramified
At a place where the character is unramified, the local zeta integral of a
spherical function is a product of local Hecke $L$-factors divided by $d_v$.
This is the unramified computation of the doubling method for the trivial
representation.

% before: asm-theta-agree-nonvanishing-pair
For a pair of definite spaces, the Euler product identifies the zeta integral
with a product of Hecke $L$-functions times an entire function. Continuing
this identity to the point where the Siegel-Weil formula holds shows that the
zeta integral is nonzero there whenever the $L$-values are, and so the theta
integral is nonzero.

% before: asm-theta-agree-nonvanishing-hecke
The $L$-values that appear lie at integers where the $L$-function of $\nu$ has
no zero. The infinity type of $\nu$ is larger than $a+a'$ at every place,
which keeps the gamma factors away from their poles there.

% before: asm-theta-agree-nonvanishing-bottom
When both spaces of a triple are definite, its two Weil pairs satisfy the
hypotheses of Lemma~\ref{asm-theta-agree-nonvanishing-pair}, with the
characters $\nu^{-1}$ and $\nu$. This starts the inductions on $p$ and on $q$.

% before: asm-theta-agree-nonvanishing
When one Hermitian space is definite, its theta integrals are integrals over a
compact quotient. Their nonvanishing on polynomial data comes from the
Siegel-Weil formula and the doubling zeta integral.

% before: asm-theta-agree-definite-volume
When both Hermitian spaces of a triple are definite, all three unitary groups
are compact modulo their rational points. The argument for this case
integrates over fundamental domains and translates on the right, so it needs
finite volume and unimodularity.

% before: asm-theta-agree-definite-products
The case of definite spaces compares theta integrals through their Petersson
products over a fundamental domain. Integrating a function on $U_W$ against
theta functions carries it back to functions on the other unitary groups.

% before: asm-theta-agree-definite-compact
For a definite space the Witt tori of reduction theory have size zero, so one
compact set meets every orbit of the rational points. Invariant continuous
functions are then bounded and rapidly decreasing, and their Petersson norms
detect whether they vanish.

% before: asm-theta-agree-definite-continuous
Petersson products of theta integrals need these integrals to be continuous on
$U_W$. The domination of theta functions by an integrable function, with a
polynomial factor in the height, gives this by dominated convergence.

% before: asm-theta-agree-definite-dist-tools
The local invariance at a split place comes down to linear forms on locally
constant functions with compact support on spaces of matrices over a $p$-adic
field. Such a form can be shown to vanish piece by piece: near each point of a
support, on products of functions, and, for a form invariant under translation,
at a single bump function.

% before: asm-theta-agree-definite-dist-group
On the open set of invertible matrices, equivariance under right translation is
to be turned into equivariance under left translation. An involution of the
group that inverts the given element exchanges two compact open subgroups of
the same index, and the stabilizer of a function decomposes it into translates
of their indicator functions.

% before: asm-theta-agree-definite-dist-degenerate
Under a condition on the character at one element, a linear form on functions
on matrices that is equivariant under right translation, and vanishes on
functions supported away from the degenerate matrices, vanishes. The proof
inducts on the number of rows, through a chart near the matrices with a nonzero
corner entry and the invariance under unipotent column operations.

% before: asm-theta-agree-definite-dist
For square matrices the degenerate ones are the singular ones, so an
equivariant form is controlled by its restriction to the invertible matrices,
where the involution $x\mapsto(x^{\mathsf T})^{-1}$ inverts every symmetric
element. Symmetric matrices generate the general linear group, which gives
equivariance under left translation by every element.

% before: asm-theta-agree-definite-dist-bilinear
The local invariance concerns a pairing between functions on two spaces of
matrices whose rows together fill a square matrix. Through tensor products of
Schwartz-Bruhat functions such a pairing is a linear form on functions on square
matrices, to which the equivariance on square matrices applies.

% before: asm-theta-agree-definite-hecke
At a place where the rational prime splits completely, the local invariance
needs the value of the character $\nu$ at a uniformizer to differ from $1$. As
$\nu$ has a nonzero infinity type, its value at a prime of degree one is not
even a root of unity: a power of the prime is principal, and the Galois
conjugates of a generator other than itself are units at the prime.

% before: asm-theta-agree-definite-inputs-tensor
The correspondence of two functions is an identity of tensors over the places
of a finite set, while theta integrals are linear in the functions themselves.
Passing between the two needs that a function determines its tensor once the
components outside the set are fixed, and that a multilinear form killed by the
kernels of the local maps factors through the tensor product of their images.

% before: asm-theta-agree-definite-inputs-constancy-hecke
The constancy criterion for functions on unitary groups reduces, through the
norm-one ideles, to a statement about Hecke characters. A unitary Hecke
character that is trivial at almost all places of prime absolute norm has an
$L$-function with a pole at $s=1$ unless it is trivial, and Hecke's theorem
excludes the pole.

% before: asm-theta-agree-definite-inputs-constancy-ideles
By duality for the compact group of idele classes of norm one, the statement
about characters becomes one about closed subgroups of the ideles. A closed
subgroup that contains the principal ideles and the local multiplicative groups
at almost all places of prime norm is everything.

% before: asm-theta-agree-definite-inputs-constancy-norm
The unitary group of rank one is the group of ideles of norm one for $E/F$.
Hilbert's Theorem 90 writes each of its elements as a quotient $z/c(z)$, which
carries the statement on closed subgroups of the ideles of $E$ over to it. The
places of $E$ of prime absolute norm lie over primes that split completely.

% before: asm-theta-agree-definite-inputs-constancy-cayley
To pass from rational unitary matrices to the unitary group at one place, the
rational points must be dense there. The Cayley transform turns unitary
matrices into skew matrices, which form a linear space in which the rational
points are dense.

% before: asm-theta-agree-definite-inputs-constancy-lines
Every idempotent matrix that is self-adjoint for the Hermitian form maps the
norm-one ideles into the unitary group. A closed subgroup that contains
all these embedded groups contains every rational unitary matrix, by
Dieudonné's generation by quasi-reflections, and then, by density at each
place, the whole unitary group.

% before: asm-theta-agree-definite-inputs-pairs-density
The space of pairs is described through tensors over a finite set of places,
on whose factors the local unitary groups act irreducibly. With the countable
dimension of these factors, the Jacobson density theorem lets a subspace stable
under the local operators be detected on pure tensors.

% before: asm-theta-agree-definite-inputs-pairs-finite
For a definite space $W$ every line of $U_W$ lies in the maximal compact
subgroup at infinity, and the derivative of an orbit is a limit of difference
quotients. These stay in a fixed finite-dimensional space, since the compact
group moves a polynomial datum within one.

% before: asm-theta-agree-definite-inputs-pairs-operators
The actions of $U_W$ on the two sides of a correspondence pass to the local
quotients $\Pi_v$ and $\Pi$. The matching conditions make the induced operators
the same for both sides, which lets them act on the tensors that define
correspondence.

% before: asm-theta-agree-definite-inputs-pairs-fill
A subspace of polynomial data stable under $U_W$, and whose tensors are not all
zero, reaches every tensor. Irreducibility at the finite places gives a pure
tensor with a nonzero factor at infinity, and irreducibility at infinity then
gives every factor there.

% before: asm-theta-agree-definite-inputs-growth-lattice
A theta function is a sum over the rational points of a function that decays
at infinity and is supported on a lattice with some denominator. Comparing with
the points of the integral lattice bounds the sum by the decay constant times a
power of the denominator.

% before: asm-theta-agree-definite-inputs-growth-inputs-denominator
The height of a matrix bounds the absolute values of its entries at each
finite place, and at almost all places this bound is $1$. A single integer
that clears all these denominators at once is a product of powers of the norms
of the places involved, and it is no larger than the product of the local
bounds.

% before: asm-theta-agree-definite-inputs-growth-inputs-entries
The map $s_{g,h}$ is built from the entries of $g^{-1}$ and $h^{-1}$, which lie
in $\mathbf A_E$. Writing these entries in a fixed basis of $\mathbf A_E$ over
$\mathbf A$ gives adelic coordinates, and the height bounds both their
archimedean size and their denominators.

% before: asm-theta-agree-definite-inputs-growth-inputs-bilinear
When the entries of two matrices have coordinates of bounded archimedean size
and denominator, so does every fixed bilinear expression in them. The map
$w\mapsto e(Me^{-1}(w)N^{\mathsf T})$ is such an expression, so its
denominators and its archimedean size are bounded by the products of the
bounds for $M$ and $N$.

% before: asm-theta-agree-definite-inputs-growth-inputs-derivatives
The rapid decay of $\omega(g,h)\Phi$ comes from a Sobolev inequality on the
adeles. Its derivatives are those of the Schrödinger operators along
archimedean directions, and it applies to functions invariant under
translation by a multiple of the integral finite adeles.

% before: asm-theta-agree-definite-inputs-growth-inputs-derivatives-bounds
Along an archimedean direction, the derivative of a Schwartz-Bruhat function
is given by a first-order differential operator on its archimedean factor, so
it is again Schwartz-Bruhat. Its $L^2$-norms grow polynomially in the size of
the directions, and its difference quotients are dominated by a square
integrable function, which allows limits under the integral.

% before: asm-theta-agree-definite-inputs-growth-inputs-transport
An operator that intertwines the Schrödinger operators through a linear map
$s$ and preserves $L^2$-norms carries the derivative along a direction to the
derivative along its image under $s$. When $s$ enlarges archimedean directions
by at most a factor $c$, iterated derivatives of the transformed function are
transformed derivatives along directions at most $c$ times larger.

% before: asm-theta-agree-definite-inputs-growth-inputs-sobolev-approximation
The bound for periodic functions below compares the measure of the integral
finite adeles with that of their multiples by $m$. The comparison rests on
writing a finite adele as an element of $F$ plus $m$ times an integral one,
which follows from the Chinese remainder theorem in $\mathcal O_F$ and the
density of $F$ in each completion.

% before: asm-theta-agree-definite-inputs-growth-inputs-sobolev
The values of a Schwartz-Bruhat function invariant under translation by $m$
times the integral finite adeles are bounded by the $L^2$-norms of finitely
many of its derivatives. The constant grows polynomially in $m$, as the
integral finite adeles are covered by polynomially many translates of their
multiples by $m$.

% before: asm-theta-agree-definite-inputs-growth-inputs-trace
The derivative along a direction $(0,b)$ multiplies a function by a linear
form in the archimedean variable. To bound a power of the size of that
variable by such a factor, one needs a direction $b$ of bounded size along
which the trace pairing is comparable to the size of the variable.

% before: asm-theta-agree-definite-inputs-growth-inputs
To bound a theta function polynomially in the heights, both the decay at
infinity and the denominator of the support must be controlled by the heights.
The action of $(g,h)$ on the symplectic space changes denominators and
archimedean sizes by at most the heights, and the Weil operators transport
this to the functions.

% before: asm-theta-agree-definite-inputs
The comparison for definite spaces rests on local properties of the Weil
representation at split places, on the stability and irreducibility of the
space of pairs, and on the doubling method with the Siegel-Weil formula. A
constancy criterion for functions on unitary groups of rank at least two
completes them.

% before: asm-theta-agree-definite-orbits
The space of pairs carries a representation of $U_W$, and its invariant
Hermitian forms are to be shown proportional. In a space of countable
dimension, the Baire category theorem confines the orbit of a vector under a
compact set to a finite-dimensional subspace.

% before: asm-theta-agree-definite-eigen
For a group acting unitarily with finite-dimensional orbit spans, an invariant
form has an eigenvector relative to the inner product. The argument reduces to
a finite-dimensional subspace built from equivariant maps.

% before: asm-theta-agree-definite-schur
An irreducibility condition upgrades one eigenvector to proportionality on the
whole space. This is a form of Schur's lemma for representations realised on
continuous functions on a locally compact group.

% before: asm-theta-agree-definite-open
The form of Schur's lemma above needs a compact open subgroup. For a definite
space the archimedean unitary groups are compact, and at the finite places the
matrices congruent to $1$ modulo an integral lattice give the rest.

% before: asm-theta-agree-definite-line
With $p=q=0$, the Petersson products of the functions $\mathcal I_{s,t}(y)$ are
invariant forms on the space of pairs. The space is irreducible modulo the
kernel of $\mathcal I_{0,1}$, so all these forms are multiples of one of them.

% before: asm-theta-agree-definite-norms
The doubling method expresses the Petersson norm of a theta lift to $U_{V'}$
through an Eisenstein series on the doubled group of $W$. Its unfolded form
sees the function only through its Petersson products with its translates, so
a proportionality of products passes to the norms of theta lifts.

% before: asm-theta-agree-definite-constant
A rapidly decreasing function on $U_W$ lifts to functions on $U_V$ and
$U_{V'}$ that are invariant under the local unitary groups at the clean
places. As the ranks are at least two, such lifts are constant.

% before: asm-theta-agree-definite
When both Hermitian spaces are definite, the induction starts. A nonzero
polynomial constant comes from the two published results on Siegel
Eisenstein series.

% before: asm-theta-agree-local-cones
The proofs of Lemmas~\ref{asm-theta-agree-local-alpha}
and~\ref{asm-theta-agree-local-beta} compare the cone with an evaluation map
one place at a time. The cone integral splits into integrals over local
pieces of the lattice stabilizer and of the ideles, and these define local
cones at the finite places and an operator on Schwartz functions at the
infinite places.

% before: asm-theta-agree-local-measures
The local cones need Haar measures on the local pieces. They exist with the
normalizations below because the lattice stabilizer is compact.

% before: asm-theta-agree-local-real-unique-diagonal
Two Weil representations with the same Schrödinger operators differ, on each
element of the doubled group, by a scalar. To compare these scalars on the
Siegel parabolic subgroup, one needs a linear form that is invariant under the
diagonal Lagrangian and nonzero.

% before: asm-theta-agree-local-real-unique-cayley
The element $\operatorname{diag}(g,1)$ of the doubled group is then reduced to
the Siegel parabolic subgroup. When $g+1$ is invertible, the Cayley transform
of $g$ writes it as a conjugate of an element of that subgroup times another
one.

% before: asm-theta-agree-local-real-unique
Two matched quadruples for the same triple datum may have different
characters. On the unitary group of the first space at the infinite places
the Weil representation depends only on the infinity types, and these are the
same for all character data.

% before: asm-theta-agree-local-cone-unique
As a consequence, the archimedean cone does not depend on the quadruple. This
lets the proofs choose one constant for all quadruples at once.

% before: asm-theta-agree-local-factor-sums
The cones are iterated integrals of functions that take only finitely many
values. Such an integral is a finite linear combination of those values, so it
stays in every subspace that contains them and commutes with linear forms.

% before: asm-theta-agree-local-factor-evaluation
At a finite place a Schrödinger operator acts through an explicit local
formula, and an evaluation map acts through its components. An identity
between global operators that holds after evaluation therefore descends to the
components at the place.

% before: asm-theta-agree-local-factor-setting
Both kinds of configuration give the same local data: two pairs of lattice
spaces with adapted planes, and an evaluation map that is invariant under
$\mathcal J_\alpha$ or under $\mathcal J_\beta$. Only these data enter the local
cones.

% before: asm-theta-agree-local-factor-directions
The map $u$ moves along one rank-one matrix $V_0$, and the evaluation map is
invariant under translation by multiples of a second matrix $V_1$ that pairs
with it. Comparing the two through the lattice $\Gamma_v$ shows that a nonzero
local value puts a multiple of $1_vu(l)$ into $\Gamma_v$.

% before: asm-theta-agree-local-ideal-lines
The vectors of an adapted plane meet the lattice in lines: the multiples of $x$
in the lattice are given by the ideal of the plane, and those of $y$ by its dual.
At the places of $P$ these ideals are prime to the place, as the frame of the
lattice there has $x$ and $y$ among its columns.

% before: asm-theta-agree-local-ideal-lattice
For two such lines the lattice group $\Lambda_v$ meets the rank-one matrices
$c\,xy^{\mathsf T}$ in those whose coefficient lies in the product of the two
ideals at $v$. This turns the lattice condition on $1_vu(l)$ into a condition on
$l$.

% before: asm-theta-agree-local-ideal-values
The value of $\nu$ on the resulting ideal is computed from three facts. The
character $\nu$ inverts under conjugation, the same holds for its values on
ideals prime to $P$, and it is trivial on the codifferent, which comes from $F$.

% before: asm-theta-agree-local-factor-ideal
The lattice condition on $1_vu(l)$ is a condition on $l$ alone. It cuts out a
fractional ideal, whose value under $\nu$ is a product of powers of the values
on the ideals of the two planes.

% before: asm-theta-agree-local-ideles-parts
The idele group is cut into its components at sets of places by idempotents of
the adeles, as the lattice stabilizer is. When each Haar measure gives the units
mass $1$, the product decompositions carry the measures to product measures with
constant $1$.

% before: asm-theta-agree-local-ideles-tate
At a place where $\nu$ is unramified, the local integral of $\nu(l)|l|^z$ over the
ideles of order at least $m$ is a geometric series. Its value is the local factor
that enters the partial $L$-value.

% before: asm-theta-agree-local-ideles-tail
Outside a finite set of places the integral runs over infinitely many components.
It is the limit of the integrals over the ideles that are units outside larger and
larger finite sets, and these limits converge because $z\geq2$.

% before: asm-theta-agree-local-factor-ideles
The integral over $\mathbf A_E^\times$ in the cone factors over the places. Away
from the ramified places the local factors are integrals of $\nu(l)|l|^z$ over
the local ideal, and their product is the partial $L$-value.

% before: asm-theta-agree-local-factor-tate
At a finite place the inner integral of the local cone runs over $A_v$. Its
integrand takes finitely many values and its weight is integrable on the support,
so the integral is a finite sum.

% before: asm-theta-agree-local-factor-finite
So the local cone at a finite place is a finite sum of Schwartz-Bruhat
functions and is linear in its argument. At an unramified place the evaluation
of a translated vacuum is the vacuum or $0$, according to the lattice
condition.

% before: asm-theta-agree-local-stabilizer-parts
An element of the lattice stabilizer is cut into its components at the places by
multiplying with idempotents of the adeles. The components are again in the
stabilizer, they commute when the idempotents are orthogonal, and their products
over larger and larger sets of places converge to the element.

% before: asm-theta-agree-local-factor-stabilizer
The integral over the lattice stabilizer factors in the same way, into
integrals over its components at the places, up to a constant. This needs the
components of an element of the stabilizer to lie in the stabilizer again.

% before: asm-theta-agree-local-vacuum-lattice
At a finite place the vacuum is described through the group $\Lambda_v$ built
from the lattices. An element of the lattice stabilizer at $v$ preserves the
local lattice $L_v$, and so it preserves $\Lambda_v$; this is what lets it act on
the vacuum by a scalar.

% before: asm-theta-agree-local-vacuum-residue
At an inert place the decomposition below needs a matrix over the residue field
to be invertible. A power of $2+N+N'$ is idempotent, and a suitable exponent
makes the matrix invertible.

% before: asm-theta-agree-local-vacuum-units
The scalar by which an element $k$ of the stabilizer acts on the vacuum is
computed in the doubled group. There $\operatorname{diag}(k,1)$ is moved into the
Siegel parabolic, which needs matrices $b$ and $b'$, adjoint to each other with
sum $-1$, such that $1+(1-k)b$ is invertible and $b$ and the inverse preserve the
local lattice.

% before: asm-theta-agree-local-vacuum-decomposition
With these matrices, $\operatorname{diag}(k,1)$ is the product of an involution, a
unipotent element, the involution again and an element of the Siegel parabolic.
The Levi part of the last factor is $1+b(1-k)$, which preserves the local
lattice together with its inverse.

% before: asm-theta-agree-local-vacuum-doubled
Each factor of this decomposition acts on the doubled vacuum by a scalar. As for
the vacuum itself, this follows from the uniqueness up to scalars of an
eigenvector of the doubled lattice group with a given character.

% before: asm-theta-agree-local-factor-vacuum
At an unramified place the local cone of the vacuum must be a multiple of the
vacuum. For this the component at the place of the lattice stabilizer must fix
the vacuum when the characters are unramified there.

% before: asm-theta-agree-local-archimedean-orthogonal
The archimedean cone shifts by a matrix built from one vector of each plane. Such
matrices, and every matrix after a correction from the smaller spaces, are
orthogonal to the matrices coming from the smaller spaces. This splits the
Heisenberg directions into those the evaluation map transports and those it
does not see.

% before: asm-theta-agree-local-archimedean-gaussians
The polynomial Fock space is generated by its Gaussians. In the frame
coordinates every Gaussian is a multiple of the image of the constant
polynomial, so the Gaussians span at most a line.

% before: asm-theta-agree-local-archimedean-slices
The archimedean cone integrates the slices of a polynomial datum over a real
group. Its values stay polynomial if all the slices of finitely many polynomial
data lie in one finite-dimensional space of polynomial data for the smaller
model. Induction over the generators of the Fock space gives this, as the slices
carry Gaussians to Gaussians and commute with the derivations along the smaller
spaces.

% before: asm-theta-agree-local-factor-archimedean
At the infinite places the cone is an integral over real groups. It is a
continuous map of Schwartz spaces that preserves the polynomial Fock spaces, on
which the real maps of the matched quadruples are defined.

% before: asm-theta-agree-local-pure-schrodinger
The integrand of the cone applies a Schrödinger operator to the image of a pure
tensor. Since the character $\psi$ is the product of its local components, a
Schrödinger operator maps pure tensors to pure tensors, and almost all of the
local factors stay the indicator functions of the integral vectors.

% before: asm-theta-agree-local-pure-partial
The operator $\omega(k,1)$ is handled through the components $k^{(v)}$ of $k$.
Each component acts only at its place, by the vacuum at the unramified places,
and continuity in $k$ passes from the finite partial products to $k$ itself.

% before: asm-theta-agree-local-factor-pure
On a pure tensor the integrand of the cone is again a pure tensor. Its factors
are the archimedean and local integrands, and the factors at the unramified
places impose the lattice condition.

% before: asm-theta-agree-local-factor-span
The linearity of the cone is first proved on pure tensors. These span the
Schwartz-Bruhat functions, even with the factors fixed to be the vacua at almost
all places.

% before: asm-theta-agree-local-factor
The cone of a pure tensor factors into the archimedean cone and the local
cones. At the places outside $T$ the local cone multiplies the vacuum by a
number, and the product of the remaining unramified factors is the partial
$L$-value.

% before: asm-theta-agree-local-ramified-values-integrals
At a ramified place the test functions are built from the conductor of the
local component of $\nu$, and their cones reduce to products of integrals in one
variable. The Gauss integral among them produces the epsilon factor, and the
integral over the local ideles reduces to its value at a power of the
uniformizer.

% before: asm-theta-agree-local-ramified-values-zeta
The two test functions at a ramified place are linked through the zeta
condition of the finite maps. For this their zeta functions must be equal, with
no further constant, and must not vanish.

% before: asm-theta-agree-local-ramified-values-coordinates
The local computations take place in the matrix coordinates of the linear models
of the zeta condition. In these coordinates a Schrödinger operator at the place
translates the argument and multiplies by a character.

% before: asm-theta-agree-local-ramified-values-linear
The linear models of the zeta condition have this form. Splitting an element at
a split place into its two components reduces a general Schrödinger operator to
the two cases in the definition of a linear model.

% before: asm-theta-agree-local-ramified-values-evaluation
The evaluation maps of a configuration lower the ranks of the spaces. In
coordinates their components at a split place integrate over the entries in the
rows and columns of the removed planes, as the invariance under the elements of
$\mathcal J$ forces.

% before: asm-theta-agree-local-ramified-values-dictionary
The integral over the lattice stabilizer in the cone is trivial on a test
function that the stabilizer fixes. To see that it fixes the standard functions,
the dictionary at a split place is needed with its matrices named, so that the
stabilizer acts through integral matrices.

% before: asm-theta-agree-local-ramified-values-standard
The test functions at a ramified place are the standard functions of the
conductor of $\nu_v$, read through the linear models of the zeta condition.
They have equal nonzero images in the local space and are fixed by the lattice
stabilizers.

% before: asm-theta-agree-local-ramified-values-cone
With these test functions the local cone can be computed outright. The integral
over the lattice stabilizer is trivial, the Schrödinger operator and the
evaluation map turn into an integral of a standard function, and the integral
over the local ideles gives a Gauss integral.

% before: asm-theta-agree-local-ramified-values
At the places of $T$ the local cone is compared with the local evaluation
map on suitable test functions. The ratio is the epsilon factor of the local
component of $\nu$, which enters the step factor.

% before: asm-theta-agree-local-bump-schwartz
The archimedean cones are integrals of families of Schwartz functions. Handling
them requires the Fourier inversion formula, weak integrals of such families
and the uniform boundedness principle on Schwartz spaces.

% before: asm-theta-agree-local-bump-positivity
The nonvanishing of an archimedean cone comes from a positivity argument on a
compact group times a complex space. A function that is positive on a fibre of
a map is detected by a bump function on the target, and the bump is a
superposition of characters. In the second form of the argument the characters
act as translations, and a local functional equation at the complex places
relates the two integrals.

% before: asm-theta-agree-local-bump-levi
On the integrand of the cone, an element that moves the first vector of the
plane to a multiple of itself acts through a character. Such an element is a
unipotent element times a Levi element, and the evaluation map transforms Levi
elements by the character.

% before: asm-theta-agree-local-bump-dense
A test vector at the infinite places has to be found for every nonzero
invariant form. The evaluation map intertwines the Schrödinger operators of the
two models, and the translates and modulations of a nonzero Schwartz function
span a dense subspace.

% before: asm-theta-agree-local-bump-slices-heisenberg
At the infinite places a Schrödinger operator is a translation followed by
multiplication by a phase. Its seminorms grow polynomially in the translation,
and it moves a function by an amount proportional to the translation. If a
family of maps changes sign under a translation of size inversely proportional
to a weight, these two bounds give decay in that weight.

% before: asm-theta-agree-local-bump-slices-adelic
An adelic Schrödinger operator is the product of its archimedean and finite
parts, and two of them commute up to a character. The finite part is locally
constant in the translation, a character that is trivial on a neighbourhood of
zero forces the dual variable into a compact set, and the trace form of
$F_\infty$ produces a translation that changes the sign of a character.

% before: asm-theta-agree-local-bump-slices
The integrand of an archimedean cone is a Schrödinger operator along a line,
followed by the evaluation map. Along directions dual to directions that the
evaluation map ignores, this composite splits into archimedean and finite
parts, and it decreases rapidly in the archimedean coordinates.

% before: asm-theta-agree-local-bump-cone
With this decomposition, the integrand of an archimedean cone is continuous and
rapidly decreasing on the compact group times the line. So the cone, tested
against a continuous linear form, is the integral of the tested integrand.

% before: asm-theta-agree-local-bump
To produce a test vector at the infinite places, the archimedean cone must
not vanish against an invariant linear form. Such a form arises from a
nonzero theta integral of the smaller quadruple.

% before: asm-theta-agree-local-augmentation-point
Part (1) of the lemma on augmentations is an invariance under the unitary group
of $V'$ at $v_1$. It comes from a law for the integrand of the cone under the
parabolic subgroup of $x_H$, which follows from the formula for the evaluation
map at Levi elements.

% before: asm-theta-agree-local-augmentation-slice
Part (2) concerns full rank at $v_1$. If a pure tensor with vacua away from
$v_1$ is a full-rank datum, its component at $v_1$ lies in the span of the
functions of full rank.

% before: asm-theta-agree-local-augmentation-subspace
Full rank is defined through some linear model. As every linear model of the
second kind detects it, the functions of full rank form a subspace.

% before: asm-theta-agree-local-invariance-parabolic
The law of the integrand at $v_1$ carries the factor $|t|^{n_1+1}$. In
coordinates adapted to the plane, the parabolic subgroup of $x_H$ is the
unipotent group times the Levi elements, and conjugation by a Levi element
scales the unipotent coordinates by $c(t)$ and $N_{E/F}(t)$. So it multiplies the
Haar measure of the unipotent group by $|t|^{n_1+1}$.

% before: asm-theta-agree-local-invariance-height
To move a function with this law from a place to the whole adelic group, the
proof needs a positive function with the same law that is invariant under the
compact group. The Iwasawa decomposition provides it, as the multipliers of the
elements of the compact group have absolute value $1$.

% before: asm-theta-agree-local-invariance-integral
By Lemma~\ref{asm-theta-agree-local-invariance-parabolic} the factor
$|t|^{n_1+1}$ is the modular quasicharacter of the parabolic subgroup. For functions that transform by it, the integral over a
compact group complementary to the parabolic subgroup is invariant under right
translation, by a lemma of Bump.

% before: asm-theta-agree-local-augmentation-invariance
The law for the integrand is the modulus law of the parabolic subgroup of
$x_H$ at $v_1$. For functions with this law the integral over the component of
the lattice stabilizer is invariant under right translation.

% before: asm-theta-agree-local-augmentation
For the second space, full-rank data must be kept, and this involves the
place $v_1$. There the local cone respects augmentations and preserves full
rank.

% before: asm-theta-agree-local-equivariant
The cone and the evaluation map induce maps between the spaces of the real
and finite maps. These maps intertwine the actions of Levi elements and of
the unipotent group of the plane $H'$, in the sense below.

% before: asm-theta-agree-local-mult-unramified-iwasawa
At a place outside $T$, an equivariant map is determined by its values on the
images of the vacuum under the group $U_W^{(v)}$. To compute these, an element
of $U_W^{(v)}$ is written as an element that preserves the line of an isotropic
vector times an element of the stabilizer of the lattice at $v$.

% before: asm-theta-agree-local-mult-unramified-levi
Equivariance at a finite place is defined through products of a unipotent
element and a Levi element of the plane $H'$, both at $v$. An element of
$U_W^{(v)}$ that preserves the line of the first vector of $H'$ has this
form.

% before: asm-theta-agree-local-mult-unramified-vacuum
The factor from the stabilizer of the lattice multiplies the vacuum by a
scalar. The reason is that the vacuum is the only eigenvector, up to scalars,
of the Schrödinger operators of the local lattice, and the stabilizer
preserves that lattice.

% before: asm-theta-agree-local-mult-unramified
Equivariant maps are unique up to a scalar, by local multiplicity one. At an
unramified place the scalar is read off on the vacuum.

% before: asm-theta-agree-local-mult-ramified-situation
At a ramified place the multiplicity one of equivariant maps reduces to a
statement about functions on matrices over the local field. The two columns of
the removed plane are singled out, and the maps in question are invariant under
the left action and under a group of matrices fixing the first of them.

% before: asm-theta-agree-local-mult-ramified-reduction
Such a map is determined by its values on the functions of full rank in the
middle columns, with the indicator function of the integers in the two other
columns. The ramification of $\nu$ makes the map vanish on functions supported
where the first of these columns is not zero.

% before: asm-theta-agree-local-mult-ramified-irreducible
The functions of full rank modulo their augmentation form an irreducible
representation of the general linear group. This rests on Bernstein's
irreducibility of the induced representations of two characters.

% before: asm-theta-agree-local-mult-ramified-gl
With the reduction and the irreducibility, an invariant map that is
equivariant under the middle block is determined up to a scalar. Dixmier's lemma
supplies the scalar.

% before: asm-theta-agree-local-mult-ramified-pair
For a Weil pair at a split place, the dictionary and a frame adapted to the
plane carry the equivariance of Definition~\ref{asm-theta-agree-local-equivariant}
to the conditions of Situation~\ref{asm-theta-agree-local-mult-ramified-situation}.
The unipotent and Levi elements become the matrices of $P^\circ$ and the middle
block.

% before: asm-theta-agree-local-mult-ramified
At the places of $T$ the space of equivariant maps has dimension at most one.
This gives the comparison of the two local maps at those places.

% before: asm-theta-agree-local-mult-archimedean-modules
At the infinite places the comparison goes through the Lie algebras
$M_{R+S}(\mathbf C)$ of all the places at once, acting on the space of a real
map. Two maps between such modules turn out to be proportional when each module
is generated by one highest-weight vector and contains no other vector of that
weight up to scalars.

% before: asm-theta-agree-local-mult-archimedean-gen
Stable subspaces are built by enlarging a subspace stable under part of an
algebra inside a larger subspace stable under another part. When the two parts
together span the algebra, the enlargement can be made stable under all of it.

% before: asm-theta-agree-local-mult-archimedean-low
A diagonal element orders the vectors of a module by its eigenvalues. The sum of
the eigenspaces below a given eigenvalue is stable under the operators that
lower the eigenvalue, and it is killed by every map on which the diagonal
element acts by that eigenvalue.

% before: asm-theta-agree-local-mult-archimedean-cyclic
With a suitable diagonal element, the two previous lemmas place a highest-weight
vector in a subspace stable under the block-diagonal matrices that lies, up to
vectors of lower eigenvalue, inside a given subspace. The diagonal element is
chosen so that each block-diagonal matrix unit lowers the eigenvalue, kills the
vector or preserves the given subspace.

% before: asm-theta-agree-local-mult-archimedean-generator
In a module generated with Casimir gap, the highest-weight vector controls the
images under intertwining maps: a sum of images that is a highest-weight vector
of the same weight is a combination of the images of the generator. The Casimir
gap moves each vector into the finite-dimensional space of the highest-weight
datum, and the previous lemma removes its components of lower weight.

% before: asm-theta-agree-local-mult-archimedean-tensor
At several infinite places the modules generated with Casimir gap are combined
through a tensor product. A quotient of the tensor product that kills the
kernels at each place is a highest-weight module of the combined weight, as
generation and the uniqueness of the highest-weight vector can be checked one
place at a time.

% before: asm-theta-agree-local-mult-archimedean-action
The derivatives of the Weil representation along lines in $U_W$, read in the
Fock model, give an action of the Lie algebras of the infinite places on the
space of a real map. At each place the Fock space divided by the action of the
Lie algebra of $V$ is a quotient generated with Casimir gap, so the space of the
real map is a highest-weight module.

% before: asm-theta-agree-local-mult-archimedean-embeddings
The plane $H'$ splits two coordinates off the skew-Hermitian space $W$. For
matrix units this becomes an embedding of the index set of $W_1$ into that of
$W$ that leaves out one index in each block, placed first or last according to
the step.

% before: asm-theta-agree-local-mult-archimedean-relations
A map equivariant at $\infty$ intertwines the Levi part of the parabolic
subalgebra defined by $H'$ with the action on the smaller space, up to
characters at the two extra indices, and kills its unipotent part. In matrix
units these are six families of relations.

% before: asm-theta-agree-local-mult-archimedean-zero
A map satisfying the parabolic relations that kills the highest-weight vector
is zero. Its kernel contains a subspace stable under the block-diagonal
matrices, and the relations make the kernel stable under the rest of the
algebra.

% before: asm-theta-agree-local-mult-archimedean-proportional
Two maps satisfying the same parabolic relations send the highest-weight vector
to multiples of the highest-weight vector of the smaller module. By the
previous lemma they are then proportional.

% before: asm-theta-agree-local-mult-archimedean-frame
To read $H'$ in matrix units, a frame of $W$ at each infinite place must carry
the two extra indices to the isotropic vectors of $H'$, up to scalars, and agree
with a given frame of $W_1$ on the image of $\iota$. Such a frame is built from
the vector $x$ and its image under the Cartan involution.

% before: asm-theta-agree-local-mult-archimedean-torus
Levi elements along a line need units of $\mathbf A_E$ with given components at
the infinite places and trivial finite components. The derivative of the
character of such a family of units enters the equivariance condition at
$\infty$.

% before: asm-theta-agree-local-mult-archimedean-levi
With frames adapted to $H'$, a line generated by matrices that preserve the
lines of $x_e$ and $y_e$ and commute with $P_e$ consists of Levi elements. A
line generated by matrices that kill $x_e$ and move the image of $P_e$ into the
line of $x_e$ lies in the unipotent group of $H'$.

% before: asm-theta-agree-local-mult-archimedean
At the infinite places, too, the space of equivariant maps has dimension at most
one. Together with Lemmas~\ref{asm-theta-agree-local-mult-unramified}
and~\ref{asm-theta-agree-local-mult-ramified} this compares the local maps of
the cone and of the evaluation at every place.

% before: asm-theta-agree-kills-class
The theta integrals of pure tensors with a component in a kernel are shown to
vanish through their norms on $U_W(F)\backslash U_W$. This needs a space of
functions on which the theta integral is defined, linear, continuous and square
integrable. For the second space of a triple only the full-rank data are known
to have these properties.

% before: asm-theta-agree-kills-class-triple
Both Weil pairs of a triple have such a space. For the first it consists of all
Schwartz-Bruhat functions, and for the second of the full-rank data; the
properties come from the convergence, continuity and square-integrability
results proved earlier.

% before: asm-theta-agree-kills-forms
Integration over a fundamental domain turns a family of theta integrals into a
Hermitian form. The arguments at the finite and at the infinite places use its
invariance under right translations and the skew-symmetry of derivatives along
lines.

% before: asm-theta-agree-kills-finite
At a finite place, Howe duality, taken here as a hypothesis on the pair, gives
a vector whose translates span the Schwartz-Bruhat functions modulo the
augmentation space. The Hermitian form of
the theta integrals then vanishes on every proper theta subspace at the place,
and the theta integral of a pure tensor with its component there vanishes.

% before: asm-theta-agree-kills-radical-maximal
At the infinite places the same argument needs a property of Hermitian forms on
the polynomial Fock space of the pair. A form that vanishes on the derivatives
along lines in $U_Y$ and is skew for the lines in $U_W$ should have a radical
that no stable subspace other than the whole space properly contains.

% before: asm-theta-agree-kills-radical-greatest
In a module generated with Casimir gap, the stable subspaces other than the
module have a largest one. A form for which the matrix units are adjoint up to
sign vanishes on it, because it lies in the image of the shifted Casimir
operator.

% before: asm-theta-agree-kills-radical-place
At one infinite place this applies to the quotients of
Lemma~\ref{asm-theta-agree-real-quotient}. Pulled back to the Fock space, the
largest proper stable subspace is killed by every form that vanishes on the
span defining the quotient and is skew for the second action.

% before: asm-theta-agree-kills-radical-tensor
The Fock space of a Weil pair is a tensor product over the infinite places.
If each factor has a stable subspace with irreducible quotient, the tensor
product of the quotients is irreducible. So a stable subspace that contains a
subspace $R$, which in turn contains every pure tensor with a factor in one of
those subspaces, is $R$ or the whole space.

% before: asm-theta-agree-kills-radical
Both Weil pairs of a triple have maximal radicals. In the Fock model of the
pair the lines act place by place through the two actions, so the form becomes
a form on a tensor product to which the three previous lemmas apply.

% before: asm-theta-agree-kills-real
With this property and Howe duality at the real places, the radical of the
Hermitian form of the theta integrals is a theta subspace that contains every
proper theta subspace. So the theta integral of a pure tensor whose
archimedean component lies in a proper theta subspace vanishes.

% before: asm-theta-agree-kills-full-rank
For the second space only components of full rank at $\tau_1$ are allowed at
$v_1$. The augmentation argument there needs that an element of the augmentation
space in the span of such components is a combination of differences of
components of full rank. This follows from the irreducibility of the functions
of full rank modulo their augmentation.

% before: asm-theta-agree-local-injective-kills
The nonvanishing of theta integrals of pure tensors is proved by contradiction
from its converse. A pure tensor with a component in the kernel of the real
map or of one of the finite maps has theta integral $0$.

% before: asm-theta-agree-local-injective-kills-zero
For the second space with $q=0$ the condition of full rank at $v_1$ can be
dropped. The component at $v_1$ is changed by an element of the augmentation
space into one of full rank, which does not change the theta integral.

% before: asm-theta-agree-local-injective-general
The vanishing of one theta integral then spreads to all pure tensors. At each
finite place, and at the infinite places, the pure tensors with vanishing theta
integral give a stable subspace of the irreducible space of the real or finite
map, and it contains a nonzero vector.

% before: asm-theta-agree-local-injective
The local maps are built from theta integrals of pure tensors. Their
definition needs that such a theta integral is nonzero when all its local
images are nonzero.

% before: asm-theta-agree-local-split-invariant
The constant terms used to compare theta integrals come in two parts with
different laws. Each part is invariant under the unipotent group of the
plane, and both vanish with the theta integral.

% before: asm-theta-agree-local-transfer
The evaluation map and the cone are treated alike, as transfers from one
side of a matched quadruple to a side of the smaller one. A transfer
factors on pure tensors, preserves vanishing of theta integrals, and
intertwines Levi elements with a character.

% before: asm-theta-agree-local-transfer-criteria
The conditions on a transfer follow from the constant-term formulas and from
the factorization of the map. The evaluation maps and the cones are covered
by the two forms of the factorization below.

% before: asm-theta-agree-local-evaluation-pol
For a test vector of an evaluation map one needs its archimedean component
to reach a polynomial function with nonzero image. This follows from the
surjectivity of the evaluation map on polynomial data.

% before: asm-theta-agree-local-maps-sides
The local maps are defined on the three kinds of sides alike. What the
construction uses of a side is collected below: its local spaces map onto the
spaces of the real and finite maps, theta integrals of pure tensors are linear
and continuous, they vanish when a component has image $0$, and the actions of
$U_W$ on the spaces of the real and finite maps are well defined.

% before: asm-theta-agree-local-maps
A transfer with a test vector induces local maps on the spaces of the real
and finite maps. Part (1) produces the archimedean half of a test vector for
the cone.

% before: asm-theta-agree-local-alpha
The cone of a polynomial datum of the first space corresponds, in the smaller
triple, to a multiple of the evaluation of its partner. The multiple is a step
factor times a partial $L$-value.

% before: asm-theta-agree-local-beta
For the removal of a plane from the second space, the evaluation and the cone
change roles. The cone of a full-rank datum is again of full rank.

% before: asm-theta-agree-end-notions
The last two lemmas of the induction compare the two theta integrals through
their difference, an automorphic function on $U_W$. Its constant terms are
taken along planes that need not be adapted, and it is shown to vanish through
cuspidality and orthogonality on $U_W(F)\backslash U_W$.

% before: asm-theta-agree-end-unipotent
Constant terms along a plane are integrals over a compact quotient of its
unipotent group. The facts used about this group are those known for adapted
planes, together with the transitivity of $U_W(F)$ on nonzero isotropic
vectors.

% before: asm-theta-agree-end-invariant
An automorphic function is bounded on the unipotent group of a plane, after
translation by rational elements, so its constant terms are integrals of
bounded functions over sets of finite measure. The constant term is again
invariant under the unipotent group.

% before: asm-theta-agree-end-constant-terms
Constant terms are linear on automorphic functions. When a constant term is
the sum of two parts with laws of different exponents, a Levi element with
$|t|>1$ separates them, and each part is invariant under the unipotent group.

% before: asm-theta-agree-end-move
Purity is defined through one plane. The first vectors of two planes are
conjugate under $U_W(F)$, and the unipotent group of a plane depends only on
the line of its first vector, so an automorphic function that is pure along
one plane is pure along every plane, with the same law.

% before: asm-theta-agree-end-iwasawa
A function with a law and invariant under the unipotent group is determined by
its values at the products of Levi elements for $(1,h_1)$ with elements of
$K^\theta_W$. This rests on a decomposition of $U_W$ along an adapted plane.

% before: asm-theta-agree-end-automorphic
The theta integrals that enter the difference are automorphic. Their
invariance under $U_W(F)$ comes from that of the theta function, and their
continuity from the domination of the theta functions on a fundamental domain.

% before: asm-theta-agree-end-stable-level
An element of $K^\theta_W$ is constrained only at the infinite places, and its
finite components can be nontrivial at infinitely many places. Correspondence
is preserved one place at a time, so the remaining factor must fix the data. A
Schwartz-Bruhat function is fixed by every element near $1$ whose archimedean
components are trivial.

% before: asm-theta-agree-end-stable
The hypotheses of Lemmas~\ref{asm-theta-agree-end-alpha}
and~\ref{asm-theta-agree-end-beta} concern values at Levi elements. To use them
at the products of Levi elements with elements $k$ of $K^\theta_W$,
correspondence must be preserved by $\omega(1,k)$ and $\omega'(1,k)$.

% before: asm-theta-agree-end-rechoose
The induction removes planes after re-choosing the archimedean and local data
of the spaces. A right translation carries the differences of a triple to
differences of the re-chosen triple.

% before: asm-theta-agree-end-tower
Removing a plane from the first space lowers $p$ by one and keeps $q$.
Lemma~\ref{asm-theta-agree-tower-alpha} does this when $q=0$, and the same
construction applies for every $q$.

% before: asm-theta-agree-end-alpha-zero-stable
The next proofs translate a difference on the right by an element of
$K^\theta_W$. Translation keeps the correspondence of polynomial data, so it
keeps the polynomial data of the first space and the full-rank polynomial data
of the second space.

% before: asm-theta-agree-end-lower-kill
The proof that an evaluation can be given full rank works at the place $v_1$,
through the matrix model of a linear model of the second kind. There, a function
is shown to lie in a space of relations by multiplying it by locally constant
functions. The argument of Lemma~\ref{asm-theta-agree-coinvariants-kill} works
with a set of points in place of a single one.

% before: asm-theta-agree-end-lower-eigen
A linear form that transforms under the Schrödinger operators of a set of
elements by given scalars sees a function only on the matrices where those
elements act by the same scalars. This is the form of
Lemma~\ref{asm-theta-agree-coinvariants-eigen}(3) for a set of elements.

% before: asm-theta-agree-end-lower-step
The relations needed below are those of the group $GL_n$ acting only on
functions that already vanish where some rows are dependent. The step of
Lemma~\ref{asm-theta-agree-coinvariants-step} uses no other relations.

% before: asm-theta-agree-end-lower-corank
A function that vanishes wherever all but one row are dependent differs from a
function of full rank by such relations. Near a matrix whose rows other than
one are independent, one step of the previous lemma suffices.

% before: asm-theta-agree-end-lower-rank
The evaluation along $H'$ reads a function on $n\times d$ matrices on the
matrices $Y_1\kappa+yr^{\mathsf T}$. When all but one row of $Y_1$ are always
dependent, these matrices never have full rank.

% before: asm-theta-agree-end-lower-lift
The relations for the smaller model must come from relations for the larger
one. A function that vanishes where the rows other than one are dependent is
the evaluation of a function of full rank, built from it and a translate of a
fixed function.

% before: asm-theta-agree-end-lower-local
In the matrix models at $v_1$, the evaluation along $H'$ commutes with
translations and reads a function only on the affine spaces
$Y_1\kappa+F_{v_1}^nr^{\mathsf T}$. These are the two inputs of the previous
lemma.

% before: asm-theta-agree-end-alpha-zero-lower
On the second space, the splitting along $H'$ of
Lemma~\ref{asm-theta-agree-split-any} reads the $\alpha$-part through the
evaluation of a full-rank datum, and needs this evaluation to have full rank. At
$v_1$ a full-rank function can be changed by an element of the augmentation
space so that its evaluation has full rank.

% before: asm-theta-agree-end-alpha-zero-move
Changing a full-rank datum at $v_1$ by an element of the augmentation space
does not change its theta integral, as theta integrals are invariant under the
group of the first space. So a full-rank polynomial datum of the second space
can be replaced by one with the same theta integral whose evaluation along
$H'$ has full rank.

% before: asm-theta-agree-end-alpha-zero-second
With such a datum, the theta integral of the second space splits along $H'$,
and its $\alpha$-part, read on Levi elements, is the theta integral of a
full-rank datum for a smaller Weil pair. Both the case $p=0$ and the removal of
a plane use this.

% before: asm-theta-agree-end-alpha-zero
When one of the two spaces is definite, one of the two parts of the constant
term of a difference vanishes. When $p=0$ this is the $\alpha$-part.

% before: asm-theta-agree-end-beta-zero-forms
The parts of the constant term that must vanish are theta integrals for
smaller Weil pairs, and these vanish for an archimedean reason. On the Fock
space, a linear form that transforms under the block-diagonal operators $d_1(Y)$
by the weight $-\frac12m\operatorname{tr}(Y)$ is zero once $|m|$ is large. The
integer $m$ will be an exponent of the character $\chi_W$.

% before: asm-theta-agree-end-beta-zero-real
Through a Fock map, a continuous linear form on the archimedean Schwartz space
that is invariant under $U_X$ at $\infty$ gives linear forms on the Fock space.
Differentiating along lines in $U_X$ shows that they transform as in the
previous lemma.

% before: asm-theta-agree-end-beta-zero-vanish
For a fixed finite part and a fixed point of $U_W$, the theta integral is such
an invariant form, since the integral over a fundamental domain is invariant
under right translation. So under the same conditions on the exponents and the
signatures, the theta integrals vanish.

% before: asm-theta-agree-end-beta-zero
When $q=0$ the $\beta$-part of the constant term of a difference vanishes.
This is the case of Lemma~\ref{asm-theta-agree-end-alpha}.

% before: asm-theta-agree-end-alpha-low-cone
After a plane has been removed from the first space, the $\alpha$-part of the
first theta integral is the theta integral of a cone. The cone of a polynomial
datum is again a polynomial datum, for the smaller model, because the cone
factors into local pieces.

% before: asm-theta-agree-end-alpha-low
After a plane has been removed from the first space, the $\alpha$-part of a
difference, read on Levi elements, is a difference for the smaller triple. The
induction on $p$ applies to it.

% before: asm-theta-agree-end-two-step-moderate
The proof that the lower part is pure compares constant terms along a
contracting element, and this needs a bound on growth. A difference grows at
most like a power of the height.

% before: asm-theta-agree-end-two-step-split
When $p+q\geq1$ the skew-Hermitian space is isotropic, so it contains a plane.
Every difference splits along it into an $\alpha$-part and a $\beta$-part,
after the space is re-chosen so that the plane is adapted.

% before: asm-theta-agree-end-kill-notions
The proof that the double constant term is invariant integrates in stages over
unipotent groups and then argues on $SL_2$ along a diagonal element. Both need
averages over lattices in locally compact groups, and the induction on $SL_2$
needs bumps, the growth of functions along the diagonal element and two
properties of reduction theory.

% before: asm-theta-agree-end-kill-averages
Averages over a lattice exist, are unique and are continuous, because a lattice
has a fundamental domain of finite positive measure. A conjugation that
preserves the group and the function also preserves the average.

% before: asm-theta-agree-end-kill-stages
When a group is the product of a normal subgroup and a complement, an average
over the group is the average over the complement of the average over the
subgroup. This is how two different orders of integration give the same
function.

% before: asm-theta-agree-end-kill-induction
On a group with a Siegel set along a contracting element, a recursion with a
factor of absolute value less than $1$ forces invariance for bounded functions.
Smoothing by a bump lowers the order of growth by one, so the same holds for
functions of any order of growth. The proof that a difference with a vanishing
constant term is cuspidal uses the same induction.

% before: asm-theta-agree-end-kill-siegel-reduction
The induction will be applied to $SL_2$ over the adeles, and its Siegel
property comes from reduction theory. Every element of $SL_2(\mathbf A_E)$ is a
rational translate of an upper triangular element of a compact set, times a
diagonal element that is large at the infinite places, times an element of a
second compact set.

% before: asm-theta-agree-end-kill-siegel-adeles
The Lipschitz property is built from the archimedean size of an adele and the
position of its finite part relative to small open subgroups. Three facts about
the adeles of a number field connect the two: boxes are compact, they form a
basis of neighbourhoods, and every adele has a rational translate in a fixed
box.

% before: asm-theta-agree-end-kill-siegel-bump
The bump of the Lipschitz property is a function of the entries of $y-1$. It
changes by at most a multiple of $b$ when the entries move by adeles of
archimedean size at most $b$ whose finite parts lie in a fixed open subgroup.

% before: asm-theta-agree-end-kill-siegel
The induction is applied to $SL_2$ over the adeles of $E$, with the upper
unipotent group and the diagonal element that is $2$ at the infinite places.
Reduction theory for $SL_2$ gives the Siegel property, and conjugation by this
element contracts the upper unipotent group at the infinite places.

% before: asm-theta-agree-end-kill-growth
The induction needs a growth bound along the diagonal element. It comes from
the bound by a power of the height, which survives averaging over a unipotent
group with a lattice.

% before: asm-theta-agree-end-kill-unipotent
The first average is over the unipotent group of the isotropic plane spanned by
$x$ and $\iota x_1$. These groups behave like those of a single isotropic line,
with the same proof.

% before: asm-theta-agree-end-kill-levi
Comparing the two orders of integration needs a Levi map that is a continuous
homomorphism with rational values, and a map back from the stabilizer of $x$
to $U_{W_1}$. Both come from a left inverse of $\iota$ that kills $x$ and $y$.

% before: asm-theta-agree-end-two-step-kill
A plane in $W$ and a plane in $W_1$ give an embedding of $SL_2$ into $U_W$,
whose diagonal elements are Levi elements for both planes. For the double
constant term of an automorphic function of moderate growth, a recursion along
such a diagonal element with a factor of absolute value less than $1$ forces
invariance.

% before: asm-theta-agree-end-two-step
For the induction hypothesis to apply, the difference read on Levi elements
must again be pure. It is, with the same law.

% before: asm-theta-agree-end-cusp-notions
Cuspidality is proved by induction on the dimension of an isotropic subspace.
The subspace is spanned by the columns of a frame, whose Levi map embeds $GL_j$
into $U_W$. The constant term along the subspace, read on this copy of $GL_j$,
is studied through its Whittaker coefficient.

% before: asm-theta-agree-end-cusp-frames
The Levi map of a frame acts on the frame by the given matrix and normalizes
the unipotent group of its span. An element that acts on the span by a matrix
and trivially on the rest is a Levi element times an element of the unipotent
group. Every isotropic subspace is the span of a frame.

% before: asm-theta-agree-end-cusp-constant
The constant term along the span of a frame is invariant under the rational
points of the Levi copy of $GL_j$. If the constant terms along the smaller spans
vanish, so do its averages over the unipotent radicals of the standard
parabolic subgroups of $GL_j$. Averaging in stages over a retraction gives this.

% before: asm-theta-agree-end-cusp-expansion
Successive integrals over the columns of the upper unipotent group of $GL_a$
against a character turn a function into its Whittaker coefficient. When the
function has vanishing averages over the unipotent radicals of the standard
parabolic subgroups, its Whittaker coefficient vanishes only if the function
does.

% before: asm-theta-agree-end-cusp-coefficient
The Whittaker coefficient along a frame is a Whittaker functional for a
standard character. If it vanishes on all translates of an automorphic function
whose constant terms along the smaller spans vanish, the constant term along the
span vanishes too.

% before: asm-theta-agree-end-cusp-line
All isotropic lines are rational translates of each other. So the vanishing of
the constant term along one plane gives the case of dimension $1$.

% before: asm-theta-agree-end-cusp-kill
In dimension $2$ the Whittaker coefficient is an integral over one copy of the
adeles, and $SL_2$ acts on it through the Levi map. The induction on $SL_2$ of
Lemma~\ref{asm-theta-agree-end-kill-induction} shows that such an integral
vanishes when its values at a sequence of translates are proportional to a
sequence tending to infinity.

% before: asm-theta-agree-end-cusp-two
The constant term of a difference has moderate growth, and its average along
the first column vanishes by the case of dimension $1$. So the reduction applies
to it on $SL_2$.

% before: asm-theta-agree-end-cusp-keys-places
The properties of Whittaker functionals on the differences are proved at one
finite place where $E$ splits, chosen outside finitely many exceptional places.
Three facts about the places of $F$ allow this choice: the projections $\pi_t$
are continuous, a given adele is integral at almost all places, and there are
infinitely many split places.

% before: asm-theta-agree-end-cusp-keys-family
At a split place $v$, a theta integral depends on the component at $v$ of its
Schwartz-Bruhat function, and a linear model turns this component into a
function on matrices over $F_v$. The resulting map from functions on matrices to
automorphic functions transforms by characters under the left action of
$GL_n(F_v)$ and under right translation by the unitary group at $v$.

% before: asm-theta-agree-end-cusp-keys-local
For the theta integrals of a Weil pair, the component at $v$ can be replaced by
any Schwartz-Bruhat function, and the theta integral of the result defines a
local family. Its laws come from the invariance of theta integrals under the
first unitary group and from the action of the second, both read through a
linear model.

% before: asm-theta-agree-end-cusp-keys-complement
To pass from a frame to coordinates at $v$, the columns of the frame are
completed to an invertible matrix. The completion can be chosen so that the
dual columns of the frame pick out the original ones.

% before: asm-theta-agree-end-cusp-keys-form
A Whittaker functional applied to a local family gives a linear form on the
functions on matrices over $F_v$. The law of the functional under the elements
at $v$ that act on the frame by upper triangular matrices becomes a law of the
form for adding multiples of one column to the other columns.

% before: asm-theta-agree-end-cusp-keys-three
The linear forms that arise transform by a character under the left action of
$GL_n$ and by a nontrivial character under column operations. With such laws
for two consecutive pairs of columns a form vanishes, which gives the case of
rank at least $3$. The proof is local on the space of matrices: near each point
a unipotent symmetry fixes the point and multiplies the form by a scalar
different from $1$.

% before: asm-theta-agree-end-cusp-keys-two
With a law for one pair of columns only, the form need not vanish. Its values
at the torus translates of the indicator function of the integral matrices are
a fixed multiple of an explicit sequence, which grows when
$|\xi(\varpi)q^{n-d}|$ exceeds $q$.

% before: asm-theta-agree-end-cusp-keys-vacuum
The computation applies to the vacuum at $v$ once a linear model carries it to
a multiple of the indicator function of the integral matrices. At almost all
places this holds, as the vacuum is fixed by the Heisenberg operators of the
integral vectors and a linear model is a Heisenberg chart.

% before: asm-theta-agree-end-cusp-keys-torus
The torus translates come from elements of $SL_2(\mathbf A_E)$ through the Levi
map of a frame of rank $2$. Such an element can be taken at $v$, where it acts
in the coordinates given by the completion of the frame by a diagonal matrix.

% before: asm-theta-agree-end-cusp-keys
The remaining inputs are properties of Whittaker functionals on the
differences. In rank at least $3$ they vanish; in rank $2$ their values on
suitable translates are proportional to a sequence tending to infinity.

% before: asm-theta-agree-end-cusp
A difference whose constant term along one plane vanishes has vanishing
integrals over the unipotent groups of all isotropic subspaces. So it is
cuspidal.

% before: asm-theta-agree-end-orth-notions
The orthogonality of a cuspidal difference to theta integrals is proved through
its theta lifts, the integrals of the difference against theta functions. The
rapid decay of cusp forms needs smoothness at the finite places and a common
growth bound for the derivatives along lines, and the invariance of the theta
lifts at a clean place needs the difference to be built from data at that
place.

% before: asm-theta-agree-end-orth-lift
A rapidly decreasing function absorbs any polynomial growth that is invariant
under the rational points, and theta functions grow polynomially. So the theta
lift of a rapidly decreasing function is defined for every Schwartz-Bruhat
datum, and it is an automorphic function on the first group.

% before: asm-theta-agree-end-orth-perp
If the theta lifts of a rapidly decreasing function against a datum vanish,
Fubini's theorem moves the integral over the first group outside. The theta
integral of the datum is then orthogonal to the function.

% before: asm-theta-agree-end-orth-unbounded
A constant theta lift against full-rank data is bounded by a multiple of the
inverse of the reduced height of the first group. On an isotropic Hermitian
space of rank at least $3$ that height is unbounded, as a Siegel Witt torus
there has positive size.

% before: asm-theta-agree-end-orth-zero
For full-rank data the theta function decays in the reduced height of the first
group, so a constant theta lift against such data vanishes. At $v_1$ every
Schwartz-Bruhat function is a full-rank one modulo the augmentation, which a
constant theta lift does not see. So a constant theta lift vanishes for every
datum.

% before: asm-theta-agree-end-orth-cusp-rapid
A cuspidal difference is rapidly decreasing by the rapid decay of cusp forms,
due to Harish-Chandra. Here it is stated for unitary groups, with uniform
moderate growth expressed through derivatives along lines and the decay
measured by the reduced height.

% before: asm-theta-agree-end-orth-smooth-pure
A theta integral of a product of a polynomial Fock vector and a dominated
finite part has the three properties the rapid decay of cusp forms needs. Its
derivatives along lines are again such theta integrals, by the archimedean
operators of the Weil representation, and a fixed finite part is invariant
under a neighbourhood of $1$ at the finite places.

% before: asm-theta-agree-end-orth-smooth
The rapid decay of cusp forms applies to a cuspidal difference once it is smooth
at the finite places and of uniform moderate growth. Both properties come from
the theta integrals of the products that make up the polynomial and full-rank
data.

% before: asm-theta-agree-end-orth-rapid
A cuspidal difference is automorphic, smooth at the finite places and of
uniform moderate growth, so the rapid decay of cusp forms applies to it. Its
theta lifts are then defined.

% before: asm-theta-agree-end-orth-local-decomp
The local decomposition of a difference rests on a fact of linear algebra. A
vector in a stable subspace that is a sum of values of intertwiners out of an
irreducible representation is already such a sum with intertwiners whose values
lie in that subspace, because a submodule of a finite sum of simple modules is
a direct summand.

% before: asm-theta-agree-end-orth-local
At a clean place other than $v_1$ the finite maps matched to the triple identify
the coinvariants of the two sides. A rapidly decreasing difference then
decomposes into images of local maps at that place, for each of the two Weil
representations.

% before: asm-theta-agree-end-orth-local-pair
The theta lift against one side of a function built from local data of the
other side is a sesquilinear form at the place. When the character of the
skew-Hermitian space is nontrivial on a uniformizer, such forms are invariant
under the general linear group of the first space, and so the theta lift is
invariant at that place.

% before: asm-theta-agree-end-orth-constant
The theta lifts of a rapidly decreasing difference are continuous, automorphic
and invariant at the places outside a finite set at which the rational prime
splits completely in $E$. A function with these properties is constant.

% before: asm-theta-agree-end-orth-iso
When the first space is isotropic, a cuspidal difference is orthogonal to the
theta integrals of the first space, and likewise for the second space.
Lemma~\ref{asm-theta-agree-end-orth-an} covers the definite spaces.

% before: asm-theta-agree-end-orth-cusp-lift-unfolding
The integral of a cusp form against an Eisenstein series is computed by
unfolding the Eisenstein sum over the orbits of a discrete group. Two general
facts about fundamental domains do this: an orbit contributes the integral over
a fundamental domain of its stabilizer, and an integral vanishes when one factor
is invariant under a normal subgroup along which the other has mean zero.

% before: asm-theta-agree-end-orth-cusp-lift-inputs
The doubling method supplies the Eisenstein series on the doubled group of the
definite space that the next proof integrates against. Its value at the
Siegel-Weil point is an average of products of theta functions. For a coset of
the Siegel parabolic subgroup outside the open orbit of the unitary group, the
conjugated parabolic subgroup contains the unipotent radical of a parabolic
subgroup of the unitary group, along which a cusp form has mean zero.

% before: asm-theta-agree-end-orth-cusp-lift
On a definite first space a nonzero pairing of a cuspidal difference with one
theta integral makes all the theta integrals cuspidal. The Siegel-Weil formula
writes the integral of the difference against an Eisenstein series of the
doubled group as a multiple of a theta integral, and that integral is cuspidal
because the difference is.

% before: asm-theta-agree-end-orth-rechoice
The tower lemmas re-choose the archimedean and local data of the spaces, which
changes the polynomial data but not the theta integrals. Cuspidality passes to
the polynomial data of the re-chosen spaces, because an element moving the
Cartan involutions carries one Fock space onto the other.

% before: asm-theta-agree-end-orth-constant-term
The constant term along an adapted plane is the integral over the unipotent
group of the isotropic line of the plane. So a cuspidal function has constant
term $0$, and constant terms are unique.

% before: asm-theta-agree-end-orth-an
For a definite space the orthogonality uses the Siegel-Weil formula and the
continuation of Eisenstein series. The two lemmas together give orthogonality
to both sides whenever $p\geq1$ or $q\geq1$.

% before: asm-theta-agree-end-square-exponent
On the side of the second space only full-rank data enter, and their lattice
sums decay with the excess. The next inequality trades that decay against the
growth of the bound on Siegel sets. It is an estimate on cones and involves no
groups.

% before: asm-theta-agree-end-square-torus
Square-integrability on a fundamental domain reduces to a bound for the theta
function on products of two Witt tori. The theta integral is bounded on the
torus of $W$ by integrating over the torus of the first space, and the square
of that bound is then integrated over the torus of $W$.

% before: asm-theta-agree-end-square
The orthogonality relations are integrals over a fundamental domain for
$U_W(F)$. The theta integrals are square-integrable there, so these integrals
converge.

% before: asm-theta-agree-end-perp-zero
A difference that is orthogonal to both sides is orthogonal to itself. Being
continuous and automorphic, it vanishes.

% before: asm-theta-agree-end-differences
Given corresponding data and a candidate constant $c$, the function
$I(\Phi)-c\,I'(\Phi')$ is a difference. Its constant term along an adapted
plane splits as the same combination of the splittings of the two theta
integrals.

% before: asm-theta-agree-end-compare
A part of a constant term is determined by its law and its values at the
products $lk$. This is the form in which the hypotheses of
Lemmas~\ref{asm-theta-agree-end-alpha} and~\ref{asm-theta-agree-end-beta} are
used.

% before: asm-theta-agree-end-zero
The vanishing of a difference follows from the vanishing of one constant term.
The difference is then cuspidal, orthogonal to both sides, and zero.

% before: asm-theta-agree-end-pure
When $q\geq1$ the $\beta$-parts are compared first, and what remains is a
difference that is pure of the $\alpha$-law. Such a difference vanishes, by
induction on $p$: planes are removed from the first space until $p=0$, and then
one plane from the second space.

% before: asm-theta-agree-end-alpha
After the plane has been removed, the constant terms of the two theta
integrals agree at the Levi elements. When the second space is definite this
agreement of $\alpha$-parts forces the agreement of the theta integrals
themselves.

% before: asm-theta-agree-end-beta
When the second space is not definite, the comparison of the $\beta$-parts
plays the same role. With these two lemmas the induction closes.

% before: asm-theta-agree-constants
The constants of Definition~\ref{asm-theta-agree-matched} are defined by
matching data, and the induction produces polynomial constants. The two
notions agree, by approximating matching data with corresponding polynomial
data.

% before: asm-theta-agree-transfer
Re-choosing the archimedean data changes the polynomial data. A theta integral
that is nonzero on the polynomial data of one choice is nonzero on those of
the other, by density.

% before: asm-theta-agree-key-beta
The lemmas on constant terms, evaluation maps, cones and local computations
combine into one comparison for the removal of a plane from the second space.
It relates the $\beta$-parts of the two theta integrals through the polynomial
constant of the smaller triple.

% before: asm-theta-agree-key-alpha
The same combination for the removal of a plane from the first space relates
the $\alpha$-parts. Here the second space of the smaller triple is definite,
and Lemma~\ref{asm-theta-agree-anisotropic} removes the full-rank condition
on its data.

% before: asm-theta-agree-claim
Nondegeneracy is proved by induction on $p+q$. The definite case starts it,
and each step removes a hyperbolic plane with the two comparisons above.

% before: asm-theta-agree-derivation
With the two published results, the agreement of theta integrals and the
local matching become consequences of proved statements. The derivation
splits them into the parts of Definition~\ref{th-theta-hypotheses}. The
Eisenstein series are used for the constant of its item~(2) and for the
formulas of its item~(3).

% before: asm-every-character-literature
Most of the published results of Chapter~\ref{chapter-literature} are proved
outright. Three hypotheses remain open at this point: the Galois
representations, used twice, and the existence of maps satisfying the local
Langlands laws. The bound on Tate-Shafarevich groups is also kept as a
hypothesis.

% before: asm-every-character
The Galois representations of Chapter~\ref{chapter-galois-representations}
give the bound on Tate-Shafarevich groups of
Chapter~\ref{chapter-tate-shafarevich}, and the local Langlands maps exist.
So the half-plane statement depends only on the existence of these Galois
representations, which Chapter~\ref{chapter-galois-representations} proves.

% before: asm-flow
The previous sections combine the links by their logical form. At a single
member, the comparison bounds the slope of the dual of the glued lattice from
below, through divisibility of periods, and from above, through covolumes. The
inert-prime sum is controlled by the gap between the two bounds.

% before: asm-parameters
The constants of the proof are chosen in a fixed order, and some of them are
numbers. Listing them in one place shows the scales at which each bound is
used. The exponents $3/2$, $401$, $1/52$ and $9$ all appear here.
